\chapter[Vertex algebra axioms, Borcherds identity, reconstruction]{Vertex algebra axioms, the Borcherds identity, and reconstruction}\label{ch:L03}
\begingroup\setlength{\emergencystretch}{3em}

\providecommand{\lthreevac}{\mathbf{1}}
\providecommand{\lthreeV}{\mathcal{V}}
\providecommand{\lthreeY}{Y}
\providecommand{\lthreeBL}{B_{\mathrm L}}
\providecommand{\lthreeBR}{B_{\mathrm R}}
\providecommand{\lthreeD}{\partial}
\providecommand{\lthreeno}[1]{{:}#1{:}}

This chapter documents the release of the LeanVOA library named in the front matter. Every Lean quotation below is checked by \texttt{scripts/textbook/sync.py}
to occur verbatim (up to whitespace and comments) in the source files of that release. Checks that were run at an earlier state of the library and not re-run for this edition are marked as such. The chapter covers the twenty-seven files assigned to it in
\texttt{Notes/Textbook/assignment.txt}, all in \leanpath{VOA/\allowbreak Core/}:
\leanpath{BareVA}, \leanpath{VertexAlgebra}, \leanpath{VertexOperatorAlgebra}, \leanpath{Fields},
\leanpath{Identities}, \leanpath{OPE}, \leanpath{Goddard}, \leanpath{Dong}, \leanpath{FieldOps},
\leanpath{FieldCombinatorics}, \leanpath{Reconstruction}, \leanpath{BorcherdsIdentity},
\leanpath{Borcherds}, \leanpath{Skew}, \leanpath{WeakAssoc}, \leanpath{SkewSymmetry},
\leanpath{FLM}, \leanpath{Commutativity}, \leanpath{ReconstructionExistence},
\leanpath{TensorFieldLift}, the five \leanpath{LocalSystem*} files, \leanpath{FieldProductIterate}
and \leanpath{JacobiLocal} (each with suffix \texttt{.lean}).

\section[What this chapter is about]{What this chapter is about, and what it found}\label{l03:sec:intro}

Chapter~\ref{ch:01} stated the library's structures \lean{PreVertexAlgebra} (Gannon's axioms VA1--VA5 on a graded
space, which, as this chapter shows, do \emph{not} make a vertex algebra in the standard sense; hence the prefix) and
\lean{VertexOperatorAlgebra} (the Lepowsky--Li/FLM vertex operator algebra, or VOA, \emph{of CFT type}, the CFT-type
layer over the general notion \lean{GradedVertexOperatorAlgebra}), and it stated the most important theorems about them --- the Borcherds identity,
skew-symmetry, Goddard uniqueness, Dong's lemma, the reconstruction theorem --- with short proofs.
Chapter~\ref{ch:L01} develops the formal calculus (formal distributions, the delta function,
locality, the decomposition of a local distribution into derivatives of the delta function) on
which all of this rests. The present chapter goes one level deeper. It has four jobs.

\begin{enumerate}
\item \textbf{Audit the axioms.} We compare the library's structures \lean{PreVertexAlgebra},
  \lean{BarePreVA} and \lean{VertexOperatorAlgebra} with the textbook definitions of Borcherds,
  Frenkel--Lepowsky--Meurman (FLM), Kac, Frenkel--Ben-Zvi and Lepowsky--Li, axiom by axiom:
  is the translation operator part of the data or derived? is truncation required for all
  pairs? is the Jacobi identity an axiom, and if not, is it proved in all indices?
\item \textbf{Present the derivation of the Borcherds identity} from locality, step by step, as
  the library actually carries it out: Goddard uniqueness, Dong's lemma, translation covariance,
  the field $n$-th product theorem for $n\ge0$ and $n<0$, the iterate formula for all $\ell$,
  and the Pascal-recursion argument that assembles the full three-index identity.
\item \textbf{Present the further structure theory} proved in these files: skew-symmetry and the
  weak-associativity argument behind it, the FLM associativity formula, commutative vertex
  algebras, the operator-first reconstruction theorem, the tensor-product lift of fields, Li's
  local systems of weak vertex operators (including a Borcherds identity for an arbitrary pair of
  local fields on an arbitrary vector space), and the localisation of the first-order Jacobi
  identity used in deformation theory.
\item \textbf{Record precisely what is conditional, what is incomplete, and where the library's
  documentation over-claims.}
\end{enumerate}

\begin{warning}[title={The main finding of this chapter: Gannon's axioms, \texttt{PreVertexAlgebra}, do not give a vertex algebra in the textbook sense}]
The library's structure \lean{VOA.PreVertexAlgebra} transcribes Gannon's Definition 5.1.3(a): a
graded space with a state--field map satisfying vacuum, creation, locality and truncation. It has
\textbf{no translation-covariance axiom}. The textbook definitions (Kac, Frenkel--Ben-Zvi) include
a translation operator $T$ with $[T,Y(u,z)] = \partial_zY(u,z)$ as part of the data; the
FLM/Lepowsky--Li definitions take the Jacobi identity as an axiom, from which translation
covariance follows. Without it, \emph{the Borcherds identity does not follow.} We exhibit in
Section~\ref{l03:sec:counter} a three-dimensional graded space satisfying every field of
\lean{PreVertexAlgebra} on which the Borcherds identity, skew-symmetry and the derivative property
$Y(Tu,z) = \partial_zY(u,z)$ all fail. The counterexample was checked in Lean for this chapter, axiom-clean, and has since been
vendored into the library as \lean{VOA.TranslationCovariantCounterexample.VA}, with
\lean{not_isTranslationCovariant} and \lean{jacobi_fails} (file
\leanpath{VOA/\allowbreak Core/\allowbreak TranslationCovariant.lean}).

The library itself is consistent with this: every theorem that needs translation covariance
(the Borcherds identity, skew-symmetry, the duality formulas, the commutative-algebra theorem) is
stated for \lean{VertexOperatorAlgebra}, whose axiom (VOA3), $Y(L_{-1}v,z) = \partial_zY(v,z)$,
supplies exactly the missing ingredient. What was wrong is \emph{documentation}: (i) an earlier
docstring of \leanpath{VOA/\allowbreak Core/\allowbreak BareVA.lean} asserted that the axioms VA2--VA5 are
``all the Borcherds/Jacobi derivation actually consumes'' and planned to restate the Borcherds chain
over \lean{BarePreVA}; this is false, and the plan is a dead end unless a translation axiom is added.
The current docstring has been corrected: it carries a \textbf{Warning} that VA2--VA5 alone do not
imply the Borcherds/Jacobi identity, cites the counterexample of Section~\ref{l03:sec:counter}, and
calls the restatement over \lean{BarePreVA} a dead end.
(ii) An earlier version of the fidelity check of the vertex-algebra definition in Chapter~\ref{ch:01}
called the library's definition ``stronger'' than Kac's and said that locality and Jacobi are equivalent
``for vertex algebras'' (Chapter~\ref{ch:01} now corrects this); the first claim is incorrect (the two definitions are incomparable: the
library's has an extra grading but lacks translation covariance), and the second holds only in the
presence of translation covariance, hence for the library's VOAs but not for its
\lean{PreVertexAlgebra}s. For VOAs, we show (on paper, Section~\ref{l03:sec:compare}) that the
library's definition is exactly the Lepowsky--Li definition plus the CFT-type condition; only the direction
``library VOA $\Rightarrow$ Lepowsky--Li VOA'' is machine-checked.

The release audit of 2026-10-27 (\texttt{Notes/\allowbreak Textbook/\allowbreak referee/\allowbreak RELEASE\_AUDIT\_FOUNDATIONS.md},
finding M-1) recompiled the counterexample and confirmed this finding. Its wording, which this chapter adopts: the
library's \lean{VertexOperatorAlgebra} is the Lepowsky--Li/FLM VOA \textbf{of CFT type} ($V_0 = \C\lthreevac$,
$V_n = 0$ for $n<0$, integer grading, finite-dimensional pieces), so every theorem about \lean{VertexOperatorAlgebra}
is a theorem about VOAs of CFT type only; \lean{PreVertexAlgebra} and \lean{BarePreVA} are \textbf{not} vertex algebras
in the standard sense, and the Borcherds identity, a theorem for \lean{VertexOperatorAlgebra}, fails for them. Where this
chapter says ``vertex algebra in the library's sense'' or ``bare vertex algebra'' it means these weaker structures
(``graded VA1--VA5 data''), never the textbook notion, and never the library's \lean{VertexAlgebra} described next.

\textbf{The library's response: names that say what the structures are.} After this finding the library renamed its
structures (28 October 2026). The structure transcribing Gannon's axioms, formerly called \lean{VertexAlgebra}, is now
\lean{PreVertexAlgebra} (projection \lean{toPreVertexAlgebra}); its ungraded companion, formerly \lean{BareVA}, is
\lean{BarePreVA}. The name \lean{VertexAlgebra} now denotes the genuine notion: a \lean{PreVertexAlgebra} together with
the field \lean{translation_covariant}, which says $Y(Tu,z) = \partial_zY(u,z)$ for the intrinsic translation operator
$Tu = u_{(-2)}\lthreevac$ (file \leanpath{VOA/\allowbreak Core/\allowbreak GenuineVertexAlgebra.lean}). On the conformal
side there are two layers: \lean{GradedVertexOperatorAlgebra} extends \lean{PreVertexAlgebra} by a conformal vector,
the axioms (VOA1)--(VOA3) and truncation from below ($V_n = 0$ for $n\ll0$), which is the FLM/Lepowsky--Li notion
(written with locality in place of the Jacobi identity); \lean{VertexOperatorAlgebra} extends it by the two
CFT-type fields \lean{vacuum_grade} and \lean{positivity}, i.e.\ (VOA4). Every VOA is a vertex algebra in the new sense,
by \lean{VertexOperatorAlgebra.toVertexAlgebra}, whose covariance field is
\lean{VertexOperatorAlgebra.isTranslationCovariant}. The structure theory of this chapter (Borcherds identity,
skew-symmetry, FLM associativity, \dots) is proved for \lean{VertexOperatorAlgebra}; the library does not restate it
for \lean{VertexAlgebra} or for \lean{GradedVertexOperatorAlgebra}.
\end{warning}

\textbf{How the chapter is organised.} Section~\ref{l03:sec:conventions} recalls the conventions
(modes, locality, fields, normal ordering) in the form used below. Section~\ref{l03:sec:axioms}
presents the three structures and Section~\ref{l03:sec:compare} compares them with the literature;
Section~\ref{l03:sec:counter} gives the counterexample. Sections~\ref{l03:sec:vafacts}--\ref{l03:sec:dong}
treat what holds at the level of the weak structures \lean{PreVertexAlgebra}/\lean{BarePreVA}: product fields, the operator product
expansion, Goddard uniqueness and Dong's lemma. Section~\ref{l03:sec:cascade} is the
``reconstruction cascade'' at the VOA level, Section~\ref{l03:sec:borcherds} the full Borcherds
identity, Section~\ref{l03:sec:skew} skew-symmetry, Section~\ref{l03:sec:flm} the FLM associativity
formula, and Section~\ref{l03:sec:comm} commutative vertex algebras.
Sections~\ref{l03:sec:recon}--\ref{l03:sec:jacobilocal} treat the existence side: reconstruction,
tensor lifts, local systems and the localised Jacobi identity. The chapter ends with the referee's
checklist and the notes.

\textbf{The logical skeleton in one paragraph.} At the bare level one has: (a) the iterated
product $Y(u,z_1)Y(v,z_2)x$ lies in $V((z_1))((z_2))$, on which $z_1-z_2$ is not a zero divisor;
(b) hence a field local with every $Y(b,\cdot)$ is determined by what it does to the vacuum
(Goddard); (c) locality is preserved by all field products (Dong). At the VOA level one adds:
(d) $Y(u,z)\lthreevac = e^{zT}u$ (from VOA3); then (e) every field product $Y(u,\cdot)_{(n)}Y(v,\cdot)$
is local with all $Y(b,\cdot)$ by (c) and creates $u_{(n)}v$ by (d), so equals $Y(u_{(n)}v,\cdot)$
by (b) --- this is the \emph{duality} or \emph{$n$-th product theorem}; (f) its $n\ge0$ case gives
the commutator formula, its all-$n$ case gives the iterate formula, and a Pascal recursion glues
these into the Borcherds identity; (g) skew-symmetry is again (b), applied to the field
$v\mapsto e^{zT}Y(v,-z)u$, whose locality is the hard computation.

%======================================================================
\section{Conventions recalled}\label{l03:sec:conventions}
%======================================================================

We use the conventions of Chapter~\ref{ch:01} and Chapter~\ref{ch:L01} throughout. Since every
fidelity check below depends on them, we restate the four that matter most.

\textbf{Distributions are indexed by powers of $z$; modes by the shifted index.} A formal
distribution with coefficients in a space $W$ is a function $f:\Z\to W$, read as
$\sum_kf(k)z^k$ (\lean{VOA.FormalDist}, written \lean{W⟦z^±⟧}). The \emph{mode} is
\lean{mode f n = f (-n - 1)}, so that for $f = Y(u,z) = \sum_nu_{(n)}z^{-n-1}$ one has
\lean{mode (Y u) n} $=u_{(n)}$, while \lean{(Y u) k} is the coefficient of $z^k$, i.e.\
$u_{(-k-1)}$. Both indexings occur in the statements below and we decode each occurrence. The
$n$-th product is \lean{V.smul u n v := mode (V.Y u) n v} $=u_{(n)}v$ (the name \lean{smul} has
nothing to do with scalar multiplication). The derivative is \lean{deriv f k = (k+1) • f (k+1)},
the coefficientwise $\partial_z$; in modes, $(\partial_zf)_{(m)} = -m\,f_{(m-1)}$.

\textbf{Two-variable distributions and locality.} A two-variable distribution is
$F:\Z\to\Z\to W$, $F(m,n)$ being the coefficient of $z_1^mz_2^n$ (\lean{FormalDist₂}).
Multiplication by $z_1-z_2$ is \lean{mulBySub F m n = F (m-1) n - F m (n-1)} and
\lean{mulBySubPow N} is its $N$-th iterate. The commutator of two operator-valued distributions is
\lean{bracketDist a b m n = a m * b n - b n * a m}, and
\begin{equation*}
\text{\lean{Local a b}}\quad:\Longleftrightarrow\quad \exists N\in\N,\ (z_1-z_2)^N[a(z_1),b(z_2)] = 0 .
\end{equation*}
The $j$-th \emph{OPE coefficient} of a two-variable distribution is
$\text{\lean{cCoeff j F}}(w) = \Res_{z_1}(z_1-w)^jF(z_1,w)$, i.e.\ \lean{cCoeff j F n = mulBySubPow j F (-1) n}
(no factorial). Chapter~\ref{ch:L01} proves (\lean{local_iff_deltaSum}, Gannon Prop.~5.1.2) that
$(z_1-z_2)^NF = 0$ if and only if $F = \sum_{j<N}c_j(z_2)\partial_{z_2}^j\delta(z_1/z_2)$ for some
one-variable $c_j$, where the building block \lean{deltaDeriv j c} has
$z_1^mz_2^n$-coefficient $\big(\prod_{i<j}(-m-1-i)\big)\,c(m+n+1+j)$; and then
$c_j = \frac1{j!}\,$\lean{cCoeff j F}.

\textbf{Fields and lower truncation.} The library's predicate is
\begin{leancode}
def IsField (a : FormalDist (End ℂ V)) : Prop :=
  ∀ x : V, ∃ B : ℤ, ∀ n : ℤ, n < B → a n x = 0
\end{leancode}
File \leanpath{VOA/\allowbreak Core/\allowbreak Dong.lean} (\lean{VOA.IsField}). In modes: for each
vector $x$, $a_{(m)}x = 0$ for $m\gg0$, i.e.\ $a(z)x\in V((z))$. This is Kac's notion of field and
Lepowsky--Li's ``truncation''.

\textbf{Normal ordering.} The two-variable normal-ordered product is
\lean{normalOrdered a b m n = if 0 ≤ m then a m * b n else b n * a m}: the coefficients of
\emph{nonnegative} powers of $z_1$ (the modes $a_{(k)}$ with $k\le-1$, which create states) go to
the left. The same-variable product ${:}a(z)b(z){:}$ of two fields is \lean{nprod ha hb}, with
$z^s$-coefficient $\sum_m\text{\lean{normalOrdered a b m (s-m)}}$ (a finite sum on each vector),
and the negative products are
\begin{leancode}
noncomputable def negProd {a b : FormalDist (End ℂ V)} (ha : IsField a) (hb : IsField b)
    (m : ℕ) : FormalDist (End ℂ V) :=
  (m.factorial : ℂ)⁻¹ • nprod (ha.deriv_iterate m) hb
\end{leancode}
File \leanpath{VOA/\allowbreak Core/\allowbreak Dong.lean}. So \lean{negProd ha hb m}
$=\frac1{m!}{:}(\partial_z^ma)(z)\,b(z){:}$, the field product $a_{(-m-1)}b$ of Kac
(§2.3).

\begin{dictionary}
\textbf{Why these conventions are the standard ones.} For $n\ge0$ Kac defines the field product
$a(w)_{(n)}b(w) = \Res_z\big(a(z)b(w)\iota_{z,w}(z-w)^n - b(w)a(z)\iota_{w,z}(z-w)^n\big)$; for
$n\ge0$ the two expansions of $(z-w)^n$ agree (it is a polynomial), so this is
$\Res_z(z-w)^n[a(z),b(w)]$, which is \lean{cCoeff n (bracketDist a b)}. For $n=-m-1<0$ Kac's
formula gives $\frac1{m!}{:}\partial^ma(w)\,b(w){:}$, which is \lean{negProd}. Hence the library's
pair (\lean{cCoeff}, \lean{negProd}) is exactly Kac's family of $n$-th products of fields. The
$z$-power indexing of distributions is Gannon's; the $+1$ shift between powers and modes is the
source of most sign and index decoding in this chapter, and we redo it in every fidelity check.
\end{dictionary}

\begin{physics}
These conventions are the algebraic shadow of the operator product expansion of two chiral fields,
$a(z)b(w) = \sum_{n\ge0}\frac{(a_{(n)}b)(w)}{(z-w)^{n+1}} + {:}a(z)b(w){:}$. The singular terms are
the OPE coefficients $a_{(n)}b$ with $n\ge0$; they are extracted by contour integrals
$\oint_w dz\,(z-w)^n a(z)b(w)$, which is what $\Res_{z}(z-w)^n[a(z),b(w)]$ computes (the commutator
is the difference of the two radially ordered expansions, i.e.\ the contour around $w$). The
regular part at $z=w$ is the normal-ordered product, and its Taylor coefficients are the
$a_{(-m-1)}b = \frac1{m!}{:}\partial^ma\,b{:}$. A field is an operator whose OPE with any state has a
pole of finite order: in radial quantisation, $a(z)$ applied to a state of bounded energy produces
only finitely many states of lower energy.
\end{physics}

%======================================================================
\section{The three structures}\label{l03:sec:axioms}
%======================================================================

\subsection{\texttt{PreVertexAlgebra}: Gannon's Definition 5.1.3(a)}

\begin{definition}[The library's \texttt{PreVertexAlgebra}: graded VA1--VA5 data]\label{l03:def:VA}
A \emph{pre-vertex algebra} (Gannon's ``vertex algebra''; in this book also called, for historical reasons, ``a
vertex algebra in the library's sense''; it is \emph{not} a vertex algebra in the standard sense of Kac, Frenkel--Ben-Zvi or Lepowsky--Li, since it lacks
translation covariance, see Section~\ref{l03:sec:counter}) is a complex vector space $V$ with a
$\Z$-grading $V = \bigoplus_{n\in\Z}V_n$ by finite-dimensional subspaces, a linear map
$Y:V\to\End(V)[[z^{\pm1}]]$, $u\mapsto Y(u,z) = \sum_nu_{(n)}z^{-n-1}$, and a vector
$\lthreevac\in V_0$, such that for all $u,v\in V$:
\begin{enumerate}[label=(VA\arabic*)]
\item $u\in V_k$, $v\in V_\ell$ $\Rightarrow$ $u_{(n)}v\in V_{k+\ell-n-1}$ for all $n\in\Z$;
\item $Y(\lthreevac,z) = \mathrm{id}_V$;
\item $u_{(-1)}\lthreevac = u$ and $u_{(n)}\lthreevac = 0$ for $n\ge0$;
\item $Y(u,z)$ and $Y(v,z)$ are mutually local;
\item there is $N = N(u,v)\in\Z$ with $u_{(n)}v = 0$ for all $n\ge N$.
\end{enumerate}
\end{definition}

\begin{leancode}
structure PreVertexAlgebra (𝒱 : Type u) [AddCommGroup 𝒱] [Module ℂ 𝒱] where
  grading : ℤ → Submodule ℂ 𝒱
  isInternal : DirectSum.IsInternal grading
  finrank : ∀ n, FiniteDimensional ℂ (grading n)
  Y : 𝒱 →ₗ[ℂ] (End ℂ 𝒱)⟦z^±⟧
  vacuum : 𝒱
  vacuum_mem : vacuum ∈ grading 0
  grading_mode : ∀ {k ℓ : ℤ} (u : 𝒱), u ∈ grading k → ∀ (v : 𝒱), v ∈ grading ℓ →
    ∀ (n : ℤ), mode (Y u) n v ∈ grading (k + ℓ - n - 1)
  vacuum_Y : Y vacuum = (1 : (End ℂ 𝒱)⟦z^±⟧)
  creation : ∀ u : 𝒱, mode (Y u) (-1) vacuum = u ∧ ∀ n : ℤ, 0 ≤ n → mode (Y u) n vacuum = 0
  locality : ∀ u v : 𝒱, Local (Y u) (Y v)
  regularity : ∀ u v : 𝒱, ∃ N : ℤ, ∀ n : ℤ, N ≤ n → mode (Y u) n v = 0
\end{leancode}
File \leanpath{VOA/\allowbreak Core/\allowbreak VertexAlgebra.lean} (\lean{VOA.PreVertexAlgebra};
field docstrings removed).

\begin{dictionary}
Chapter~\ref{ch:01} decodes every field; we add only what the audit needs.
\begin{itemize}
\item \textbf{The data} are the grading (three fields: a family of submodules, the internal
  direct-sum property \lean{DirectSum.IsInternal} --- every vector is uniquely a finite sum of
  homogeneous components --- and finite-dimensionality of each piece), the map \lean{Y} and the
  vector \lean{vacuum}. \textbf{There is no translation operator among the data, and no axiom
  mentions a derivative.}
\item \textbf{Truncation is for all pairs}: \lean{regularity} is quantified over all $u,v$,
  and the bound $N$ may depend on both. Applied for fixed $u$ and all $v$ it says that $Y(u,z)$ is
  a field (\lean{isField_Y} below).
\item \textbf{Locality is for all pairs, including $u = v$}, with an order $N$ depending on the
  pair (the existential is inside \lean{Local}).
\item \textbf{Vacuum}: \lean{1} is the library's own unit distribution (coefficient
  $\mathrm{id}$ at $z^0$, zero elsewhere), so \lean{vacuum_Y} is $Y(\lthreevac,z) = \mathrm{id}$,
  not the constant series $\delta(z)$.
\item \textbf{No Jacobi identity} is assumed, and \textbf{no axiom relates $u_{(-2)}\lthreevac$ to
  anything}.
\end{itemize}
How the definition compares with the literature is the subject of Section~\ref{l03:sec:compare}.
\end{dictionary}

\begin{physics}
The five axioms are the minimal chiral-algebra data a physicist writes down: a space of states
graded by scaling dimension with finite degeneracies (VA1), the identity operator corresponding to
the vacuum (VA2), the state--operator map $u = \lim_{z\to0}Y(u,z)\lthreevac$ (VA3), commutativity of
chiral fields at distinct points, i.e.\ single-valuedness of radially ordered correlators
$\langle\cdots a(z)b(w)\cdots\rangle$ up to poles on the diagonal (VA4), and finite-order poles in
the OPE (VA5). What is \emph{missing} is the statement that the field at $z$ is the translate of
the field at $0$, i.e.\ that the Hamiltonian-like operator $L_{-1}$ generating translations acts as
$\partial_z$ on fields. Section~\ref{l03:sec:counter} shows that this is not a cosmetic omission.
\end{physics}

\begin{definition}[The $n$-th product and the translation operator]\label{l03:def:smulT}
In a vertex algebra in the library's sense, $u_{(n)}v := $ \lean{V.smul u n v}, and the
\emph{translation operator} is the linear map $T:V\to V$, $Tu := u_{(-2)}\lthreevac$. One has
$T\lthreevac = 0$ and $T(V_k)\subseteq V_{k+1}$.
\end{definition}
\begin{leancode}
def smul (V : PreVertexAlgebra 𝒱) (u : 𝒱) (n : ℤ) (v : 𝒱) : 𝒱 := mode (V.Y u) n v

def T (V : PreVertexAlgebra 𝒱) : 𝒱 →ₗ[ℂ] 𝒱 where
  toFun u := V.smul u (-2) V.vacuum

@[simp] theorem T_vacuum (V : PreVertexAlgebra 𝒱) : V.T V.vacuum = 0

theorem T_mem_grading (V : PreVertexAlgebra 𝒱) {k : ℤ} {u : 𝒱} (hu : u ∈ V.grading k) :
    V.T u ∈ V.grading (k + 1)
\end{leancode}
File \leanpath{VOA/\allowbreak Core/\allowbreak VertexAlgebra.lean} (\lean{smul}) and
\leanpath{VOA/\allowbreak Core/\allowbreak Identities.lean} (\lean{T}, \lean{T_vacuum},
\lean{T_mem_grading}; linearity proofs omitted).

\begin{dictionary}
$T$ is \emph{defined}, not postulated, and defined by the formula that holds in every textbook
vertex algebra ($Y(u,z)\lthreevac = e^{zT}u$ gives $u_{(-2)}\lthreevac = Tu$). \lean{T_vacuum} is
VA2: $\lthreevac_{(-2)} = \delta_{-2,-1}\mathrm{id} = 0$. \lean{T_mem_grading} is VA1 with
$\lthreevac\in V_0$ and $n=-2$: $k+0+2-1 = k+1$. Nothing about $T$ beyond these two facts holds in
general (Section~\ref{l03:sec:counter}).
\end{dictionary}

\begin{physics}
$T$ is meant to be the generator of translations, $L_{-1}$, which raises the scaling dimension by
one ($\partial_z$ has dimension one) and annihilates the translation-invariant vacuum. In a genuine
CFT $Tu$ is the state of the derivative field $\partial u(z)$.
\end{physics}

\subsection{\texttt{BarePreVA}: the axioms without the grading}

\begin{definition}[Bare pre-vertex algebra (the library's \texttt{BarePreVA})]\label{l03:def:bare}
A \emph{bare pre-vertex algebra} (in this book also ``bare vertex algebra''; not a vertex algebra in the standard sense:
like \lean{PreVertexAlgebra} it has no translation-covariance axiom, and the Borcherds identity fails for it) is a
complex vector space with $Y$ and $\lthreevac$ as above satisfying
(VA2)--(VA5), with no grading. Every pre-vertex algebra gives one by forgetting
the grading.
\end{definition}
\begin{leancode}
structure BarePreVA (𝒱 : Type u) [AddCommGroup 𝒱] [Module ℂ 𝒱] where
  Y : 𝒱 →ₗ[ℂ] (End ℂ 𝒱)⟦z^±⟧
  vacuum : 𝒱
  vacuum_Y : Y vacuum = (1 : (End ℂ 𝒱)⟦z^±⟧)
  creation : ∀ u : 𝒱, mode (Y u) (-1) vacuum = u ∧ ∀ n : ℤ, 0 ≤ n → mode (Y u) n vacuum = 0
  locality : ∀ u v : 𝒱, Local (Y u) (Y v)
  regularity : ∀ u v : 𝒱, ∃ N : ℤ, ∀ n : ℤ, N ≤ n → mode (Y u) n v = 0

def PreVertexAlgebra.toBarePreVA (V : PreVertexAlgebra 𝒱) : BarePreVA 𝒱 where

theorem isField_Y (V : BarePreVA 𝒱) (u : 𝒱) : IsField (V.Y u)
\end{leancode}
File \leanpath{VOA/\allowbreak Core/\allowbreak BareVA.lean} (\lean{VOA.BarePreVA},
\lean{VOA.PreVertexAlgebra.toBarePreVA}, \lean{VOA.BarePreVA.isField_Y}; the body of \lean{toBarePreVA} copies
the five fields). The same file also defines \lean{VertexOperatorAlgebra.toBarePreVA} through the
underlying vertex algebra.

\begin{dictionary}
The fields are those of \lean{PreVertexAlgebra} minus \lean{grading}, \lean{isInternal},
\lean{finrank}, \lean{vacuum_mem}, \lean{grading_mode}, verbatim otherwise. \lean{isField_Y}: VA5
for the pair $(u,x)$ gives $N$ with $u_{(n)}x = 0$ for $n\ge N$; in $z$-power indexing the
coefficient \lean{(Y u) k x} $=u_{(-k-1)}x$ vanishes for $k<-N$, which is \lean{IsField} with
$B = -N$.

\textbf{The file's motivation, and why its original plan fails.} An earlier docstring explained that the
grading is ``dead weight'' for the Borcherds theory --- correctly: none of the Borcherds files uses the
grading --- and that the Borcherds chain should be restated over \lean{BarePreVA} so that it can be
applied to Li's local systems (Section~\ref{l03:sec:local}), which carry no grading. It asserted
that VA2--VA5 are ``all the Borcherds/Jacobi derivation actually consumes''. The current docstring
reverses this: under the heading \textbf{Warning} it says that VA2--VA5 alone do not imply the
Borcherds/Jacobi identity, because neither \lean{PreVertexAlgebra} nor \lean{BarePreVA} has a
translation-covariance axiom; it cites the counterexample of Section~\ref{l03:sec:counter}; it calls the
restatement over \lean{BarePreVA} a dead end; and its status note says that nothing Borcherds-related is
proved over \lean{BarePreVA}. The old sentence was
false: the derivation also consumes translation covariance, which the library takes from the VOA
axiom (VOA3). The counterexample of Section~\ref{l03:sec:counter} is a \lean{PreVertexAlgebra}, hence
(by \lean{toBarePreVA}) a \lean{BarePreVA}, violating the Borcherds identity. The restatement plan can
only succeed if a translation operator with $[T,Y(u,z)] = \partial_zY(u,z)$ and $T\lthreevac = 0$
is added to \lean{BarePreVA}. (For the local-system application this is harmless: a local system has
the natural translation $a\mapsto\partial a$, see \lean{LocalFieldSystem.translation_val} in
Section~\ref{l03:sec:local}.)
\end{dictionary}

\begin{physics}
Removing the grading means forgetting the scaling dimensions, keeping only the operator algebra.
That is the natural setting for, e.g., the algebra of all chiral fields acting on a single sector
(a module), where no vacuum-module grading is available.
\end{physics}

\begin{status}
\lean{BarePreVA} and \lean{toBarePreVA} are definitions; \lean{BarePreVA.isField_Y} is proved,
\texttt{\#print axioms} $=$ \lean{[propext, Classical.choice, Quot.sound]} (run for this
chapter). No Borcherds-type theorem is stated over \lean{BarePreVA} at this commit.
\end{status}

\subsection{\texttt{VertexOperatorAlgebra}}

\begin{definition}[Vertex operator algebra of CFT type, the library's \texttt{VertexOperatorAlgebra}]\label{l03:def:VOA}
A VOA (in the library's sense: a Lepowsky--Li/FLM VOA \emph{of CFT type}) is a structure as in
Definition~\ref{l03:def:VA} with $\omega\in V_2$ and
$c\in\C$ such that, with $L_n := \omega_{(n+1)}$: (VOA1) the Virasoro relations with central charge
$c$; (VOA2) $L_0 = n$ on $V_n$; (VOA3) $Y(L_{-1}v,z) = \partial_zY(v,z)$ for all $v$; (VOA4)
$V_0 = \C\lthreevac$ and $V_n = 0$ for $n<0$.
\end{definition}
Axiom (VOA4) is the CFT-type condition (Lean fields \lean{vacuum_grade} and \lean{positivity}). It is a
\emph{restriction} relative to FLM and Lepowsky--Li, who require only $\dim V_n<\infty$ and $V_n = 0$ for $n\ll0$: it
excludes VOAs with $\dim V_0>1$ (for example non-simple direct sums), VOAs with negative weights (the $\beta\gamma$
system in its standard grading, shifted conformal vectors) and indefinite-lattice examples. Within CFT type the
definition imposes no further restriction. Every theorem the library proves about \lean{VertexOperatorAlgebra} is
therefore a theorem about VOAs of CFT type. Note also that the zero space satisfies every axiom (if
$\lthreevac = 0$, creation forces $V = 0$, and $V_0 = \C\cdot 0 = 0$ is allowed); the literature differs on whether to
admit it, and statements meant for nonzero VOAs must assume $\lthreevac\neq 0$ or \lean{Nontrivial} (release audit
2026-10-27, finding m-4).
The library keeps the two layers apart. \lean{VOA.GradedVertexOperatorAlgebra} extends \lean{PreVertexAlgebra} by
$\omega$, $c$, (VOA1)--(VOA3) and the FLM/Lepowsky--Li truncation $V_n = 0$ for $n\ll 0$ (field
\lean{grading_bounded_below}); \lean{VOA.VertexOperatorAlgebra} extends it by (VOA4), and its truncation field is filled
automatically ($N = 0$). So \lean{VertexOperatorAlgebra} is of CFT type \emph{by definition}, and the general notion is
\lean{GradedVertexOperatorAlgebra}; for the latter the library provides the $L_n$ form of (VOA1)--(VOA3) and the map
\lean{GradedVertexOperatorAlgebra.toNearVOA}, while every theorem of this chapter is stated for
\lean{VertexOperatorAlgebra}. Both structures are printed in Chapter~\ref{ch:01}.
The Lean structure \lean{VOA.VertexOperatorAlgebra} (file
\leanpath{VOA/\allowbreak Core/\allowbreak VertexOperatorAlgebra.lean}) is printed and decoded in
full in Chapter~\ref{ch:01}; for the audit we need only the translation field:
\begin{leancode}
  translation : ∀ v : 𝒱, Y (mode (Y conformalVector) 0 v) = deriv (Y v)
\end{leancode}
File \leanpath{VOA/\allowbreak Core/\allowbreak VertexOperatorAlgebra.lean}: \lean{mode (Y ω) 0}
$=\omega_{(0)} = L_{-1}$, and \lean{deriv} is $\partial_z$; so this is the $L_{-1}$-derivative
property of Lepowsky--Li and FLM, stated as an equality of fields for every $v$.

\begin{physics}
The conformal vector is the stress tensor $T(z) = Y(\omega,z) = \sum L_nz^{-n-2}$. Axiom (VOA3) is
the translation Ward identity: $L_{-1} = \oint T(z)\,dz$ generates $z\mapsto z+\epsilon$, and the
state $L_{-1}v$ corresponds to the field $\partial v(z)$. It is precisely this axiom that the bare
definition lacks, and it is through it that all of the structure theory of this chapter enters.
\end{physics}

(\lean{VOA.NearVOA}, the same axioms without finite-dimensionality and without (VOA4), is described in
Chapter~\ref{ch:01}; no theorem of this chapter is stated for it.)

%======================================================================
\section{Comparison with the textbook definitions}\label{l03:sec:compare}
%======================================================================

The standard definitions differ in which properties are axioms and which are theorems. We
summarise them, then compare.

\begin{itemize}
\item \textbf{Borcherds (1986)}: a vector space with bilinear products $u_{(n)}v$ ($n\in\Z$), a
  vacuum, $u_{(n)}v = 0$ for $n\gg0$, $\lthreevac_{(n)}v = \delta_{n,-1}v$, $u_{(-1)}\lthreevac = u$,
  $u_{(n)}\lthreevac = 0$ ($n\ge0$), and the Borcherds identity in all three indices as the main
  axiom. (Borcherds originally wrote it with a derivation $D$, $Du = u_{(-2)}\lthreevac$, and
  skew-symmetry among the axioms; the identity in all indices implies both.)
\item \textbf{FLM (1988) and Lepowsky--Li (2004, Def.~3.1.1)}: $(V,Y,\lthreevac)$ with
  truncation for all pairs, vacuum property, creation property, and the Jacobi identity
  \eqref{l03:eq:jacobi} for all $u,v$; the $D$-derivative property is a consequence
  (Lepowsky--Li, §3.1). A VOA (LL Def.~3.1.22) adds a $\Z$-grading with finite-dimensional pieces,
  $V_n = 0$ for $n\ll0$, $\omega\in V_2$ with Virasoro relations, $L_0$-grading and the
  $L_{-1}$-derivative property.
\item \textbf{Kac (1998, Def.~1.3) and Frenkel--Ben-Zvi (2004, Def.~1.3.1)}: $(V,\lthreevac,T,Y)$
  with $T\in\End V$ \emph{part of the data}; each $Y(u,z)$ a field; vacuum axiom
  $Y(\lthreevac,z) = \mathrm{id}$, $Y(u,z)\lthreevac\in V[[z]]$ with constant term $u$;
  \emph{translation axiom} $[T,Y(u,z)] = \partial_zY(u,z)$ and $T\lthreevac = 0$; locality. No
  grading.
\item \textbf{Gannon (2006, Def.~5.1.3(a))}: as Definition~\ref{l03:def:VA}, with grading and
  locality, without $T$ and without Jacobi.
\end{itemize}

The Jacobi identity referred to is
\begin{equation}\label{l03:eq:jacobi}
z_0^{-1}\delta\Big(\tfrac{z_1-z_2}{z_0}\Big)Y(u,z_1)Y(v,z_2)
 - z_0^{-1}\delta\Big(\tfrac{z_2-z_1}{-z_0}\Big)Y(v,z_2)Y(u,z_1)
 = z_2^{-1}\delta\Big(\tfrac{z_1-z_0}{z_2}\Big)Y(Y(u,z_0)v,z_2),
\end{equation}
whose $z_0^{-\ell-1}z_1^{-m-1}z_2^{-n-1}$-coefficient is the Borcherds identity
\eqref{l03:eq:borcherds}. (To extract it, multiply by $z_0^{\ell}z_1^{m}z_2^{n}$, take
$\operatorname{Res}_{z_0}\operatorname{Res}_{z_1}\operatorname{Res}_{z_2}$, and expand each $\delta$-function binomially in the
direction fixed by its argument; this is the standard computation behind Kac (4.8.3). The library states and proves only the
coefficient form, Theorem~\ref{l03:thm:borcherds}.)

\begin{longtable}{>{\raggedright\arraybackslash}p{0.27\textwidth}>{\raggedright\arraybackslash}p{0.2\textwidth}>{\raggedright\arraybackslash}p{0.2\textwidth}>{\raggedright\arraybackslash}p{0.23\textwidth}}
\toprule
Ingredient & Kac / FBZ & FLM / LL & library \lean{PreVertexAlgebra} (\lean{VOA})\\
\midrule\endhead
grading, $\dim V_n<\infty$ & no (VA); yes (VOA) & VOA only & required (both)\\
translation $T$ & data, axiom $[T,Y]=\partial Y$ & derived, theorem & derived $Tu = u_{(-2)}\lthreevac$; \textbf{no axiom} (VOA: from (VOA3))\\
truncation, all pairs & yes (``field'') & yes & yes (\lean{regularity})\\
vacuum, creation & yes & yes & yes\\
locality & axiom & theorem & axiom\\
Jacobi / Borcherds & theorem & axiom & \textbf{not implied}; theorem for VOAs, all $(\ell,m,n)$\\
skew-symmetry & theorem & theorem & \textbf{not implied}; theorem for VOAs\\
CFT type & --- & no ($V_n=0$ for $n\ll0$) & VOA: $V_0 = \C\lthreevac$, $V_{<0}=0$\\
\bottomrule
\end{longtable}

\begin{proposition}[The library's VOA is the Lepowsky--Li VOA of CFT type]\label{l03:prop:equiv}
Let $(V,Y,\lthreevac,\omega)$ be a graded vector space with a state--field map, vacuum and
conformal vector. The following are equivalent:
\begin{enumerate}[label=(\alph*)]
\item it is a \lean{VertexOperatorAlgebra} (Definition~\ref{l03:def:VOA});
\item it is a VOA in the sense of Lepowsky--Li Definition~3.1.22, with $V_0 = \C\lthreevac$ and
  $V_n = 0$ for $n<0$.
\end{enumerate}
Direction (a)$\Rightarrow$(b) is formalised (through \lean{jacobi}); direction (b)$\Rightarrow$(a), ``every standard
Lepowsky--Li VOA of CFT type is a library VOA'', is \textbf{proved on paper only}.
\end{proposition}

\begin{physics}
The proposition says that, at the level of chiral algebras with a stress tensor, the library
axiomatises exactly the standard object: a holomorphic CFT chiral algebra whose spectrum of
$L_0$ is nonnegative with a unique vacuum. The extra condition (CFT type) excludes chiral
algebras with negative-dimension states (some non-unitary and logarithmic theories, the
$\beta\gamma$ system in its standard grading) and those with several dimension-zero states.
\end{physics}

\textbf{Proof on paper.} (a)$\Rightarrow$(b): truncation, vacuum, creation, the grading, Virasoro,
$L_0$- and $L_{-1}$-properties are fields of the structure, and $V_n=0$ for $n<0$ is (VOA4). The
Jacobi identity is the content of the theorem \lean{jacobi} (Section~\ref{l03:sec:borcherds}), whose
mode form for all $\ell,m,n$ is equivalent to \eqref{l03:eq:jacobi} coefficient by coefficient.
(b)$\Rightarrow$(a): the only fields of the library structure that are not literally LL axioms are
VA1 (\lean{grading_mode}) and VA4 (\lean{locality}). Locality follows from Jacobi by the
classical argument (Lepowsky--Li, §3.2: multiply \eqref{l03:eq:jacobi} by $z_0^N$ with $N$ beyond the
truncation order of $u_{(n)}v$ and take $\Res_{z_0}$). VA1 follows from the commutator formula with
$u = \omega$, $m = 1$: $[L_0,v_{(n)}] = (L_{-1}v)_{(n+1)} + (L_0v)_{(n)} = (-n-1+\wt v)\,v_{(n)}$
(using $(L_{-1}v)_{(n+1)} = -(n+1)v_{(n)}$, the derivative property), so $v_{(n)}$ shifts the
$L_0$-eigenvalue by $\wt v-n-1$, and the grading is the $L_0$-eigenspace grading by
(VOA2). \qed

\begin{status}
Direction (a)$\Rightarrow$(b) is machine-checked in the sense that every LL axiom is either a field
of \lean{VertexOperatorAlgebra} or the theorem \lean{VOA.VertexOperatorAlgebra.jacobi}
(axiom-clean). Direction (b)$\Rightarrow$(a) is argued on paper only: the library contains no
structure in the Lepowsky--Li form, so there is nothing to state; in particular the release audit of 2026-10-27
(finding m-2) found no Lean declaration deriving \lean{Local} from the Jacobi identity. The claim ``every standard VOA of
CFT type is an instance of \lean{VertexOperatorAlgebra}'' therefore rests on the literature (FLM, Lepowsky--Li), not on
the machine.
\end{status}

\begin{remark}
Dropping (VOA4) on both sides, the same argument identifies \lean{GradedVertexOperatorAlgebra} with the VOA of
Lepowsky--Li Definition~3.1.22 on paper. Neither direction of that statement is formalised: the theorem
\lean{jacobi} is proved only for \lean{VertexOperatorAlgebra}.
\end{remark}

%======================================================================
\section[A counterexample to the Borcherds identity]{A vertex algebra in the library's sense that violates the Borcherds identity}\label{l03:sec:counter}
%======================================================================

\begin{theorem}[Gannon's axioms do not imply translation covariance]\label{l03:thm:counter}
Let $A = \C\lthreevac\oplus\C x\oplus\C y$ be the commutative unital algebra with
$x^2 = xy = y^2 = 0$, graded by $\deg\lthreevac = \deg x = 0$, $\deg y = 2$. For $u = a\lthreevac + bx + cy$
put $\varphi_z(u) := u + bz^2y\in A[z]$ and let $Y(u,z)$ be multiplication by $\varphi_z(u)$:
\begin{equation*}
Y(u,z)v = \varphi_z(u)\,v ,\qquad\text{so}\quad Y(\lthreevac,z) = \mathrm{id},\quad Y(x,z) = x\cdot{} + z^2\,y\cdot{},\quad Y(y,z) = y\cdot{} .
\end{equation*}
Then $(A,Y,\lthreevac)$ with this grading satisfies (VA1)--(VA5), i.e.\ it is a
\lean{PreVertexAlgebra}. But $Tx = x_{(-2)}\lthreevac = 0$ while $\partial_zY(x,z) = 2z\,y\cdot{}\ne0$, so
$Y(Tx,z)\ne\partial_zY(x,z)$; there is \emph{no} operator $T'$ with $[T',Y(u,z)] = \partial_zY(u,z)$
for all $u$; and the Borcherds identity \eqref{l03:eq:borcherds} fails for
$(u,v,w,\ell,m,n) = (x,\lthreevac,\lthreevac,-2,0,-2)$: its left side is $2y$ and its right side is $0$.
Consequently skew-symmetry fails as well: $Y(x,z)\lthreevac = x+z^2y$, whereas
$e^{zT}Y(\lthreevac,-z)x = x$.
\end{theorem}

\begin{physics}
Physically this is an ``operator algebra'' in which the field attached to $x$ at the point $z$ is
\emph{not} the translate of the field at the origin: $Y(x,z) = x + z^2y$ instead of $e^{zL_{-1}}$
acting. All fields commute, so locality is trivially satisfied, and the vacuum reproduces every
state; but there is no consistent notion of moving an operator, so OPE associativity (crossing
symmetry) and Bose exchange symmetry fail. In conformal field theory translation invariance is
always present (it is the $L_{-1}$ Ward identity), so physicists never meet this pathology; it is
a reminder that the bare axioms VA1--VA5 are not a complete axiomatisation of a chiral algebra.
\end{physics}

\textbf{Proof on paper.} The modes of $Y(u,z) = u\cdot{} + bz^2y\cdot{}$ are $u_{(-1)} = u\cdot{}$,
$u_{(-3)} = b\,y\cdot{}$ and $u_{(n)} = 0$ otherwise (the coefficient of $z^k$ is $u_{(-k-1)}$).

(VA2): $\varphi_z(\lthreevac) = \lthreevac$. (VA3): $u_{(-1)}\lthreevac = u$, and $u_{(n)} = 0$ for
$n\ge0$. (VA5): $u_{(n)}v = 0$ for $n\ge0$. (VA4): all $Y(u,z)$ are multiplication operators by
elements of the commutative algebra $A[z]$, so $[Y(u,z_1),Y(v,z_2)] = 0$, local with order $0$.
(VA1): it suffices to check homogeneous $u,v$ and the two nonzero modes. For $n=-1$ the target
degree is $\deg u+\deg v$, and $u\cdot v$ is homogeneous of that degree because $A$ is a graded
algebra ($\deg$ is additive on the nonzero products $\lthreevac\cdot w = w$). For $n = -3$ the target is
$\deg u+\deg v+2$; the mode is $b\,y\cdot{}$ with $b$ the $x$-coefficient of $u$, nonzero only for
$u\in\C x$ (degree $0$), and $y\cdot v$ is nonzero only for $v\in\C\lthreevac$, giving $y$ of degree
$2 = 0+0+2$. The pieces $V_0 = \mathrm{span}(\lthreevac,x)$ and $V_2 = \C y$ are finite-dimensional
and $\lthreevac\in V_0$.

Failure of translation covariance: $Tx = x_{(-2)}\lthreevac$ is the $z^1$-coefficient of
$\varphi_z(x)\lthreevac = x+z^2y$, which is $0$; so $Y(Tx,z) = 0$, while $\partial_zY(x,z)$ has
$z^1$-coefficient $2y\cdot{}\ne0$. For an arbitrary $T'$, the $z^1$-coefficient of $[T',Y(x,z)]$ is
$[T',x_{(-2)}] = [T',0] = 0$, again different from $2y\cdot{}$.

Failure of Borcherds: with $u = x$, $v = w = \lthreevac$, $\ell=-2$, $m=0$, $n=-2$, the left side of
\eqref{l03:eq:borcherds} is $\sum_{i\ge0}(-1)^i\binom{-2}{i}\big(x_{(-2-i)}\lthreevac_{(-2+i)}\lthreevac -
\lthreevac_{(-4-i)}x_{(i)}\lthreevac\big)$. The second term vanishes ($x_{(i)}\lthreevac = 0$ for $i\ge0$),
and $\lthreevac_{(-2+i)}\lthreevac = \delta_{i,1}\lthreevac$; so only $i = 1$ survives, giving
$(-1)\binom{-2}{1}x_{(-3)}\lthreevac = 2y$. The right side is $\sum_i\binom0i(x_{(-2+i)}\lthreevac)_{(-2-i)}\lthreevac
= (x_{(-2)}\lthreevac)_{(-2)}\lthreevac = 0_{(-2)}\lthreevac = 0$. Skew-symmetry: $T = 0$ on $x$ and
$Y(\lthreevac,-z) = \mathrm{id}$. \qed

\begin{dictionary}
\textbf{Why this example is the general mechanism.} For a \emph{commutative} vertex algebra in the
library's sense (all fields commute and have no negative powers of $z$), $Y(u,z) = \varphi_z(u)\cdot{}$
for a map $\varphi_z:A\to A[[z]]$ which the axioms force to be an algebra homomorphism with
$\varphi_0 = \mathrm{id}$ --- a \emph{Hasse--Schmidt derivation}. The Borcherds identity at $m=0$,
$\ell<0$ says $Y(u_{(\ell)}v,z) = Y(u,\cdot)_{(\ell)}Y(v,\cdot)$, which here is the group law
$\varphi_w\circ\varphi_z = \varphi_{z+w}$, i.e.\ $\varphi_z = e^{zT}$ for an ordinary derivation $T$. A
Hasse--Schmidt derivation that is not exponential (here $\varphi_z = \mathrm{id}+z^2D$ with $D(x) = y$, $D(\lthreevac) = D(y) = 0$)
satisfies VA1--VA5 and nothing more. In particular the correspondence behind Gannon's
Theorem~5.1.4 --- a commutative vertex algebra \emph{is} a commutative algebra with a derivation $T$,
via $Y(u,z)v = (e^{zT}u)\,v$ --- fails for Gannon's own definition: the example is commutative, its
product $u_{(-1)}v$ is commutative and associative, its $T$ is $0$, yet $Y(x,z)\neq x\cdot{}$. The
library proves the theorem only for VOAs (Section~\ref{l03:sec:comm}), and only the three
conditions, not the reconstruction of $Y$ from $(A,T)$.

\textbf{The machine check.} For this chapter the example was formalised in Lean outside the library
(not built by \texttt{lake build}, compiled with
\texttt{lake env lean -DmaxSynthPendingDepth=3} at an earlier commit, when this chapter was first written, and not re-run for
this edition; the check is not part of this release) as an inhabitant
\lean{L03Counter.VA : PreVertexAlgebra (Fin 3 → ℂ)} together with two refutations, \lean{jacobi_fails} and
\lean{translation_fails}, all depending only on the axioms \texttt{propext}, \texttt{Classical.choice} and
\texttt{Quot.sound}. The library's own copy,
\lean{VOA.TranslationCovariantCounterexample}, is described in the warning of Section~\ref{l03:sec:intro}.
Here \lean{V = Fin 3 → ℂ} with coordinates $(\lthreevac,x,y)$, \lean{grading 0 = {v | v 2 = 0}},
\lean{grading 2 = {v | v 0 = 0 ∧ v 1 = 0}}, all other pieces $\bot$; \lean{jacobi_fails} is the
negation of the conclusion of \lean{VertexOperatorAlgebra.jacobi} (Theorem~\ref{l03:thm:borcherds})
at fixed indices, with \lean{VA.smul} in place of \lean{W.smul}.
\end{dictionary}

\begin{status}
Machine-checked for this chapter in a scratch file (not in the library, not part of the aggregate
build); axiom-clean. The paper proof is complete and independent of the scratch
file.
\end{status}

\begin{warning}[title={Consequences for the reader}]
\begin{enumerate}
\item \emph{Do not read ``vertex algebra'' in a Lean statement as the textbook notion.} A theorem
  with hypothesis \lean{V : PreVertexAlgebra 𝒱} is a theorem about Gannon-style graded local
  state--field systems, a class strictly larger than that of the vertex algebras of Kac or
  Lepowsky--Li equipped with a compatible grading (every such graded vertex algebra is a
  \lean{PreVertexAlgebra}, but not conversely). The theorems of this chapter
  stated at that level (Goddard uniqueness, the OPE decomposition, Dong's lemma, product fields)
  are true in the larger class, so nothing false is proved.
\item \emph{Theorems that need translation covariance are, correctly, stated for VOAs.} The
  Borcherds identity, skew-symmetry, the duality formulas and Theorem~5.1.4 are all hypothesised
  on \lean{VertexOperatorAlgebra}. In the source, the docstrings of
  \leanpath{VOA/\allowbreak Core/\allowbreak Identities.lean} say this explicitly (``at the bare-VA
  level $Y(Tu,z) = \partial_zY(u,z)$ is circular''); the counterexample shows it is not merely
  circular but false.
\item \emph{Dead end.} Restating the Borcherds chain over \lean{BarePreVA} (the plan of an earlier
  docstring of \leanpath{BareVA.lean}; the current docstring records it as a dead end) cannot succeed
  without adding a translation axiom.
\end{enumerate}
\end{warning}

%======================================================================
\section[The bare level: product fields, cancellation]{What holds at the bare level: product fields and the cancellation principle}\label{l03:sec:vafacts}
%======================================================================

Everything in this and the next three sections is proved for an arbitrary \lean{PreVertexAlgebra}
(and, where noted, for arbitrary distributions on an arbitrary vector space); by
Section~\ref{l03:sec:counter} none of it needs, or yields, translation covariance.

\subsection{The vertex operators are fields}

\begin{theorem}\label{l03:thm:isfield}
In a vertex algebra in the library's sense, every $Y(u,z)$ is a field.
\end{theorem}
\begin{leancode}
theorem isField_Y {𝒱 : Type u} [AddCommGroup 𝒱] [Module ℂ 𝒱] (V : PreVertexAlgebra 𝒱) (u : 𝒱) :
    IsField (V.Y u)
\end{leancode}
File \leanpath{VOA/\allowbreak Core/\allowbreak Reconstruction.lean}
(\lean{VOA.PreVertexAlgebra.isField_Y}).

\begin{dictionary}
Same content as \lean{BarePreVA.isField_Y}; the proof converts VA5 for $(u,x)$ to the $z$-power
indexing (bound $B = -N$). The only hypothesis is the structure.
\end{dictionary}

\begin{physics}
Every chiral operator has finite-order poles against every state; this is what allows normal
ordering ${:}u(z)v(z){:}$ to be defined at all.
\end{physics}

\textbf{Proof on paper.} If $u_{(n)}x = 0$ for $n\ge N$, then the $z^k$-coefficient
$u_{(-k-1)}x$ vanishes for $-k-1\ge N$, i.e.\ $k<-N$ (indeed $k\le-N-1$). \qed

\subsection{The iterated product field and the cancellation principle}

\begin{definition}[Iterated product field]\label{l03:def:yprod}
For $u,v,x\in V$, $Y(u,z_1)Y(v,z_2)x$ is the two-variable distribution whose
$z_1^mz_2^n$-coefficient is $u_{(-m-1)}v_{(-n-1)}x$.
\end{definition}
\begin{leancode}
def Yprod (V : PreVertexAlgebra 𝒱) (u v x : 𝒱) : FormalDist₂ 𝒱 :=
  fun m n => (V.Y u) m ((V.Y v) n x)
\end{leancode}
File \leanpath{VOA/\allowbreak Core/\allowbreak Fields.lean} (\lean{VOA.PreVertexAlgebra.Yprod}).

\begin{dictionary}
\lean{(V.Y u) m} is the coefficient of $z_1^m$ in $Y(u,z_1)$, i.e.\ $u_{(-m-1)}$; similarly for
$v$. No hypothesis; this is a definition, and each coefficient is a single vector (no sum).
\end{dictionary}

\begin{physics}
$Y(u,z_1)Y(v,z_2)x$ is the radially ordered product $|z_1|>|z_2|$ of two operators acting on the
state $x$: the formal-series counterpart of a three-point function
$\langle x'|u(z_1)v(z_2)|x\rangle$ expanded in the region $|z_1|>|z_2|$.
\end{physics}

\begin{definition}[Lower truncation in the second variable]\label{l03:def:bdd2}
A two-variable distribution $F$ is \emph{lower-truncated in $z_2$} if there is $B$ with
$F(m,n) = 0$ for all $m$ and all $n<B$. Such distributions form the space $W((z_1))((z_2))$ \emph{only
in the second variable}; nothing is assumed in $z_1$.
\end{definition}
\begin{leancode}
def IsBddBelow₂ (f : FormalDist₂ 𝒲) : Prop := ∃ B : ℤ, ∀ m n : ℤ, n < B → f m n = 0
\end{leancode}
File \leanpath{VOA/\allowbreak Foundations/\allowbreak IteratedLaurent.lean}
(\lean{VOA.IsBddBelow₂}; treated in Chapter~\ref{ch:L01}).

\begin{physics}
Lower truncation in $z_2$ uniformly in $z_1$ is the formal analogue of analyticity of a correlator
near $z_2 = 0$ in the domain $|z_1|>|z_2|$: the inner operator, acting on a fixed state, has a pole
of bounded order.
\end{physics}

\begin{theorem}[The product field is lower-truncated; cancellation]\label{l03:thm:yprodext}
For all $u,v,x$, $Y(u,z_1)Y(v,z_2)x$ is lower-truncated in $z_2$. Consequently, if
$(z_1-z_2)^NY(u,z_1)Y(v,z_2)x = (z_1-z_2)^NY(u',z_1)Y(v',z_2)x'$ for some $N$, then
$Y(u,z_1)Y(v,z_2)x = Y(u',z_1)Y(v',z_2)x'$.
\end{theorem}
\begin{leancode}
theorem Yprod_bddBelow (V : PreVertexAlgebra 𝒱) (u v x : 𝒱) : IsBddBelow₂ (V.Yprod u v x)

theorem Yprod_ext (V : PreVertexAlgebra 𝒱) {u v x u' v' x' : 𝒱} (N : ℕ)
    (h : mulBySubPow N (V.Yprod u v x) = mulBySubPow N (V.Yprod u' v' x')) :
    V.Yprod u v x = V.Yprod u' v' x'
\end{leancode}
File \leanpath{VOA/\allowbreak Core/\allowbreak Fields.lean}. The engine is the cancellation lemma of
Chapter~\ref{ch:L01}:
\begin{leancode}
theorem eq_zero_of_mulBySubPow_of_bddBelow :
    ∀ (N : ℕ) (f : FormalDist₂ 𝒲), mulBySubPow N f = 0 → IsBddBelow₂ f → f = 0
\end{leancode}
File \leanpath{VOA/\allowbreak Foundations/\allowbreak IteratedLaurent.lean}.

\begin{dictionary}
\lean{Yprod_bddBelow}: the bound is uniform in $m$ (the $z_1$-power) because it comes from VA5 for
the single pair $(v,x)$: if $v_{(n)}x = 0$ for $n\ge N$, the whole $z_2^q$-column vanishes for
$q<-N$. \lean{Yprod_ext}: the hypothesis is one equation between two-variable distributions; the
six vectors are arbitrary. The cancellation lemma is the precise statement ``$z_1-z_2$ is a non-zero
divisor on $W((z_1))((z_2))$''; without the truncation hypothesis it is false
($(z_1-z_2)\delta(z_1/z_2) = 0$).
\end{dictionary}

\begin{physics}
The cancellation principle is the formal version of a familiar analytic fact: two functions
holomorphic in the annulus region $|z_1|>|z_2|$ that agree after multiplication by $(z_1-z_2)^N$
agree. The point of formal calculus is that this fails for distributions supported on the diagonal
($\delta(z_1/z_2)$ is the formal ``contact term''), and truncation in $z_2$ excludes exactly those.
\end{physics}

\textbf{Proof on paper.} Truncation is explained in the dictionary. For cancellation, let
$f\ne0$ be lower-truncated in $z_2$ with $(z_1-z_2)f = 0$, and let $q$ be the least $z_2$-power that
occurs. The $z_1^{m+1}z_2^q$-coefficient of $(z_1-z_2)f$ is $f(m,q) - f(m+1,q-1) = f(m,q)$, so
$f(m,q) = 0$ for all $m$, contradicting the choice of $q$. Induction on $N$ handles $(z_1-z_2)^N$, and
\lean{Yprod_ext} applies this to the difference of the two product fields, which is again
lower-truncated. \qed

\begin{status}
\lean{Yprod_bddBelow}, \lean{Yprod_ext}, \lean{eq_zero_of_mulBySubPow_of_bddBelow}: unconditional;
\texttt{\#print axioms} \lean{VOA.PreVertexAlgebra.Yprod_ext} $=$
\lean{[propext, Classical.choice, Quot.sound]} (run for this chapter).
\end{status}

\subsection{The vacuum expansion}

\begin{lemma}\label{l03:lem:vacexp}
$Y(b,w)\lthreevac$ has no negative powers of $w$, and its $w^0$-coefficient is $b$.
\end{lemma}
\begin{leancode}
theorem Y_apply_vacuum_eq_zero (V : PreVertexAlgebra 𝒱) (b : 𝒱) {n : ℤ} (hn : n < 0) :
    (V.Y b) n V.vacuum = 0

theorem Y_apply_vacuum_zero (V : PreVertexAlgebra 𝒱) (b : 𝒱) : (V.Y b) 0 V.vacuum = b
\end{leancode}
File \leanpath{VOA/\allowbreak Core/\allowbreak Goddard.lean}.

\begin{dictionary}
\lean{(V.Y b) n} with $n<0$ is the mode $b_{(-n-1)}$ with $-n-1\ge0$, killed on $\lthreevac$ by VA3;
\lean{(V.Y b) 0} is $b_{(-1)}$, and $b_{(-1)}\lthreevac = b$. This is all of VA3, re-indexed. The
higher coefficients $b_{(-k-1)}\lthreevac$, $k\ge1$, are \emph{not} determined at this level (in the
counterexample $x_{(-3)}\lthreevac = y$ while $T^2x/2 = 0$).
\end{dictionary}

\begin{physics}
$\lim_{w\to0}Y(b,w)\lthreevac = b$: inserting the operator at the origin creates the state. In a
genuine CFT the full expansion is $Y(b,w)\lthreevac = e^{wL_{-1}}b$; only its constant term is an
axiom here.
\end{physics}

%======================================================================
\section{The operator product expansion}\label{l03:sec:ope}
%======================================================================

\begin{theorem}[OPE decomposition and the structural commutator formula]\label{l03:thm:ope}
For all $u,v$ in a vertex algebra in the library's sense there are $N$ and one-variable
distributions $c_0,\dots,c_{N-1}$ with
\begin{equation*}
[Y(u,z_1),Y(v,z_2)] = \sum_{j<N}c_j(z_2)\,\partial_{z_2}^j\delta(z_1/z_2),
\end{equation*}
and then, for all $a,b\in\Z$,
\begin{equation*}
[u_{(a)},v_{(b)}] = \sum_{j<N}a(a-1)\cdots(a-j+1)\,(c_j)_{(a+b-j)} .
\end{equation*}
\end{theorem}
\begin{leancode}
theorem bracket_deltaSum (V : PreVertexAlgebra 𝒱) (u v : 𝒱) :
    ∃ (N : ℕ) (c : Fin N → FormalDist (End ℂ 𝒱)),
      bracketDist (V.Y u) (V.Y v) = ∑ j : Fin N, deltaDeriv (j : ℕ) (c j)

theorem mode_bracket_ope (V : PreVertexAlgebra 𝒱) (u v : 𝒱) {N : ℕ}
    {c : Fin N → FormalDist (End ℂ 𝒱)}
    (h : bracketDist (V.Y u) (V.Y v) = ∑ j : Fin N, deltaDeriv (j : ℕ) (c j)) (a b : ℤ) :
    ⁅mode (V.Y u) a, mode (V.Y v) b⁆ = ∑ j : Fin N,
      (∏ i ∈ Finset.range (j : ℕ), (a - (i : ℤ))) • mode (c j) (a + b - ((j : ℕ) : ℤ))

theorem mode_commutator_structural (V : PreVertexAlgebra 𝒱) (u v : 𝒱) (a b : ℤ) :
    ∃ (N : ℕ) (c : Fin N → FormalDist (End ℂ 𝒱)),
      bracketDist (V.Y u) (V.Y v) = ∑ j : Fin N, deltaDeriv (j : ℕ) (c j) ∧
      ⁅mode (V.Y u) a, mode (V.Y v) b⁆ = ∑ j : Fin N,
        (∏ i ∈ Finset.range (j : ℕ), (a - (i : ℤ))) • mode (c j) (a + b - ((j : ℕ) : ℤ))
\end{leancode}
File \leanpath{VOA/\allowbreak Core/\allowbreak OPE.lean}.

\begin{dictionary}
\lean{Fin N} is $\{0,\dots,N-1\}$; \lean{deltaDeriv j c} is $c(z_2)\partial_{z_2}^j\delta(z_1/z_2)$
with \emph{no} $1/j!$ (Section~\ref{l03:sec:conventions}); \lean{⁅·,·⁆} is the commutator in
$\End V$; the product $\prod_{i<j}(a-i)$ is the falling factorial $a^{\underline j}$, an integer acting
by \lean{•}. The decomposition is \lean{local_iff_deltaSum} (Chapter~\ref{ch:L01}) applied to VA4.

\textbf{What is not claimed.} The coefficients $c_j$ are \emph{not} identified with
$\frac1{j!}Y(u_{(j)}v,z)$ here; they are only existential. The identification is the duality
theorem \lean{nth_product} (Theorem~\ref{l03:thm:nthproduct}), proved for VOAs. Hence this theorem is
the commutator formula ``up to the identification of the OPE coefficients'', which is why the
library calls it \emph{structural}.
\end{dictionary}

\begin{physics}
This is the singular part of the OPE in its commutator form: the commutator of two chiral fields is
supported on the diagonal, a finite sum of derivatives of the delta function whose coefficients are
the residues $\oint_w(z-w)^j u(z)v(w)\,dz$, i.e.\ the pole coefficients of $u(z)v(w)\sim
\sum_j\frac{(u_{(j)}v)(w)}{(z-w)^{j+1}}$. The mode formula is the familiar recipe ``the OPE poles
determine the commutation relations of the modes'', e.g.\ $J(z)J(w)\sim\frac{k}{(z-w)^2}$ giving
$[J_m,J_n] = km\,\delta_{m+n,0}$.
\end{physics}

\textbf{Proof on paper.} The decomposition is VA4 plus Gannon's Prop.~5.1.2. For the modes: the
$z_1^{-a-1}z_2^{-b-1}$-coefficient of the commutator is $[u_{(a)},v_{(b)}]$; that of
$c(z_2)\partial_{z_2}^j\delta(z_1/z_2) = \sum_{k}c(z_2)\,\partial_{z_2}^j(z_1^kz_2^{-k-1})$ is
$\prod_{i<j}(a-i)$ times the $z_2^{-a-b+j-1}$-coefficient of $c$, i.e.\ $a^{\underline j}c_{(a+b-j)}$ (the
library's formula with $m = -a-1$: $\prod_{i<j}(-m-1-i) = \prod_{i<j}(a-i)$). \qed

\begin{status}
All three unconditional, axiom-clean (\texttt{\#print axioms} \lean{VOA.bracket_deltaSum} run for this
chapter).
\end{status}

%======================================================================
\section{Goddard uniqueness}\label{l03:sec:goddard}
%======================================================================

Goddard's theorem is the engine of everything that follows. Its in-algebra version is form~(i) of
Theorem~\ref{l03:thm:goddard} below; we present all three versions the library proves, because each is used.

\begin{theorem}[Goddard uniqueness, three forms]\label{l03:thm:goddard}
Let $V$ be a vertex algebra in the library's sense, $a\in V$ and $A(z)\in\End(V)[[z^{\pm1}]]$ with
$A(z)\lthreevac = Y(a,z)\lthreevac$ coefficientwise.
\begin{enumerate}[label=(\roman*)]
\item If $A$ is local with every $Y(b,z)$, then $A = Y(a,\cdot)$.
\item It suffices that for every $b$ some $(z_1-z_2)^N$ kills $[A(z_1),Y(b,z_2)]\lthreevac$ (the
  commutator applied to the vacuum only).
\end{enumerate}
(iii) (\emph{Abstract test fields.}) Let $V$ be any vector space with a vector $\lthreevac$, $A$ a
distribution with $A(z)\lthreevac = 0$, and $(G_b)_{b\in V}$ distributions with $G_b(w)\lthreevac\in V[[w]]$
of constant term $b$. If $A$ is local with every $G_b$, then $A = 0$.
\end{theorem}
\begin{leancode}
theorem goddard_uniqueness (V : PreVertexAlgebra 𝒱) (A : FormalDist (End ℂ 𝒱)) (a : 𝒱)
    (hAloc : ∀ b : 𝒱, Local A (V.Y b))
    (hAvac : ∀ m : ℤ, A m V.vacuum = (V.Y a) m V.vacuum) :
    A = V.Y a

theorem goddard_zero (vac : 𝒱) (A : FormalDist (End ℂ 𝒱)) (G : 𝒱 → FormalDist (End ℂ 𝒱))
    (hloc : ∀ b, Local A (G b))
    (hcre0 : ∀ b, (G b) 0 vac = b)
    (hcreneg : ∀ b, ∀ q : ℤ, q < 0 → (G b) q vac = 0)
    (hAvac : ∀ p : ℤ, A p vac = 0)
    (m : ℤ) (b : 𝒱) : A m b = 0
\end{leancode}
File \leanpath{VOA/\allowbreak Core/\allowbreak Goddard.lean} (\lean{VOA.PreVertexAlgebra.goddard_uniqueness},
\lean{VOA.goddard_zero}).
\begin{leancode}
theorem goddard_uniqueness_vac (V : PreVertexAlgebra 𝒱) (A : FormalDist (End ℂ 𝒱)) (a : 𝒱)
    (hAloc : ∀ b : 𝒱, ∃ N : ℕ, mulBySubPow N (applyVec V.vacuum (bracketDist A (V.Y b))) = 0)
    (hAvac : ∀ m : ℤ, A m V.vacuum = (V.Y a) m V.vacuum) :
    A = V.Y a
\end{leancode}
File \leanpath{VOA/\allowbreak Core/\allowbreak WeakAssoc.lean} (\lean{VOA.PreVertexAlgebra.goddard_uniqueness_vac}).
Here \lean{applyVec x F m n = (F m n) x} evaluates an operator-valued two-variable distribution on a
vector (\lean{VOA.applyVec}, \leanpath{VOA/\allowbreak Core/\allowbreak Goddard.lean}).

\begin{dictionary}
\textbf{Hypotheses.} In (i), \lean{hAloc} is locality with every vertex operator and \lean{hAvac}
equality on the vacuum of the coefficient of every power $z^m$. $A$ is not assumed to be a field. In
(ii), \lean{hAloc} is weaker: only $[A(z_1),Y(b,z_2)]\lthreevac$, a $V$-valued two-variable distribution,
is required to be killed. (ii) implies (i) (apply a locality relation to $\lthreevac$) and is what
skew-symmetry uses (Section~\ref{l03:sec:skew}). In (iii) the conclusion \lean{A m b = 0} for all
$m,b$ is $A = 0$; no vertex algebra is present; \lean{hcre0}/\lean{hcreneg} say $G_b(w)\lthreevac\in
b+wV[[w]]$. (iii) is the version needed \emph{before} a vertex algebra exists, in the reconstruction
theorem (Section~\ref{l03:sec:recon}).

\textbf{Strength.} (i) is Kac's uniqueness theorem (Kac, §4.4) and the uniqueness half of the
reconstruction theorem of Frenkel--Ben-Zvi (§4.4); it is stated here with no translation operator, and it is \emph{true} without one:
the proof uses only Lemma~\ref{l03:lem:vacexp}. In the counterexample of
Section~\ref{l03:sec:counter} it holds (the reader may check that the field $A = Y(x,\cdot)+z\,y\cdot{}$,
which is local with everything, does not create $Y(x,z)\lthreevac$: it creates $x+zy+z^2y$).
\end{dictionary}

\begin{physics}
``The state determines the field.'' A local operator in a chiral CFT is uniquely fixed by the state
it creates from the vacuum, because the vacuum is cyclic and separating for local fields; two
local operators creating the same state would differ by a local operator annihilating the vacuum,
which must vanish (a chiral Reeh--Schlieder argument). This is why one may speak of \emph{the} field
of a state.
\end{physics}

\textbf{Proof on paper.} (iii) first. Fix $b$ and put $D(z,w) := A(z)G_b(w)\lthreevac$. Since
$A(z)\lthreevac = 0$, $D = [A(z),G_b(w)]\lthreevac$, so $(z-w)^ND = 0$ by locality. Since
$G_b(w)\lthreevac$ has no negative powers of $w$, $D$ is lower-truncated in $w$ (bound $0$). By the
cancellation lemma $D = 0$, and its $w^0$-coefficient is $A(z)b = 0$. (i) and (ii): apply the same
argument to $D = A(z)Y(b,w)\lthreevac - Y(a,z)Y(b,w)\lthreevac = [A(z),Y(b,w)]\lthreevac -
[Y(a,z),Y(b,w)]\lthreevac$ (using $A\lthreevac = Y(a)\lthreevac$), killed by $(z-w)^{N_1+N_2}$, with $N_2$ the
locality order of $(a,b)$ (VA4); reading off $w^0$ gives $A(z)b = Y(a,z)b$. \qed

\begin{status}
All three unconditional, axiom-clean (\texttt{\#print axioms} of
\lean{VOA.PreVertexAlgebra.goddard_uniqueness} and \lean{VOA.goddard_zero} run for this chapter;
\lean{goddard_uniqueness_vac} is in the dependency cone of \lean{skew}, also clean).
\end{status}

%======================================================================
\section{Dong's lemma and the calculus of local fields}\label{l03:sec:dong}
%======================================================================

The second engine is that locality propagates. All results of this section are about arbitrary
operator-valued distributions on an arbitrary complex vector space $V$ --- no vertex algebra, no
vacuum.

\subsection{Elementary closure properties}

\begin{proposition}\label{l03:prop:localcalc}
Let $a,b$ be distributions on $V$.
\begin{enumerate}[label=(\alph*)]
\item If $a,b$ are local then so are $b,a$, and so are $\partial a,b$ and $a,\partial b$ (and all
  iterated derivatives).
\item The identity field $1(z) = \mathrm{id}\cdot z^0$ is a field, local with every distribution,
  and ${:}a\,1{:} = a$ for every field $a$.
\item If $a,b$ are fields then $\partial a$ and ${:}ab{:}$ are fields.
\end{enumerate}
\end{proposition}
\begin{leancode}
theorem Local.derivL {a b : FormalDist (End ℂ V)} (h : Local a b) : Local (deriv a) b

theorem Local.deriv_iterate_left {a b : FormalDist (End ℂ V)} (h : Local a b) (i : ℕ) :
    Local ((deriv^[i]) a) b

theorem Local.symm {a b : FormalDist (End ℂ V)} (h : Local a b) : Local b a

noncomputable def idField : FormalDist (End ℂ V) := fun j => if j = 0 then LinearMap.id else 0

theorem Local_idField_left (b : FormalDist (End ℂ V)) :
    Local (idField : FormalDist (End ℂ V)) b

theorem nprod_idField {a : FormalDist (End ℂ V)} (ha : IsField a) :
    nprod ha (idField_isField) = a

theorem IsField.nprod {a b : FormalDist (End ℂ V)} (ha : IsField a) (hb : IsField b) :
    IsField (nprod ha hb)
\end{leancode}
File \leanpath{VOA/\allowbreak Core/\allowbreak FieldOps.lean}. And
\begin{leancode}
theorem IsField.deriv {a : FormalDist (End ℂ V)} (ha : IsField a) :
    IsField (VOA.deriv a)
\end{leancode}
File \leanpath{VOA/\allowbreak Core/\allowbreak Dong.lean}.

\begin{dictionary}
\lean{deriv^[i]} is the $i$-fold iterate of $\partial_z$. The right-hand versions
(\lean{Local.derivR}, \lean{Local.deriv_iterate_right}, \lean{Local_idField_right}) are in the same
file. \lean{idField} has coefficient $\mathrm{id}$ at $z^0$: it is the unit field (the same object as
the \lean{1} of VA2, defined separately for an arbitrary space). Locality of $\partial a$ with $b$
costs one extra power: $(z_1-z_2)^{N+1}[\partial a,b] = 0$ if $(z_1-z_2)^N[a,b] = 0$.
\end{dictionary}

\begin{physics}
Derivatives of local fields are local (the derivative $\partial\phi$ of a primary is a descendant
field of the same theory), the identity operator commutes with everything, and normal ordering with
the identity is trivial. These are the bookkeeping facts behind ``the operator content of a chiral
algebra is closed under $\partial$ and normal ordering''.
\end{physics}

\textbf{Proof on paper.} (a) Symmetry: swapping the variables turns $(z_1-z_2)^N$ into
$(-1)^N(z_1-z_2)^N$ and $[a,b]$ into $-[b,a]$. Derivative: differentiate $(z_1-z_2)^N[a(z_1),b(z_2)] = 0$
in $z_1$ to get $(z_1-z_2)^N[\partial a,b] = -N(z_1-z_2)^{N-1}[a,b]$, and multiply by $z_1-z_2$.
(b) $\mathrm{id}$ is central; ${:}a1{:}_s = \sum_m{:}a_m1_{s-m}{:} = a_s$ since only $m = s$
contributes. (c) The truncation bound of $\partial a$ shifts by one; for ${:}ab{:}$ on a vector $x$ the
finitely many $a_m$ with $m<0$ that do not kill $x$, and the truncation of $b$ on finitely many vectors,
give a bound. \qed

\begin{status}
All unconditional and axiom-clean (in the dependency cone of \lean{gen_local} and
\lean{reconstruct}, both checked for this chapter).
\end{status}

\subsection{Dong's lemma}

\begin{theorem}[Dong's lemma, all products]\label{l03:thm:dong}
Let $a,b,c$ be pairwise local distributions on $V$. Then for every $k\ge0$ the field product
$a_{(k)}b := \Res_{z_1}(z_1-z_2)^k[a(z_1),b(z_2)]$ is local with $c$. If moreover $a,b$ are fields,
then ${:}ab{:}$ is local with $c$, and hence so are all negative products
$a_{(-m-1)}b = \frac1{m!}{:}(\partial^ma)b{:}$.
\end{theorem}
\begin{leancode}
theorem jacobi₃ (a b c : FormalDist (End ℂ V)) :
    outerC₁₂ (bracketDist a b) c
      = outerA a (bracketDist b c) - outerB b (bracketDist a c)

theorem dong_core (a b c : FormalDist (End ℂ V)) {r : ℕ}
    (hab : mulBySubPow r (bracketDist a b) = 0)
    (hac : mulBySubPow r (bracketDist a c) = 0)
    (hbc : mulBySubPow r (bracketDist b c) = 0)
    (k M : ℕ) (hM : 2 * r ≤ M) :
    mulBySub₂₃Pow (M + r) (mulBySub₁₂Pow k (outerC₁₂ (bracketDist a b) c)) = 0

theorem dong_nonneg {a b c : FormalDist (End ℂ V)}
    (hab : Local a b) (hac : Local a c) (hbc : Local b c) (k : ℕ) :
    Local (cCoeff k (bracketDist a b)) c

theorem dong_neg {a b c : FormalDist (End ℂ V)} (ha : IsField a) (hb : IsField b)
    (hab : Local a b) (hac : Local a c) (hbc : Local b c) :
    Local (nprod ha hb) c
\end{leancode}
File \leanpath{VOA/\allowbreak Core/\allowbreak Dong.lean}.

\begin{dictionary}
\textbf{Three-variable objects.} \lean{FormalDist₃ W} is $\Z^3\to W$, the coefficient of
$z_1^mz_2^nz_3^p$ (\leanpath{VOA/\allowbreak Foundations/\allowbreak FormalDist3.lean}).
\lean{outerC₁₂ X c} is $[X(z_1,z_2),c(z_3)]$, \lean{outerA a X} is $[a(z_1),X(z_2,z_3)]$,
\lean{outerB b X} is $[b(z_2),X(z_1,z_3)]$; \lean{mulBySub₁₂Pow k}, \lean{mulBySub₂₃Pow k} multiply
by $(z_1-z_2)^k$, $(z_2-z_3)^k$. So \lean{jacobi₃} is the operator Jacobi identity
$[[a,b],c] = [a,[b,c]] - [b,[a,c]]$ with variables attached, and \lean{dong_core} says
$(z_2-z_3)^{M+r}(z_1-z_2)^k[[a(z_1),b(z_2)],c(z_3)] = 0$ once $M\ge2r$, with $r$ a common locality
order. \lean{dong_nonneg} is the $k\ge0$ case of the theorem (\lean{cCoeff k} as in
Section~\ref{l03:sec:conventions}); \lean{dong_neg} the normal-ordered case. The negative products
follow by combining \lean{dong_neg} with \lean{Local.deriv_iterate_left} (this combination is
\lean{Fld.prod_local_right_of_pairwise}, Theorem~\ref{l03:thm:fldlaws}).

\textbf{Fidelity.} This is Kac's Lemma~3.2 (Dong's lemma) for all $n\in\Z$. The hypotheses are
exactly pairwise locality, plus the field property where it is needed to form the product.
\end{dictionary}

\begin{physics}
Dong's lemma is the closure of local fields under the OPE: if $a$, $b$, $c$ are mutually local
chiral fields, then every coefficient of the OPE $a(z)b(w)$ --- singular or regular --- is a chiral
field local with $c$. It is why the composite operators of a CFT (${:}JJ{:}$, $\partial J$,
${:}TT{:}$, \dots) are again legitimate chiral fields with the same commutativity properties as their
constituents.
\end{physics}

\textbf{Proof on paper.} Let $r$ bound the three locality orders and $M\ge2r$. By the operator
Jacobi identity,
\begin{equation*}
(z_2-z_3)^{M+r}(z_1-z_2)^k[[a_1,b_2],c_3] = (z_2-z_3)^{M+r}(z_1-z_2)^k\big([a_1,[b_2,c_3]] - [b_2,[a_1,c_3]]\big),
\end{equation*}
with subscripts naming the variables. The first term dies because $(z_2-z_3)^r$ passes through
$a(z_1)$ and kills $[b_2,c_3]$. For the second write $(z_2-z_3)^{M} = \big((z_2-z_1)+(z_1-z_3)\big)^{M} =
\sum_i\binom Mi(z_2-z_1)^{i}(z_1-z_3)^{M-i}$: terms with $M-i\ge r$ die by locality of $a,c$ (the factor
passes through $b(z_2)$); the others have $i>M-r\ge r$, and $(z_2-z_1)^i$ together with the identity
$[b_2,[a_1,c_3]] = [a_1,[b_2,c_3]] - [[a_1,b_2],c_3]$ splits into a term killed by $(z_1-z_2)^r$
(locality of $a,b$) and a term killed by the reserve factor $(z_2-z_3)^r$ (locality of $b,c$). This
proves \lean{dong_core}. Taking $\Res_{z_1}$ commutes with multiplication by $(z_2-z_3)$ and turns
$(z_1-z_2)^k[[a_1,b_2],c_3]$ into $[(a_{(k)}b)(z_2),c(z_3)]$, giving \lean{dong_nonneg} with order
$3r$. For \lean{dong_neg}, one first proves $(z_1-z_3)^r(z_2-z_3)^{3r}[{:}a(z_1)b(z_2){:},c(z_3)] = 0$
by splitting ${:}a_1b_2{:} = a_+(z_1)b(z_2) + b(z_2)a_-(z_1)$ and reusing \lean{dong_core} for the
part that is a genuine double commutator; then one restricts to the diagonal $z_1 = z_2$, which is
legitimate on each vector because $a,b$ are fields (\lean{diag3}, a finite sum), and both factors
become powers of $(z-z_3)$: order $4r$. \qed

\begin{status}
\lean{dong_nonneg}, \lean{dong_neg}: unconditional; \texttt{\#print axioms} $=$
\lean{[propext, Classical.choice, Quot.sound]} for both (run for this chapter).
\end{status}

%======================================================================
\section{The reconstruction cascade in a VOA}\label{l03:sec:cascade}
%======================================================================

From here on $W$ is a \lean{VertexOperatorAlgebra}. The single new input is (VOA3). We follow the
dependency order of \leanpath{VOA/\allowbreak Core/\allowbreak Identities.lean} and
\leanpath{VOA/\allowbreak Core/\allowbreak Reconstruction.lean}.

\subsection{Creation: \texorpdfstring{$Y(u,z)\lthreevac = e^{zT}u$}{Y(u,z)1 = exp(zT)u}}

\begin{theorem}[Derivative property and creation]\label{l03:thm:creation}
In a VOA: $T = L_{-1}$; $Y(T^ju,z) = \partial_z^jY(u,z)$ for all $j\in\N$; and
$u_{(-k-1)}\lthreevac = \frac1{k!}T^ku$ for all $k\in\N$.
\end{theorem}
\begin{leancode}
theorem T_eq_Lneg1 (W : VertexOperatorAlgebra 𝒱) : W.toPreVertexAlgebra.T = W.L (-1)

theorem Y_Tpow (W : VertexOperatorAlgebra 𝒱) (u : 𝒱) (j : ℕ) :
    W.Y ((W.toPreVertexAlgebra.T ^ j) u) = (deriv^[j]) (W.Y u)

theorem creation_pow (W : VertexOperatorAlgebra 𝒱) (a : 𝒱) (k : ℕ) :
    W.smul a (-(k : ℤ) - 1) W.vacuum
      = ((k.factorial : ℂ)⁻¹) • ((W.toPreVertexAlgebra.T ^ k) a)
\end{leancode}
File \leanpath{VOA/\allowbreak Core/\allowbreak Identities.lean}.

\begin{dictionary}
\lean{W.toPreVertexAlgebra.T} is the translation operator of Definition~\ref{l03:def:smulT} of the
underlying \lean{PreVertexAlgebra}; \lean{W.L (-1)} $=\omega_{(0)}$. \lean{T ^ j} is the $j$-th power in
$\End V$, \lean{deriv^[j]} the $j$-th iterate of $\partial_z$. \lean{-(k : ℤ) - 1} is the mode
index $-k-1$. The short proof of $T = L_{-1}$ is the proof on paper after Theorem~\ref{l03:thm:TL}; the point for us is that \emph{every step
uses (VOA3)}, and that the counterexample of Section~\ref{l03:sec:counter} violates all three
conclusions (with $T = 0$ and $x_{(-3)}\lthreevac = y$).
\end{dictionary}

\begin{physics}
The vacuum is translation invariant and the state created at $z$ is the translate of the state
created at $0$: $u(z)|0\rangle = e^{zL_{-1}}u(0)|0\rangle$. This is the translation Ward identity on
the vacuum module.
\end{physics}

\begin{status}
Unconditional, axiom-clean (\texttt{\#print axioms} run for \lean{T_eq_Lneg1},
\lean{creation_pow}, \lean{Y_Tpow}).
\end{status}

\begin{theorem}[The derivative property $Y(Tu,z) = \partial_zY(u,z)$]\label{l03:thm:TL}
In a VOA, for every $u\in V$: $Y(Tu,z) = \partial_zY(u,z)$, where $T = L_{-1}$ is the translation operator of
Definition~\ref{l03:def:smulT}. In modes, $(Tu)_{(n)} = -n\,u_{(n-1)}$ for all $n\in\Z$.
\end{theorem}
\noindent File \leanpath{VOA/\allowbreak Core/\allowbreak Identities.lean}.
\begin{leancode}
theorem Y_T (W : VertexOperatorAlgebra 𝒱) (u : 𝒱) :
    W.Y (W.toPreVertexAlgebra.T u) = deriv (W.Y u)

theorem mode_Y_T (W : VertexOperatorAlgebra 𝒱) (u : 𝒱) (n : ℤ) :
    mode (W.Y (W.toPreVertexAlgebra.T u)) n = (-n) • mode (W.Y u) (n - 1)
\end{leancode}

\begin{dictionary}
\lean{W.toPreVertexAlgebra.T} is the translation operator $Tu = u_{(-2)}\lthreevac$ of the underlying \lean{PreVertexAlgebra}
(Definition~\ref{l03:def:smulT}); \lean{W.Y} is the state--field map; \lean{deriv} is the formal derivative
$\partial_z$ of Definition~\ref{l01:def:mode} (Chapter~\ref{ch:L01}), so \lean{deriv (W.Y u)} is
$\sum_n(-n-1)u_{(n)}z^{-n-2}$; the equality is an equality of formal distributions, i.e.\ of all coefficients.
\lean{mode f n} is the coefficient $f_{-n-1}$ (the library convention of Chapter~\ref{ch:L01}), and
\lean{(-n) • ...} is multiplication by the integer $-n$ cast to $\C$. \textbf{Hypotheses:} none beyond ``$W$ is a
VOA''. \textbf{Same, stronger or weaker?} The statement is the textbook one, Gannon (5.1.8a) and Kac's translation
axiom: the field of $Tu$ is the derivative of the field of $u$. What deserves attention is where the statement comes
from. The weaker structure \lean{PreVertexAlgebra} does not imply it (the counterexample of
Theorem~\ref{l03:thm:counter} has $T=0$ but a field that is not constant); for a VOA the statement is the axiom (VOA3),
which is written with $L_{-1}$, combined with the identification $T = L_{-1}$ (Theorem~\ref{l03:thm:creation}). \textbf{Could it be vacuous?} No: it is an
equation between nonzero fields in every nonzero VOA (for the Heisenberg VOA, $T J = J_{(-2)}\lthreevac$ and
$Y(TJ,z) = \partial_zJ(z)$).
\end{dictionary}

\begin{physics}
Derivatives of a field are generated by the stress tensor: $[L_{-1},\phi(z)] = \partial\phi(z)$, and, since
$L_{-1}u = \partial u$ is the state of the derivative operator, the field of that state is the derivative of
the field. This is the ``state--operator'' form of the translation Ward identity in a chiral CFT. In modes it is
the familiar rule that differentiating $\sum_nu_{(n)}z^{-n-1}$ multiplies the $n$-th mode by $-n$ and lowers its
index, e.g.\ $\partial J(z)=\sum_n(-n)J_{(n-1)}z^{-n-1}$ for a current.
\end{physics}

\textbf{Proof on paper.} By the axiom (VOA3), $Y(L_{-1}u,z) = \partial_zY(u,z)$, so it suffices to show
$T = L_{-1}$. Evaluate (VOA3) on the vacuum and use creation $Y(v,z)\lthreevac\in v+zV[[z]]$: the $z^0$-coefficient of
$Y(L_{-1}u,z)\lthreevac$ is $L_{-1}u$, while the $z^0$-coefficient of $\partial_zY(u,z)\lthreevac$ is the
$z^1$-coefficient of $Y(u,z)\lthreevac$, which is $u_{(-2)}\lthreevac = Tu$. Hence $L_{-1}u = Tu$. The mode form
is obtained by comparing the coefficient of $z^{-n-1}$ on both sides of $Y(Tu,z) = \partial_zY(u,z)$:
$\partial_z\sum_ku_{(k)}z^{-k-1} = \sum_k(-k-1)u_{(k)}z^{-k-2} = \sum_n(-n)u_{(n-1)}z^{-n-1}$ after the substitution $k=n-1$. \qed

\begin{status}
Unconditional. \texttt{\#print axioms} \lean{VOA.VertexOperatorAlgebra.Y\_T} $=$
\lean{[propext, Classical.}\allowbreak\lean{choice, Quot.}\allowbreak\lean{sound]} (run for this chapter at the
commit of the front matter, with \texttt{lake env lean -DmaxSynthPendingDepth=3}); \lean{mode_Y_T} proved from it by
rewriting. The PBW theorem \lean{VOA.pbw} (Chapter~\ref{ch:L02}; formerly a project axiom, now a theorem) is not used.
\end{status}

\subsection{Translation covariance}

\begin{theorem}[Translation covariance]\label{l03:thm:transcov}
In a VOA, $[L_{-1},Y(u,z)] = Y(L_{-1}u,z) = \partial_zY(u,z)$. In modes:
$[L_{-1},u_{(n)}] = (L_{-1}u)_{(n)} = -n\,u_{(n-1)}$.
\end{theorem}
\begin{leancode}
theorem translation_covariance (W : VertexOperatorAlgebra 𝒱) (u : 𝒱) :
    cCoeff 0 (bracketDist (W.Y W.conformalVector) (W.Y u)) = W.Y (W.L (-1) u)

theorem L_neg_one_deriv (W : VertexOperatorAlgebra 𝒱) (u w : 𝒱) (n : ℤ) :
    W.L (-1) (W.smul u n w) - W.smul u n (W.L (-1) w) = W.smul (W.L (-1) u) n w

theorem L_comm_mode (W : VertexOperatorAlgebra 𝒱) (u : 𝒱) (n : ℤ) :
    ⁅W.L (-1), mode (W.Y u) n⁆ = (-n) • mode (W.Y u) (n - 1)
\end{leancode}
File \leanpath{VOA/\allowbreak Core/\allowbreak Reconstruction.lean}.

\begin{dictionary}
\lean{cCoeff 0 (bracketDist (Y ω) (Y u))} is $\Res_{z_1}[Y(\omega,z_1),Y(u,z_2)] = [\omega_{(0)},Y(u,z_2)]
= [L_{-1},Y(u,z_2)]$ (library lemma \lean{cCoeff_zero_omega_apply}). Note that (VOA3) is
$Y(L_{-1}u) = \partial Y(u)$, an equality of the field \emph{of a state}; translation covariance is the
\emph{operator} statement $[L_{-1},Y(u)] = \partial Y(u)$, which Kac and FBZ take as their axiom. It
is a theorem here, and the combination of the two is Kac's axiom. \lean{L_neg_one_deriv} is the
$z^{-n-1}$-coefficient applied to $w$; \lean{L_comm_mode} rewrites $(L_{-1}u)_{(n)} = (Tu)_{(n)} =
-n\,u_{(n-1)}$.
\end{dictionary}

\begin{physics}
The stress tensor generates translations of every field: $[L_{-1},\phi(z)] = \partial\phi(z)$, the
conformal Ward identity for $z\mapsto z+\epsilon$, obtained by contour-integrating $T(w)\phi(z)$
around $z$.
\end{physics}

\textbf{Proof on paper.} The candidate field $A := [L_{-1},Y(u,\cdot)] = \Res_{z_1}[Y(\omega,z_1),Y(u,\cdot)]$
is local with every $Y(b,\cdot)$ by Dong's lemma (\lean{dong_nonneg} with $k = 0$, applied to the
pairwise local $Y(\omega),Y(u),Y(b)$). On the vacuum, $A(z)\lthreevac = L_{-1}Y(u,z)\lthreevac$ because
$L_{-1}\lthreevac = T\lthreevac = 0$; and $L_{-1}(u_{(s)}\lthreevac) = (L_{-1}u)_{(s)}\lthreevac$ for all $s$
(both are $T^{j+1}u/j!$ for $s = -j-1$ by \lean{creation_pow}, and $0$ for $s\ge0$). So $A$ creates
$Y(L_{-1}u,z)\lthreevac$, and Goddard (Theorem~\ref{l03:thm:goddard}(i)) gives $A = Y(L_{-1}u,\cdot)$. \qed

\begin{status}
Unconditional, axiom-clean (\texttt{\#print axioms} \lean{translation_covariance}
run for this chapter).
\end{status}

\subsection{Duality for \texorpdfstring{$n\ge0$}{n >= 0} and the commutator formula}

\begin{theorem}[The $n$-th product theorem, $n\ge0$]\label{l03:thm:nthproduct}
In a VOA, for all $u,v$ and $k\in\N$:
\begin{equation*}
Y(u_{(k)}v,z) = \Res_{z_1}(z_1-z)^k[Y(u,z_1),Y(v,z)] .
\end{equation*}
\end{theorem}
\begin{leancode}
theorem cCoeff_vacuum_eq (W : VertexOperatorAlgebra 𝒱) (u v : 𝒱) (k : ℕ) (m : ℤ) :
    cCoeff k (bracketDist (W.Y u) (W.Y v)) m W.vacuum = (W.Y (W.smul u k v)) m W.vacuum

theorem nth_product (W : VertexOperatorAlgebra 𝒱) (u v : 𝒱) (k : ℕ) :
    cCoeff k (bracketDist (W.Y u) (W.Y v)) = W.Y (W.smul u k v)

theorem assoc_nonneg (W : VertexOperatorAlgebra 𝒱) (u v : 𝒱) (k : ℕ) :
    W.Y (W.smul u k v) = cCoeff k (bracketDist (W.Y u) (W.Y v))
\end{leancode}
File \leanpath{VOA/\allowbreak Core/\allowbreak Reconstruction.lean}.

\begin{dictionary}
\lean{W.smul u k v} with \lean{k : ℕ} coerced to $\Z$ is $u_{(k)}v$. \lean{assoc_nonneg} is
\lean{nth_product} read backwards (Gannon (5.1.9b)). This is Kac's ``$n$-th product identity'' $Y(a_{(n)}b,z) =
Y(a,z)_{(n)}Y(b,z)$ (Kac, §4.6) for $n\ge0$: the state--field map is a homomorphism for the nonnegative
products.
\end{dictionary}

\begin{physics}
The field of the state $u_{(k)}v$ --- the coefficient of $(z-w)^{-k-1}$ in the OPE $u(z)v(w)$ --- is
the corresponding OPE coefficient \emph{field}: OPE coefficients of states and of operators agree.
This is the state--operator correspondence applied to the singular part of the OPE.
\end{physics}

\textbf{Proof on paper.} The left field is local with every $Y(b,\cdot)$ by Dong's lemma. It
creates $u_{(k)}v$: on the vacuum its $z^m$-coefficient vanishes for $m<0$ (all modes $v_{(q)}$,
$q\ge0$, kill $\lthreevac$, and $(z_1-z)^k$ has no negative powers of $z$), its $z^0$-coefficient is
$\Res_{z_1}(z_1)^k Y(u,z_1)v = u_{(k)}v$, and the higher coefficients follow by induction from
$[L_{-1},\cdot\,]$ acting as $\partial_z$ on both sides (\lean{L_comm_cCoeff},
\lean{L_comm_Ycoeff}) together with $L_{-1}\lthreevac = 0$. This is \lean{cCoeff_vacuum_eq}; Goddard
finishes. \qed

\begin{status}
Unconditional, axiom-clean.
\end{status}

\begin{theorem}[Commutator formula]\label{l03:thm:comm}
In a VOA, for all $u,v$ and $a,b\in\Z$,
$[u_{(a)},v_{(b)}] = \sum_{j\ge0}\binom aj(u_{(j)}v)_{(a+b-j)}$ (a finite sum).
\end{theorem}
\begin{leancode}
theorem mode_commutator (W : VertexOperatorAlgebra 𝒱) (u v : 𝒱) (a b : ℤ) :
    ∃ N : ℕ, ⁅mode (W.Y u) a, mode (W.Y v) b⁆
      = ∑ j : Fin N, gBinom a (j : ℕ) • mode (W.Y (W.smul u (j : ℕ) v)) (a + b - (j : ℕ))

theorem commBound (W : VertexOperatorAlgebra 𝒱) (u v : 𝒱) (m n : ℤ) :
    ∃ N : ℕ, (∀ k : ℕ, N ≤ k → W.smul u k v = 0)
      ∧ ⁅mode (W.Y u) m, mode (W.Y v) n⁆
        = ∑ j : Fin N, gBinom m (j : ℕ) • mode (W.Y (W.smul u (j : ℕ) v)) (m + n - (j : ℕ))
\end{leancode}
File \leanpath{VOA/\allowbreak Core/\allowbreak Reconstruction.lean} (\lean{mode_commutator}) and
\leanpath{VOA/\allowbreak Core/\allowbreak Borcherds.lean} (\lean{commBound}). The theorem
\lean{jacobi_commutator} of \leanpath{VOA/\allowbreak Core/\allowbreak BorcherdsIdentity.lean} is
\lean{mode_commutator} under another name.

\begin{dictionary}
\lean{gBinom a j} is $\binom aj = a(a-1)\cdots(a-j+1)/j!$ for $a\in\Z$ (Chapter~\ref{ch:L01}, Definition~\ref{l01:def:gbinom}).
\textbf{The existential.} \lean{mode_commutator} asserts the identity for \emph{some} $N$; read
alone it does not say that the omitted terms $j\ge N$ vanish, so it is formally weaker than the
textbook ``$\sum_{j\ge0}$''. \lean{commBound} closes the gap: the same $N$ also satisfies
$u_{(k)}v = 0$ for all $k\ge N$, so the finite sum \emph{is} the full sum. Together they are exactly
the commutator formula of Gannon (5.1.7c), Kac (§4.6) and Lepowsky--Li (§3.1).
\end{dictionary}

\begin{physics}
``Commutators of modes are determined by the singular part of the OPE.'' For the Virasoro field,
$T(z)T(w)\sim\frac{c/2}{(z-w)^4}+\frac{2T(w)}{(z-w)^2}+\frac{\partial T(w)}{z-w}$ produces the Virasoro
relations; for currents $J^a(z)J^b(w)\sim\frac{k\delta^{ab}}{(z-w)^2}+\frac{f^{abc}J^c(w)}{z-w}$ the
affine Kac--Moody relations.
\end{physics}

\textbf{Proof on paper.} Theorem~\ref{l03:thm:ope} gives the structural formula with coefficients
$c_j$; the delta-decomposition determines $c_j = \frac1{j!}\Res_{z_1}(z_1-z)^j[Y(u,z_1),Y(v,z)]$
(\lean{cCoeff_deltaSum}), which is $\frac1{j!}Y(u_{(j)}v,z)$ by Theorem~\ref{l03:thm:nthproduct}; and
$a^{\underline j}/j! = \binom aj$. For \lean{commBound}, the same decomposition shows
$\Res_{z_1}(z_1-z)^k[\cdots] = 0$ for $k\ge N$, i.e.\ $Y(u_{(k)}v) = 0$, and creation gives $u_{(k)}v = 0$. \qed

\begin{status}
Both unconditional, axiom-clean.
\end{status}

\subsection{Duality for \texorpdfstring{$n<0$}{n < 0}}

\begin{theorem}[The $n$-th product theorem, $n<0$]\label{l03:thm:assocneg}
In a VOA, for all $u,v$ and $k\in\N$,
\begin{equation*}
Y(u_{(-k-1)}v,z) = \frac1{k!}\,{:}\partial_z^kY(u,z)\,Y(v,z){:}.
\end{equation*}
\end{theorem}
\begin{leancode}
theorem assoc_neg (W : VertexOperatorAlgebra 𝒱) (u v : 𝒱) (k : ℕ) :
    W.Y (W.smul u (-(k:ℤ) - 1) v) = negProd (W.isField_Y u) (W.isField_Y v) k
\end{leancode}
File \leanpath{VOA/\allowbreak Core/\allowbreak Reconstruction.lean}.

\begin{dictionary}
\lean{negProd (W.isField_Y u) (W.isField_Y v) k} is $\frac1{k!}{:}(\partial^kY(u))Y(v){:}$
(Section~\ref{l03:sec:conventions}); the two \lean{isField_Y} arguments are the proofs that the
vertex operators are fields, needed to form the normal-ordered product. With $k = 0$:
$Y(u_{(-1)}v,z) = {:}Y(u,z)Y(v,z){:}$, Gannon (5.1.9a), Kac (§4.6). Together with
Theorem~\ref{l03:thm:nthproduct}: $Y(u_{(n)}v,z) = Y(u,z)_{(n)}Y(v,z)$ for \emph{all} $n\in\Z$.
\end{dictionary}

\begin{physics}
The regular part of the OPE: the field of the state $u_{(-k-1)}v$ is the $k$-th Taylor coefficient of
$u(z)v(w)$ at $z = w$ after subtracting the poles, i.e.\ ${:}\partial^ku\,v{:}/k!$. For instance the
Sugawara or free-boson stress tensor is the field of the state $\frac12 J_{(-1)}J$, namely
$\frac12{:}JJ{:}$.
\end{physics}

\textbf{Proof on paper.} Replace $\partial^kY(u)$ by $Y(T^ku)$ (Theorem~\ref{l03:thm:creation}). The
field $\frac1{k!}{:}Y(T^ku)Y(v){:}$ is local with every $Y(b)$ by \lean{dong_neg}. On the vacuum,
its negative $z$-powers vanish (in each normal-ordered term the factor acting first on
$\lthreevac$ is a coefficient of a negative power of $z$, which kills $\lthreevac$), its $z^0$-coefficient is $\frac1{k!}(T^ku)_{(-1)}v = u_{(-k-1)}v$ (iterate
\lean{mode_Y_T}: $(Tu)_{(n)} = -n\,u_{(n-1)}$), and the higher coefficients are pinned by translation
covariance of the normal-ordered product (\lean{nprod_L_comm}) and $L_{-1}\lthreevac = 0$ via the
recursion $L_{-1}(F_s\lthreevac) = (s+1)F_{s+1}\lthreevac$ (\lean{vacuum_recursion_eq}). So it creates
$Y(u_{(-k-1)}v,z)\lthreevac$, and Goddard finishes. \qed

\begin{status}
Unconditional, axiom-clean (\texttt{\#print axioms} \lean{VOA.VertexOperatorAlgebra.assoc_neg} run for this
chapter).
\end{status}

\subsection{The iterate formula for all \texorpdfstring{$\ell$}{l}}

\begin{theorem}[Iterate (associativity) formula]\label{l03:thm:iterate}
In a VOA, for all $u,v,x$ and $\ell,n\in\Z$,
\begin{equation*}
(u_{(\ell)}v)_{(n)}x = \sum_{i\ge0}(-1)^i\binom\ell i\Big(u_{(\ell-i)}v_{(n+i)}x - (-1)^\ell\,v_{(\ell+n-i)}u_{(i)}x\Big).
\end{equation*}
\end{theorem}
\begin{leancode}
theorem mode_iterate_borcherds_nonneg (W : VertexOperatorAlgebra 𝒱) (u v : 𝒱) (k : ℕ) (n : ℤ) :
    mode (W.Y (W.smul u (k : ℤ) v)) n =
      ∑ i ∈ Finset.range (k + 1), ((-1 : ℂ) ^ i * gBinom (k : ℤ) i) •
        (mode (W.Y u) ((k : ℤ) - i) * mode (W.Y v) (n + i)
          - ((-1 : ℂ) ^ k) • (mode (W.Y v) ((k : ℤ) + n - i) * mode (W.Y u) (i)))
\end{leancode}
File \leanpath{VOA/\allowbreak Core/\allowbreak Reconstruction.lean}.
\begin{leancode}
theorem mode_iterate_neg (W : VertexOperatorAlgebra 𝒱) (u v : 𝒱) (k : ℕ) (N : ℤ) (x : 𝒱) :
    W.smul (W.smul u (-(k:ℤ) - 1) v) N x
      = ∑ᶠ i : ℕ, ((-1:ℂ)^i * gBinom (-(k:ℤ)-1) i) •
          (W.smul u (-(k:ℤ)-1-i) (W.smul v (N+i) x)
            - ((-1:ℂ)^(-(k:ℤ)-1)) • W.smul v (-(k:ℤ)-1+N-i) (W.smul u i x))

theorem mode_iterate (W : VertexOperatorAlgebra 𝒱) (u v : 𝒱) (ℓ N : ℤ) (x : 𝒱) :
    W.smul (W.smul u ℓ v) N x
      = ∑ᶠ i : ℕ, ((-1:ℂ)^i * gBinom ℓ i) •
          (W.smul u (ℓ - i) (W.smul v (N + i) x)
            - ((-1:ℂ)^ℓ) • W.smul v (ℓ + N - i) (W.smul u i x))
\end{leancode}
File \leanpath{VOA/\allowbreak Core/\allowbreak BorcherdsIdentity.lean}.

\begin{dictionary}
\textbf{Operator vs.\ vector form.} For $\ell = k\ge0$ the identity is stated as an equality of
operators (no $x$), a sum over $i\in\{0,\dots,k\}$ since $\binom ki = 0$ for $i>k$. For $\ell<0$ the
binomial $\binom\ell i$ never vanishes and the sum over $i$ is infinite as an operator identity; it is
finite only after applying to $x$ (the factors $v_{(N+i)}x$ and $u_{(i)}x$ vanish for $i\gg0$ by VA5).
Hence \lean{mode_iterate} is stated on vectors with \lean{∑ᶠ} (\lean{finsum}: a genuine sum when the
support is finite, $0$ otherwise). Its support is finite for the reason just given, so \lean{∑ᶠ} is a
true finite sum here (the library proves the corresponding finiteness lemmas \lean{Puv_fin},
\lean{Pvu_fin}, Section~\ref{l03:sec:borcherds}).

\textbf{Signs.} \lean{(-1:ℂ)^ℓ} with \lean{ℓ : ℤ} is the integer power ($\pm1$); for $\ell = -k-1$ it is
$(-1)^{k+1}$ (library lemma \lean{neg_one_zpow_negsucc}). The formula is Gannon (5.1.7d), the iterate formula of Lepowsky--Li (§3.1), and the
$m = 0$ case of the Borcherds identity in Kac (§4.8).
\end{dictionary}

\begin{physics}
The iterate formula expresses the modes of a composite field $(u_{(\ell)}v)(z)$ through the modes of
its constituents, with radial ordering taken into account (the two sums are the two orderings
$|z_1|>|z_2|$ and $|z_2|>|z_1|$). It is the mode form of associativity of the OPE: expanding
$u(z_1)v(z_2)x$ around $z_2$ first, or computing the OPE $u\times v$ first and then acting on $x$, gives
the same answer.
\end{physics}

\textbf{Proof on paper.} For $\ell = k\ge0$: by Theorem~\ref{l03:thm:nthproduct},
$(u_{(k)}v)_{(n)}$ is the $z^{-n-1}$-coefficient of $\Res_{z_1}(z_1-z)^k[Y(u,z_1),Y(v,z)]$; expanding
$(z_1-z)^k = \sum_r\binom kr z_1^r(-z)^{k-r}$ gives $\sum_r(-1)^{k-r}\binom kr[u_{(r)},v_{(n+k-r)}]$, and
reindexing $r = k-i$ in the first half (with $\binom k{k-i} = \binom ki$) and writing
$(-1)^{k-i} = (-1)^i(-1)^k$ in the second half gives the displayed form. For $\ell = -k-1$: by Theorem~\ref{l03:thm:assocneg} the left side
is the $z^{-N-1}$-coefficient of $\frac1{k!}{:}\partial^kY(u)Y(v){:}$ applied to $x$, i.e.\
$\frac1{k!}\sum_m{:}(\partial^kY(u))_mY(v)_{-N-1-m}{:}x$. Split the sum into $m\ge0$ (creation part,
on the left) and $m\le-k-1$ (annihilation part, on the right; the range $-k-1<m<0$ contributes zero
because the falling factorial from $\partial^k$ vanishes there). Reindex the two ranges by $m = i$ and
$m = -k-1-i$; the coefficients are $\frac1{k!}\prod_{j<k}(i+1+j) = \binom{k+i}{i} = (-1)^i\binom{-k-1}{i}$
and its signed analogue (\lean{posfac_eq}, \lean{negfac_eq}). This is the displayed formula. \qed

\begin{status}
\lean{mode_iterate} (and both branches): unconditional; \texttt{\#print axioms} \lean{VOA.VertexOperatorAlgebra.mode_iterate} $=$ \lean{[propext, Classical.choice, Quot.sound]} (run for
this chapter).
\end{status}

%======================================================================
\section{The Borcherds identity}\label{l03:sec:borcherds}
%======================================================================

\begin{definition}[The two sides]\label{l03:def:BLBR}
For $u,v,x\in V$ and $\ell,m,n\in\Z$ put
\begin{align*}
\lthreeBR(\ell,m,n) &= \sum_{i\ge0}\binom mi(u_{(\ell+i)}v)_{(m+n-i)}x,\\
P_{uv}(\ell,m,n) &= \sum_{i\ge0}(-1)^i\binom\ell i\,u_{(\ell+m-i)}v_{(n+i)}x,\qquad
P_{vu}(\ell,m,n) = \sum_{i\ge0}(-1)^i\binom\ell i\,v_{(\ell+n-i)}u_{(m+i)}x,\\
\lthreeBL(\ell,m,n) &= P_{uv}(\ell,m,n) - (-1)^\ell P_{vu}(\ell,m,n).
\end{align*}
\end{definition}
\begin{leancode}
noncomputable def BR (W : VertexOperatorAlgebra 𝒱) (u v x : 𝒱) (ℓ m n : ℤ) : 𝒱 :=
  ∑ᶠ i : ℤ, cbz m i • W.smul (W.smul u (ℓ + i) v) (m + n - i) x

noncomputable def Puv (W : VertexOperatorAlgebra 𝒱) (u v x : 𝒱) (ℓ m n : ℤ) : 𝒱 :=
  ∑ᶠ i : ℤ, (negOnePow i * cbz ℓ i) • W.smul u (ℓ + m - i) (W.smul v (n + i) x)

noncomputable def Pvu (W : VertexOperatorAlgebra 𝒱) (u v x : 𝒱) (ℓ m n : ℤ) : 𝒱 :=
  ∑ᶠ i : ℤ, (negOnePow i * cbz ℓ i) • W.smul v (ℓ + n - i) (W.smul u (m + i) x)

noncomputable def BL (W : VertexOperatorAlgebra 𝒱) (u v x : 𝒱) (ℓ m n : ℤ) : 𝒱 :=
  Puv W u v x ℓ m n - negOnePow ℓ • Pvu W u v x ℓ m n
\end{leancode}
File \leanpath{VOA/\allowbreak Core/\allowbreak Borcherds.lean}.

\begin{dictionary}
\lean{cbz n i} is $\binom ni$ for $i\ge0$ and $0$ for $i<0$ (\leanpath{VOA/\allowbreak Foundations/\allowbreak
Substitution.lean}), so the sums over $i\in\Z$ are sums over $i\ge0$. \lean{negOnePow i} is $(-1)^i$
as a complex number. Each of the three \lean{∑ᶠ} has finite support (lemmas \lean{BR_fin},
\lean{Puv_fin}, \lean{Pvu_fin}): for $\lthreeBR$ by VA5 for $(u,v)$; for $P_{uv}$ by VA5 for $(v,x)$; for
$P_{vu}$ by VA5 for $(u,x)$. So none of them is a silently zero ``divergent'' \lean{finsum}.
\end{dictionary}

\begin{physics}
$\lthreeBL$ and $\lthreeBR$ are the two sides of OPE associativity in a three-point function
$\langle x'|u(z_1)v(z_2)|x\rangle$: on the left, the two radial orderings of $u(z_1)v(z_2)$ (the
difference is a contour around $z_2$); on the right, the OPE $u(z_1)v(z_2)$ expanded around $z_2$ and
then $Y(\cdot,z_2)$ acting on $x$. The weights $(z_1-z_2)^\ell$ and $z_1^m$ select one moment of each.
\end{physics}

\begin{lemma}[Pascal recursion and the anti-diagonal pin]\label{l03:lem:pascal}
Both sides satisfy $B(\ell,m+1,n) = B(\ell,m,n+1) + B(\ell+1,m,n)$. Moreover, if a family
$D:\Z^2\to M$ in an abelian group satisfies $D(\ell,n+1) = -D(\ell+1,n)$ for all $\ell,n$ and $D(0,n) = 0$
for all $n$, then $D\equiv0$.
\end{lemma}
\begin{leancode}
theorem BR_rec (W : VertexOperatorAlgebra 𝒱) (u v x : 𝒱) (ℓ m n : ℤ) :
    BR W u v x ℓ (m + 1) n = BR W u v x ℓ m (n + 1) + BR W u v x (ℓ + 1) m n

theorem BL_rec (W : VertexOperatorAlgebra 𝒱) (u v x : 𝒱) (ℓ m n : ℤ) :
    BL W u v x ℓ (m + 1) n = BL W u v x ℓ m (n + 1) + BL W u v x (ℓ + 1) m n

theorem antidiag_pin {M : Type*} [AddCommGroup M] (D : ℤ → ℤ → M)
    (hrec : ∀ ℓ n : ℤ, D ℓ (n + 1) = -D (ℓ + 1) n) (hbd : ∀ n : ℤ, D 0 n = 0) :
    ∀ ℓ n : ℤ, D ℓ n = 0
\end{leancode}
File \leanpath{VOA/\allowbreak Core/\allowbreak Borcherds.lean}.

\begin{dictionary}
Pure combinatorics; \lean{BR_rec} and \lean{BL_rec} hold in every VOA without using any identity
between modes (only Pascal's rule for $\binom\cdot\cdot$ and reindexing). \lean{antidiag_pin} is a
statement about arbitrary families.
\end{dictionary}

\begin{physics}
The recursion is the identity $z_1^{m+1} = z_1^m\big((z_1-z_2)+z_2\big)$ inside the Jacobi identity:
raising the power of $z_1$ is the same as raising the power of $z_2$ or of $z_1-z_2$. Physically, the
three moments of a correlator are not independent.
\end{physics}

\textbf{Proof on paper.} For $\lthreeBR$: $\binom{m+1}i = \binom mi + \binom m{i-1}$, and the shift
$i\mapsto i+1$ in the second sum turns it into $\lthreeBR(\ell+1,m,n)$. For $P_{uv}$ and $P_{vu}$ use
$\binom{\ell+1}{i} = \binom\ell i+\binom\ell{i-1}$ in the form $\binom\ell{i-1} = \binom{\ell+1}i - \binom\ell i$:
one gets $P_{uv}(\ell,m+1,n) = P_{uv}(\ell,m,n+1) + P_{uv}(\ell+1,m,n)$ and
$P_{vu}(\ell,m+1,n) = P_{vu}(\ell,m,n+1) - P_{vu}(\ell+1,m,n)$, and the sign $(-1)^{\ell+1} = -(-1)^\ell$
makes $\lthreeBL$ satisfy the same recursion as $\lthreeBR$. For the pin: $D(\ell,n) = -D(\ell+1,n-1) =
D(\ell+2,n-2) = \cdots$, so $D(\ell,n) = \pm D(0,n+\ell) = 0$ (and similarly moving left for $\ell<0$). \qed

\begin{status}
Unconditional, axiom-clean (in the cone of \lean{borcherds}).
\end{status}

\begin{theorem}[The Borcherds identity]\label{l03:thm:borcherds}
In a VOA, for all $u,v,x\in V$ and all $\ell,m,n\in\Z$, $\lthreeBL(\ell,m,n) = \lthreeBR(\ell,m,n)$, i.e.
\begin{equation}\label{l03:eq:borcherds}
\sum_{i\ge0}(-1)^i\binom\ell i\Big(u_{(\ell+m-i)}v_{(n+i)}x - (-1)^\ell v_{(\ell+n-i)}u_{(m+i)}x\Big)
 = \sum_{i\ge0}\binom mi(u_{(\ell+i)}v)_{(m+n-i)}x .
\end{equation}
\end{theorem}
\begin{leancode}
theorem borcherds_base (W : VertexOperatorAlgebra 𝒱) (u v x : 𝒱) (ℓ n : ℤ) :
    BL W u v x ℓ 0 n = BR W u v x ℓ 0 n

theorem borcherds_boundary (W : VertexOperatorAlgebra 𝒱) (u v x : 𝒱) (m n : ℤ) :
    BL W u v x 0 m n = BR W u v x 0 m n

theorem borcherds (W : VertexOperatorAlgebra 𝒱) (u v x : 𝒱) (ℓ m n : ℤ) :
    BL W u v x ℓ m n = BR W u v x ℓ m n

theorem jacobi (W : VertexOperatorAlgebra 𝒱) (u v x : 𝒱) (ℓ m n : ℤ) :
    (∑ᶠ i : ℤ, (negOnePow i * cbz ℓ i) •
        (W.smul u (ℓ + m - i) (W.smul v (n + i) x)
          - negOnePow ℓ • W.smul v (ℓ + n - i) (W.smul u (m + i) x)))
      = ∑ᶠ i : ℤ, cbz m i • W.smul (W.smul u (ℓ + i) v) (m + n - i) x
\end{leancode}
File \leanpath{VOA/\allowbreak Core/\allowbreak Borcherds.lean}.

\begin{dictionary}
\textbf{All indices.} \lean{borcherds} and \lean{jacobi} are quantified over all $u,v,x$ and all
$\ell,m,n\in\Z$, with no side condition: this is the Borcherds identity in full, Gannon (5.1.7b),
Kac (4.8.3), equivalent coefficient by coefficient to the Jacobi identity \eqref{l03:eq:jacobi}
(the extraction is sketched after \eqref{l03:eq:jacobi}). \lean{jacobi} is \lean{borcherds} with
$\lthreeBL$ written as one \lean{finsum} (lemma \lean{BL_as_finsum}). The base case $m = 0$ is
Theorem~\ref{l03:thm:iterate}; the boundary $\ell = 0$ is Theorem~\ref{l03:thm:comm} applied to $x$.

\textbf{What it does not say.} It is stated for \lean{VertexOperatorAlgebra} only. For a
\lean{PreVertexAlgebra} it is false in general (Theorem~\ref{l03:thm:counter}).
\end{dictionary}

\begin{physics}
The Borcherds identity is OPE associativity together with locality --- crossing symmetry of the
chiral three-point functions, in mode form. Every relation among OPE coefficients of a chiral algebra
is a consequence of it; its $\ell = 0$ slice gives the mode algebra (current, Virasoro, $W$-algebras),
its $m = 0$ slice the modes of composite fields. Proving it from locality is the rigorous version of
``single-valued correlators with finite poles are automatically associative''.
\end{physics}

\textbf{Proof on paper.} Let $D_m(\ell,n) = \lthreeBL(\ell,m,n) - \lthreeBR(\ell,m,n)$. By
Lemma~\ref{l03:lem:pascal}, $D_{m+1}(\ell,n) = D_m(\ell,n+1) + D_m(\ell+1,n)$. For $m\ge0$: $D_0\equiv0$
(base case), so $D_m\equiv0$ by induction upward. For $m = -k-1<0$, induct downward: if $D_{m+1}\equiv0$
then $D_m(\ell,n+1) = -D_m(\ell+1,n)$, and $D_m(0,n) = 0$ by the boundary case; the anti-diagonal pin
gives $D_m\equiv0$. No delta-function calculus and no rearrangement of infinite sums is used. \qed

\begin{status}
\lean{borcherds}, \lean{jacobi}: unconditional; \texttt{\#print axioms} $=$
\lean{[propext, Classical.choice, Quot.sound]} for both (run for this chapter). The PBW theorem \lean{VOA.pbw} is not used.
\end{status}

\begin{theorem}[The commutator formula, as re-exposed in the Borcherds layer]\label{l03:thm:jacobicomm}
In a VOA, for all $u,v\in V$ and $a,b\in\Z$ there is a natural number $N$ such that
\begin{equation*}
 [u_{(a)},v_{(b)}] \;=\; \sum_{j=0}^{N-1}\binom aj\,(u_{(j)}v)_{(a+b-j)} .
\end{equation*}
It is the case $\ell = 0$ of the Borcherds identity (Theorem~\ref{l03:thm:borcherds}); the case $m=0$ is the iterate
formula (Theorem~\ref{l03:thm:iterate}).
\end{theorem}
\noindent File \leanpath{VOA/\allowbreak Core/\allowbreak BorcherdsIdentity.lean}.
\begin{leancode}
theorem jacobi_commutator (W : VertexOperatorAlgebra 𝒱) (u v : 𝒱) (a b : ℤ) :
    ∃ N : ℕ, ⁅mode (W.Y u) a, mode (W.Y v) b⁆
      = ∑ j : Fin N, gBinom a (j : ℕ) • mode (W.Y (W.smul u (j : ℕ) v)) (a + b - (j : ℕ))
\end{leancode}

\begin{dictionary}
\lean{⁅x, y⁆} is the commutator $xy-yx$ in the ring $\End(V)$ (here of two mode operators);
\lean{mode (W.Y u) a} is $u_{(a)}$; \lean{W.smul u j v} is $u_{(j)}v$; \lean{gBinom a j} is the binomial coefficient
$\binom aj = a(a-1)\cdots(a-j+1)/j!$ for $a\in\Z$ (Definition~\ref{l01:def:gbinom}); \lean{Fin N} indexes $j = 0,\dots,N-1$;
\lean{(j : ℕ)} is the coercion $\mathrm{Fin}\,N\to\N$. The statement is word for word that of \lean{mode_commutator}
(Theorem~\ref{l03:thm:comm}); in the source the proof is the single term \lean{mode_commutator W u v a b}, and the
declaration exists so that the Borcherds file can name the $\ell = 0$ slice without importing the reconstruction file.
Hence the same fidelity remarks apply: because of the existential quantifier, read alone the statement does not say that
the omitted terms $j\ge N$ vanish, so it is formally weaker than the textbook ``$\sum_{j\ge0}$''. The strengthening is
\lean{commBound} of Theorem~\ref{l03:thm:comm}, which supplies the same $N$ with $u_{(k)}v = 0$ for $k\ge N$. The
\emph{iterate} formula, the case $m=0$, is a different statement, \lean{mode_iterate} of
Theorem~\ref{l03:thm:iterate}, and is not part of this theorem. \textbf{Could it be vacuous?} No: for the
Virasoro field $u=v=\omega$ one finds $N = 4$ and the Virasoro relations.
\end{dictionary}

\begin{physics}
Commutators of modes are the singular part of the OPE: $[u_{(a)},v_{(b)}]$ is the contour integral of the
singular terms $\sum_j(u_{(j)}v)(w)/(z-w)^{j+1}$, which is why only finitely many $j$ occur (a pole of bounded order).
For currents this is the affine Kac--Moody algebra; for $\omega$ it is the Virasoro algebra with its central term
$\tfrac c{12}(m^3-m)\delta_{m+n,0}$.
\end{physics}

\textbf{Proof on paper.} Put $\ell=0$ in the Borcherds identity~\eqref{l03:eq:borcherds}. Since
$\binom0i = \delta_{i,0}$ the left side collapses to $u_{(m)}v_{(n)}x - v_{(n)}u_{(m)}x$, and the right side is
$\sum_i\binom mi(u_{(i)}v)_{(m+n-i)}x$ with $m=a$, $n=b$. The sum is finite because $u_{(i)}v = 0$ for $i\gg0$ (regularity),
and the bound $N$ may be chosen independent of $x$ since the right side is a finite sum of operators. \qed

\begin{status}
Unconditional. \texttt{\#print axioms} \lean{VOA.VertexOperatorAlgebra.jacobi\_commutator} $=$
\lean{[propext, Classical.}\allowbreak\lean{choice, Quot.}\allowbreak\lean{sound]} (run for this chapter). The PBW theorem \lean{VOA.pbw} is not used.
\end{status}

%======================================================================
\section{Skew-symmetry}\label{l03:sec:skew}
%======================================================================

Skew-symmetry is the hardest of the classical identities in the library's route, because the
natural proof (Goddard uniqueness for a reflected field) requires a locality statement for a field
that is not a vertex operator. The library proves it in three files: \leanpath{Skew.lean} defines
the candidate and reduces the theorem to a locality gate, \leanpath{WeakAssoc.lean} proves the gate,
and \leanpath{SkewSymmetry.lean} assembles.

\subsection{The candidate field}

\begin{definition}[The reflected field]\label{l03:def:skewdist}
For $u\in V$ define $\tilde Y(u,z)\in\End(V)[[z^{\pm1}]]$ by $\tilde Y(u,z)b := e^{zT}Y(b,-z)u$, i.e.\
its $z^t$-coefficient is
\begin{equation*}
\tilde Y(u)_t\,b = \sum_{j\ge0}\frac{(-1)^{t-j}}{j!}\,T^j\big(Y(b)_{t-j}\,u\big),
\end{equation*}
where $Y(b)_s$ is the $z^s$-coefficient of $Y(b,z)$. The sum is finite for each $b$.
\end{definition}
\begin{leancode}
def sgn (n : ℤ) : ℂ := if Even n then 1 else -1

noncomputable def skewCoeffFun (W : VertexOperatorAlgebra 𝒱) (u : 𝒱) (t : ℤ) (b : 𝒱) : 𝒱 :=
  ∑ᶠ j : ℕ, ((j.factorial : ℂ)⁻¹) • (sgn (t - j)) •
    ((W.toPreVertexAlgebra.T ^ j) ((W.Y b) (t - j) u))

theorem skew_support_finite (W : VertexOperatorAlgebra 𝒱) (u b : 𝒱) (t : ℤ) :
    (Function.support (fun j : ℕ => ((j.factorial : ℂ)⁻¹) • (sgn (t - j)) •
      ((W.toPreVertexAlgebra.T ^ j) ((W.Y b) (t - j) u)))).Finite

noncomputable def skewDist (W : VertexOperatorAlgebra 𝒱) (u : 𝒱) : FormalDist (End ℂ 𝒱) :=
  fun t => skewOp W u t
\end{leancode}
File \leanpath{VOA/\allowbreak Core/\allowbreak Skew.lean} (\lean{skewOp W u t} is
\lean{skewCoeffFun W u t} packaged as a linear map in $b$).

\begin{dictionary}
\lean{sgn n} $=(-1)^n$. \lean{(W.Y b) (t - j) u} is the $z^{t-j}$-coefficient of $Y(b,z)u$; replacing
$z$ by $-z$ multiplies it by $(-1)^{t-j}$, and $e^{zT} = \sum_jz^jT^j/j!$ shifts the power by $j$. So
\lean{skewCoeffFun W u t b} is the $z^t$-coefficient of $e^{zT}Y(b,-z)u$, as displayed. Finiteness
(\lean{skew_support_finite}): $Y(b)_{t-j}u = b_{(j-t-1)}u = 0$ once $j-t-1\ge N(b,u)$ (VA5), so the
\lean{∑ᶠ} is a genuine finite sum.
\end{dictionary}

\begin{physics}
$\tilde Y(u,z)b = e^{zT}Y(b,-z)u$ is the operator obtained by \emph{exchanging} the two insertions in
$u(z)b(0)$: put $b$ at $0$ and $u$ at $z$, swap their roles, and translate back by $z$. For bosonic
chiral fields, OPE exchange symmetry says this is the same operator as $Y(u,z)$.
\end{physics}

\begin{lemma}\label{l03:lem:skewvac}
$\tilde Y(u,z)\lthreevac = Y(u,z)\lthreevac$.
\end{lemma}
\begin{leancode}
theorem skewCoeffFun_vacuum (W : VertexOperatorAlgebra 𝒱) (u : 𝒱) (t : ℤ) :
    skewCoeffFun W u t W.vacuum = (W.Y u) t W.vacuum
\end{leancode}
File \leanpath{VOA/\allowbreak Core/\allowbreak Skew.lean}.

\begin{dictionary}
For $b = \lthreevac$, $Y(\lthreevac,-z)u = u$ (VA2), so only $j = t$ survives for $t\ge0$ and the value is
$T^tu/t! = u_{(-t-1)}\lthreevac$ by \lean{creation_pow}; for $t<0$ both sides vanish. Uses (VOA3)
through \lean{creation_pow}.
\end{dictionary}

\begin{physics}
Both operators create the same state from the vacuum: $e^{zL_{-1}}u$.
\end{physics}

\subsection{The reduction to a locality statement}

\begin{proposition}[Skew-symmetry reduces to locality of $\tilde Y$]\label{l03:prop:skewred}
If, for every $c$, some power $(z_1-z_2)^N$ kills $[\tilde Y(u,z_1),Y(c,z_2)]$ applied to every
vector (uniformly), then $\tilde Y(u,z) = Y(u,z)$. It suffices that for every $c$ some
$(z_1-z_2)^N$ kills $[\tilde Y(u,z_1),Y(c,z_2)]\lthreevac$.
\end{proposition}
\begin{leancode}
theorem skew_eq_Y_of_local (W : VertexOperatorAlgebra 𝒱) (u : 𝒱)
    (hloc : ∀ c : 𝒱, Local (skewDist W u) (W.Y c)) :
    skewDist W u = W.Y u

def WeakAssocLoc (W : VertexOperatorAlgebra 𝒱) (u : 𝒱) : Prop :=
  ∀ c : 𝒱, ∃ N : ℕ, ∀ b : 𝒱,
    mulBySubPow N (applyVec b (bracketDist (skewDist W u) (W.Y c))) = 0

theorem skewDist_eq_Y_of_weakAssoc (W : VertexOperatorAlgebra 𝒱) (u : 𝒱)
    (hWA : WeakAssocLoc W u) :
    skewDist W u = W.Y u
\end{leancode}
File \leanpath{VOA/\allowbreak Core/\allowbreak Skew.lean}.
\begin{leancode}
def SkewVacLocal (W : VertexOperatorAlgebra 𝒱) (u : 𝒱) : Prop :=
  ∀ c : 𝒱, ∃ N : ℕ,
    mulBySubPow N (applyVec W.vacuum (bracketDist (skewDist W u) (W.Y c))) = 0

theorem skewDist_eq_Y_of_vacLocal (W : VertexOperatorAlgebra 𝒱) (u : 𝒱)
    (h : SkewVacLocal W u) : skewDist W u = W.Y u

theorem weakAssocLoc_of_vacLocal (W : VertexOperatorAlgebra 𝒱) (u : 𝒱)
    (h : SkewVacLocal W u) : WeakAssocLoc W u
\end{leancode}
File \leanpath{VOA/\allowbreak Core/\allowbreak WeakAssoc.lean}.

\begin{dictionary}
\lean{WeakAssocLoc W u}: for each $c$ a single $N$ works for all test vectors $b$ --- this
uniformity is what upgrades the vector statements to the operator statement \lean{Local}
(\lean{local_of_applyVec_annihilation}). \lean{SkewVacLocal W u}: the same with only the vacuum as
test vector, which is much easier to verify; Goddard's vacuum form
(Theorem~\ref{l03:thm:goddard}(ii)) turns it into $\tilde Y = Y$, and \emph{then}
\lean{WeakAssocLoc} holds a posteriori by VA4. The file \leanpath{Skew.lean} also contains a
predicate \lean{FLMCoeff} described in its docstring as ``the single remaining open lemma''; that
description is out of date --- the gate is now closed by \lean{skewVacLocal} below, and
\lean{FLMCoeff} is not used by \lean{skew}.
\end{dictionary}

\begin{physics}
To prove exchange symmetry it is enough to show that the exchanged operator is a legitimate local
field; uniqueness of the field of a state then does the rest. The locality check, moreover, only has
to be done inside vacuum correlators.
\end{physics}

\begin{status}
These reductions are unconditional and axiom-clean; their hypotheses are discharged next.
\end{status}

\subsection{The vacuum locality computation}

\begin{lemma}[$T$ and modes]\label{l03:lem:smulTpow}
In a VOA, $c_{(p)}(Tw) = T(c_{(p)}w) + p\,c_{(p-1)}w$, and more generally
\begin{equation*}
c_{(p)}(T^ku) = \sum_{r=0}^k\binom kr\,p^{\underline r}\;T^{k-r}\big(c_{(p-r)}u\big),\qquad
p^{\underline r} = p(p-1)\cdots(p-r+1).
\end{equation*}
\end{lemma}
\begin{leancode}
theorem smul_comm_T (W : VertexOperatorAlgebra 𝒱) (c w : 𝒱) (p : ℤ) :
    W.smul c p (W.toPreVertexAlgebra.T w)
      = W.toPreVertexAlgebra.T (W.smul c p w) + p • W.smul c (p - 1) w

def fallI (p : ℤ) (r : ℕ) : ℤ := ∏ i ∈ Finset.range r, (p - (i : ℤ))

def mcoeff (p : ℤ) (k r : ℕ) : ℤ := (k.choose r : ℤ) * fallI p r

theorem smul_Tpow_expand (W : VertexOperatorAlgebra 𝒱) (c u : 𝒱) (p : ℤ) (k : ℕ) :
    W.smul c p ((W.toPreVertexAlgebra.T ^ k) u)
      = ∑ r ∈ Finset.range (k + 1),
          (mcoeff p k r) • ((W.toPreVertexAlgebra.T ^ (k - r)) (W.smul c (p - (r : ℤ)) u))
\end{leancode}
File \leanpath{VOA/\allowbreak Core/\allowbreak WeakAssoc.lean}.

\begin{dictionary}
\lean{fallI p r} is the falling factorial $p^{\underline r}$, \lean{mcoeff p k r} $=\binom kr p^{\underline r}$, an
integer acting by \lean{•}. \lean{smul_comm_T} is translation covariance in the form
$[c_{(p)},T] = -[T,c_{(p)}] = p\,c_{(p-1)}$ (from \lean{L_neg_one_deriv} and \lean{T_eq_Lneg1});
\lean{smul_Tpow_expand} is its $k$-fold iterate, proved by induction with Pascal's rule.
\end{dictionary}

\begin{physics}
$e^{-zT}Y(c,w)e^{zT} = Y(c,w-z)$: conjugating by a translation shifts the insertion point. The lemma
is this identity coefficient by coefficient.
\end{physics}

\begin{theorem}[Skew-locality on the vacuum, and weak associativity]\label{l03:thm:skewvac}
In a VOA, for all $u$: \lean{SkewVacLocal W u} and \lean{WeakAssocLoc W u} hold.
\end{theorem}
\begin{leancode}
theorem skewVacLocal (W : VertexOperatorAlgebra 𝒱) (u : 𝒱) : SkewVacLocal W u

theorem weakAssocLoc (W : VertexOperatorAlgebra 𝒱) (u : 𝒱) : WeakAssocLoc W u :=
  weakAssocLoc_of_vacLocal W u (skewVacLocal W u)
\end{leancode}
File \leanpath{VOA/\allowbreak Core/\allowbreak WeakAssoc.lean}.

\begin{dictionary}
No hypotheses beyond the VOA. The order $N$ for \lean{skewVacLocal} is the regularity bound $M_0$ of
the pair $(c,u)$ (\lean{W.toPreVertexAlgebra.regularity c u}), i.e.\ the pole order of $Y(c,w)u$.
\end{dictionary}

\begin{physics}
This is the computation that the exchanged operator has the same singularity structure as an honest
field: inside vacuum correlators, $[\tilde Y(u,z_1),Y(c,z_2)]|0\rangle$ is a contact term supported on
$z_1 = z_2$ of order at most the pole order of the OPE $c\times u$.
\end{physics}

\textbf{Proof on paper.} Let $G(z_1,z_2) := [\tilde Y(u,z_1),Y(c,z_2)]\lthreevac$. Using
Lemma~\ref{l03:lem:skewvac} and creation,
\begin{equation*}
G = \tilde Y(u,z_1)\,e^{z_2T}c - Y(c,z_2)\,e^{z_1T}u .
\end{equation*}
Second term: by Lemma~\ref{l03:lem:smulTpow}, $Y(c,z_2)e^{z_1T}u = e^{z_1T}\,Y(c,z_2-z_1)u$ with
$(z_2-z_1)^{-s-1}$ expanded in nonnegative powers of $z_1$ (closed form \lean{Term2_expand}). First
term: by definition and $Y(T^jc,w) = \partial^jY(c,w)$,
$\tilde Y(u,z_1)e^{z_2T}c = e^{z_1T}\sum_j\frac{z_2^j}{j!}\partial^j_{w}Y(c,w)u\big|_{w=-z_1} =
e^{z_1T}\,Y(c,z_2-z_1)u$ with $(z_2-z_1)^{-s-1}$ expanded in nonnegative powers of $z_2$ (closed form
\lean{Term1_expand}). Since $Y(c,w)u = \sum_{s<M_0}(c_{(s)}u)w^{-s-1} + (\text{regular})$, the
regular parts agree, and $G$ is $e^{z_1T}$ applied to $\sum_{s<M_0}(c_{(s)}u)\big(\iota_{z_2,z_1} -
\iota_{z_1,z_2}\big)(z_2-z_1)^{-s-1}$, a finite sum of derivatives of $\delta(z_1/z_2)$ of order $<M_0$
(the ``jump sum'' of \lean{Gfield_eq_jumpSum}). Each is killed by $(z_1-z_2)^{M_0}$
(\lean{mulBySubPow_jumpSum_eq_zero}). The library carries this out entirely by coefficient
comparison in the four sign sectors of $(z_1,z_2)$-exponents, without formal delta-function
substitution. \qed

\begin{status}
Unconditional, axiom-clean (\texttt{\#print axioms} \lean{VOA.VertexOperatorAlgebra.weakAssocLoc} run for this
chapter). The file's history (a gate once sanctioned as an axiom \lean{weakAssocLoc_gate}, later a single
\lean{sorry} in \lean{Gfield_eq_jumpSum}) is closed: neither occurs at this commit.
\end{status}

\subsection{Skew-symmetry}

\begin{theorem}[Skew-symmetry]\label{l03:thm:skew}
In a VOA, $Y(u,z)v = e^{zT}Y(v,-z)u$ for all $u,v$. In modes,
$u_{(n)}v = \sum_{j\ge0}\frac{(-1)^{n+j+1}}{j!}T^j(v_{(n+j)}u)$.
\end{theorem}
\begin{leancode}
theorem skew (W : VertexOperatorAlgebra 𝒱) (u : 𝒱) : skewDist W u = W.Y u :=
  skewDist_eq_Y_of_weakAssoc W u (weakAssocLoc W u)

theorem skew_mode (W : VertexOperatorAlgebra 𝒱) (u v : 𝒱) (n : ℤ) :
    W.smul u n v
      = ∑ᶠ j : ℕ, ((j.factorial : ℂ)⁻¹) • (sgn (-n - 1 - j)) •
          ((W.toPreVertexAlgebra.T ^ j) ((W.Y v) (-n - 1 - j) u))
\end{leancode}
File \leanpath{VOA/\allowbreak Core/\allowbreak SkewSymmetry.lean}.

\begin{dictionary}
\lean{skew} is an equality of operator-valued distributions, $\tilde Y(u,\cdot) = Y(u,\cdot)$, which by
Definition~\ref{l03:def:skewdist} is $Y(u,z)v = e^{zT}Y(v,-z)u$ for all $v$. \lean{skew_mode} reads
off the $z^{-n-1}$-coefficient ($t = -n-1$): \lean{(W.Y v) (-n - 1 - j)} $=v_{(n+j)}$ and
\lean{sgn (-n - 1 - j)} $=(-1)^{n+j+1}$. Same as Gannon (5.1.8), Kac (4.2.1), Lepowsky--Li (§3.1). The
one-line proof term is shown because it is informative: skew-symmetry \emph{is} the reduction applied
to the proved gate.
\end{dictionary}

\begin{physics}
Bose exchange symmetry of the chiral algebra: $u(z)v(0) = e^{zL_{-1}}v(-z)u(0)$ inside correlators.
In terms of OPE coefficients, the poles of $v(z)u(0)$ determine those of $u(z)v(0)$; e.g.\ for
currents $J^a_{(0)}J^b = -J^b_{(0)}J^a$ (antisymmetry of the Lie bracket on $V_1$) and
$J^a_{(1)}J^b = J^b_{(1)}J^a$ (symmetry of the level form).
\end{physics}

\textbf{Proof on paper.} Theorem~\ref{l03:thm:skewvac} and Proposition~\ref{l03:prop:skewred}. \qed

\begin{status}
Unconditional; \texttt{\#print axioms} $=$ \lean{[propext, Classical.choice, Quot.sound]} for
\lean{skew} and \lean{skew_mode} (run for this chapter).
\end{status}

%======================================================================
\section[The FLM formulation]{The FLM formulation: associativity after a change of variables}\label{l03:sec:flm}
%======================================================================

FLM and Lepowsky--Li formulate the heart of the Jacobi identity as \emph{weak associativity}:
$(z_0+z_2)^NY(u,z_0+z_2)Y(v,z_2)x = (z_0+z_2)^NY(Y(u,z_0)v,z_2)x$. The library proves a variant in
which the polynomial factor is $(z_1-z_2)^N$ with $N$ the locality order, and the substitution
$z_1 = z_0+z_2$ is done on the doubly truncated product.

\begin{definition}[Substitution $z_1\mapsto z_0+z_2$; composite field]\label{l03:def:subst2}
(a) A two-variable distribution is \emph{lower-truncated in $z_1$} if its $z_1$-support is bounded
below uniformly in the $z_2$-power. (b) For $F$ lower-truncated in both variables,
$\iota_{z_2,z_0}F(z_0+z_2,z_2)$ is the distribution whose $z_0^az_2^b$-coefficient is
$\sum_{k}\binom{a+k}{a}F(a+k,b-k)$ (expand $(z_0+z_2)^{m}$ in nonnegative powers of $z_0$). (c) The
\emph{composite field} $Y(Y(u,z_0)v,z_2)x$ has $z_0^az_2^b$-coefficient $(u_{(-a-1)}v)_{(-b-1)}x$.
\end{definition}
\begin{leancode}
def IsBddBelow₁ (f : FormalDist₂ 𝒲) : Prop := ∃ B : ℤ, ∀ m n : ℤ, m < B → f m n = 0

noncomputable def subst2 (f : FormalDist₂ 𝒲) : FormalDist₂ 𝒲 :=
  fun a b => ∑ᶠ k : ℤ, cbz (a + k) a • f (a + k) (b - k)

noncomputable def comp (W : VertexOperatorAlgebra 𝒱) (u v x : 𝒱) : FormalDist₂ 𝒱 :=
  fun a b => mode (W.Y (W.smul u (-a-1) v)) (-b-1) x
\end{leancode}
File \leanpath{VOA/\allowbreak Core/\allowbreak Skew.lean}.

\begin{dictionary}
In \lean{subst2}, the term $F(m,n)z_1^mz_2^n$ with $z_1 = z_0+z_2$ contributes
$\binom ma z_0^az_2^{m-a+n}$; matching $z_0^az_2^b$ gives $m = a+k$, $n = b-k$, coefficient
\lean{cbz (a+k) a} $=\binom{a+k}a$ (zero for $a<0$). Without truncation in both variables the sum
over $k$ would be infinite; the library proves finiteness on the doubly truncated class
(\lean{subst2_summand_finite}). \lean{comp}: \lean{W.smul u (-a-1) v} is the $z_0^a$-coefficient of
$Y(u,z_0)v$, and \lean{mode (Y w) (-b-1) x} is the $z_2^b$-coefficient of $Y(w,z_2)x$.
\end{dictionary}

\begin{physics}
$z_1 = z_0+z_2$ re-expands a correlator around the second insertion point: $z_0$ is the separation
$z_1-z_2$, small compared with $z_2$. The composite field is the OPE $u\times v$ evaluated first and
then inserted at $z_2$.
\end{physics}

\begin{theorem}[FLM associativity]\label{l03:thm:flm}
In a VOA, let $N$ satisfy $(z_1-z_2)^N[Y(u,z_1),Y(v,z_2)] = 0$. Then for every $x$,
\begin{equation*}
\iota_{z_2,z_0}\Big((z_1-z_2)^NY(u,z_1)Y(v,z_2)x\Big)\Big|_{z_1 = z_0+z_2} = z_0^N\,Y(Y(u,z_0)v,z_2)x .
\end{equation*}
\end{theorem}
\begin{leancode}
theorem flm (W : VertexOperatorAlgebra 𝒱) (u v x : 𝒱) {N : ℕ}
    (hN : mulBySubPow N (bracketDist (W.Y u) (W.Y v)) = 0) (a b : ℤ) :
    subst2 (mulBySubPow N (W.toPreVertexAlgebra.Yprod u v x)) a b = comp W u v x (a - N) b
\end{leancode}
File \leanpath{VOA/\allowbreak Core/\allowbreak FLM.lean}.

\begin{dictionary}
The right side \lean{comp W u v x (a - N) b} is the $z_0^{a-N}z_2^b$-coefficient of the composite,
i.e.\ the $z_0^az_2^b$-coefficient of $z_0^N\,Y(Y(u,z_0)v,z_2)x$. The left side is legitimate because
$(z_1-z_2)^NY(u,z_1)Y(v,z_2)x$ is lower-truncated in both variables (in $z_2$ by VA5; in $z_1$ because
it equals $(z_1-z_2)^NY(v,z_2)Y(u,z_1)x$ by locality, and that is truncated in $z_1$ by VA5 for
$(u,x)$; lemma \lean{exists_mulBySubPow_Yprod_doublyBdd}). The hypothesis \lean{hN} is locality for
$(u,v)$ with an explicit order --- not an extra assumption, since VA4 provides such an $N$. The statement
holds for every integer $a$ (both the ``singular'' $a<N$ and ``normal'' $a\ge N$ regimes).

\textbf{Comparison with FLM.} FLM (§§8.8--8.10) and Lepowsky--Li (§3.6, ``weak associativity'') use the
factor $(z_0+z_2)^{k}$ with $k$ depending on $(u,x)$; the library's version uses $(z_1-z_2)^N =
z_0^N$ with $N$ depending on $(u,v)$. The two are equivalent formulations of the same associativity
(each follows from the Jacobi identity by multiplying by the appropriate power and taking a residue);
neither is used as an axiom.
\end{dictionary}

\begin{physics}
Associativity of the OPE in its most concrete form: the product $u(z_1)v(z_2)$ acting on $x$,
re-expanded around $z_2$ in the separation $z_0 = z_1-z_2$, is the OPE of $u$ and $v$ followed by
insertion at $z_2$ --- once the poles on the diagonal have been cleared by $(z_1-z_2)^N$. This is the
statement that the three-point function is a single rational function with the correct expansions in
the two channels.
\end{physics}

\textbf{Proof on paper.} Expand the left side by the binomial theorem and compare
coefficients. For $a\ge N$ (the ``normal'' columns, $z_0$-power $a-N\ge0$, i.e.\ negative products)
the composite coefficient is a normal-ordered product by Theorem~\ref{l03:thm:assocneg}, whose
coefficients are exactly what the substitution produces after the $(z_1-z_2)^N$ prefactor has been
absorbed (\lean{flm_normal}); the remaining boundary terms are shown to vanish because the relevant
binomials $\binom{p}{a-N}$ with $0\le p<a-N$ are zero (\lean{FLMBoundaryResidual}). For $a<N$ (the
singular columns, nonnegative products) the composite coefficient is $\Res_{z_1}(z_1-z_2)^{N-a-1}$ of the
commutator by Theorem~\ref{l03:thm:nthproduct}, and the substitution of the commutator, killed by
$(z_1-z_2)^N$, reproduces it (\lean{flm_lt}). \qed

\begin{status}
Unconditional, axiom-clean (\texttt{\#print axioms} \lean{VOA.VertexOperatorAlgebra.flm} run for this
chapter). The file is imported by \leanpath{WeakAssoc.lean}.
\end{status}

%======================================================================
\section{Commutative vertex algebras (Gannon Theorem 5.1.4)}\label{l03:sec:comm}
%======================================================================

\begin{definition}[Commutative; the classical product]\label{l03:def:comm}
A vertex algebra is \emph{commutative} if $[Y(u,z_1),Y(v,z_2)] = 0$ for all $u,v$. Its
\emph{classical product} is $u\cdot v := u_{(-1)}v$.
\end{definition}
\begin{leancode}
def IsCommutative (V : PreVertexAlgebra 𝒱) : Prop :=
  ∀ u v : 𝒱, bracketDist (V.Y u) (V.Y v) = 0

def cmul (V : PreVertexAlgebra 𝒱) (u v : 𝒱) : 𝒱 := V.smul u (-1) v

@[simp] theorem vacuum_cmul (V : PreVertexAlgebra 𝒱) (u : 𝒱) : V.cmul V.vacuum u = u
\end{leancode}
File \leanpath{VOA/\allowbreak Core/\allowbreak Commutativity.lean} (also \lean{cmul_vacuum}:
$u\cdot\lthreevac = u$).

\begin{dictionary}
\lean{IsCommutative} is locality with order $0$ for every pair. The counterexample of
Section~\ref{l03:sec:counter} is commutative.
\end{dictionary}

\begin{physics}
A commutative vertex algebra is a chiral algebra with no singular OPEs: all fields are holomorphic
at coincident points, as for the algebra of functions on the jet space of a target (the classical
limit of a chiral algebra, or the chiral algebra of a topologically twisted theory in which all
fields are regular).
\end{physics}

\begin{theorem}[Gannon Theorem 5.1.4 for VOAs]\label{l03:thm:commalg}
For a VOA $W$ the following are equivalent: (1) $W$ is commutative; (3) $u_{(n)}v = 0$ for all
$u,v$ and $n\ge0$. If they hold, then $u\cdot v = u_{(-1)}v$ is commutative and associative with
unit $\lthreevac$, so $W$ is a commutative ring, and $T = L_{-1}$ is a derivation of it (the derivation
law holds in every VOA).
\end{theorem}
\begin{leancode}
theorem comm_iff_nonneg_modes (W : VertexOperatorAlgebra 𝒱) :
    W.toPreVertexAlgebra.IsCommutative ↔ ∀ (u v : 𝒱) (n : ℤ), 0 ≤ n → W.smul u n v = 0

theorem Lneg1_cmul_leibniz (W : VertexOperatorAlgebra 𝒱) (u v : 𝒱) :
    W.L (-1) (W.toPreVertexAlgebra.cmul u v)
      = W.toPreVertexAlgebra.cmul (W.L (-1) u) v + W.toPreVertexAlgebra.cmul u (W.L (-1) v)

theorem cmul_assoc_of_comm (W : VertexOperatorAlgebra 𝒱)
    (hcomm : W.toPreVertexAlgebra.IsCommutative) (u v w : 𝒱) :
    W.toPreVertexAlgebra.cmul (W.toPreVertexAlgebra.cmul u v) w
      = W.toPreVertexAlgebra.cmul u (W.toPreVertexAlgebra.cmul v w)

theorem cmul_comm_of_comm (W : VertexOperatorAlgebra 𝒱)
    (hcomm : W.toPreVertexAlgebra.IsCommutative) (u v : 𝒱) :
    W.toPreVertexAlgebra.cmul u v = W.toPreVertexAlgebra.cmul v u

def commRing_of_comm (W : VertexOperatorAlgebra 𝒱)
    (hcomm : W.toPreVertexAlgebra.IsCommutative) : CommRing 𝒱 :=

theorem commRing_of_comm_isDerivation (W : VertexOperatorAlgebra 𝒱)
    (hcomm : W.toPreVertexAlgebra.IsCommutative) (a b : 𝒱) :
    letI := commRing_of_comm W hcomm
    W.L (-1) (a * b) = (W.L (-1) a) * b + a * (W.L (-1) b)

theorem comm_iff_assoc_alg (W : VertexOperatorAlgebra 𝒱) :
    W.toPreVertexAlgebra.IsCommutative ↔
      ((∀ (u v : 𝒱) (n : ℤ), 0 ≤ n → W.smul u n v = 0)
        ∧ (∀ u v : 𝒱, W.toPreVertexAlgebra.cmul u v = W.toPreVertexAlgebra.cmul v u)
        ∧ (∀ u v w : 𝒱, W.toPreVertexAlgebra.cmul (W.toPreVertexAlgebra.cmul u v) w
              = W.toPreVertexAlgebra.cmul u (W.toPreVertexAlgebra.cmul v w)))
\end{leancode}
File \leanpath{VOA/\allowbreak Core/\allowbreak Commutativity.lean}.

\begin{dictionary}
\textbf{\lean{commRing_of_comm}} builds a Mathlib \lean{CommRing} structure on the underlying additive
group with multiplication \lean{cmul} and one \lean{vacuum}; its fields (associativity, commutativity,
unit laws, distributivity) are the theorems listed (distributivity is bilinearity of \lean{smul}).
\lean{letI := …} makes this ring structure the one used by \lean{*} in the next line. It is a
\lean{def} whose type mentions \lean{hcomm}, so it exists only for commutative VOAs.

\textbf{What \lean{comm_iff_assoc_alg} says, precisely.} The right side is the conjunction of
condition (3) with commutativity and associativity of \lean{cmul}. The direction $\Leftarrow$ uses only
the first conjunct (it is (3)$\Rightarrow$(1)). So the theorem is ``(1)$\iff$(3), and (1) implies that
\lean{cmul} is commutative and associative''. Gannon's converse --- \emph{every} commutative
associative unital algebra with a derivation $T$ is a commutative vertex algebra via
$Y(u,z)v = (e^{zT}u)v$ --- is \textbf{not} stated in the library. It is in any case not a statement
about VOAs (a commutative differential algebra has no conformal vector in general), and at the level
of \lean{PreVertexAlgebra} the correspondence fails (Section~\ref{l03:sec:counter}).
\end{dictionary}

\begin{physics}
A chiral algebra with no singular OPEs is the same thing as a commutative algebra of local operators
with an action of $\partial$: the OPE reduces to the ordinary product, $u(z)v(w) = {:}u(z)v(w){:}$, and
all the structure is the commutative product and the derivative. In a VOA the stress tensor always
has a singular OPE with itself unless everything is trivial, which is the content of the warning
below.
\end{physics}

\textbf{Proof on paper.} (1)$\Rightarrow$(3): by Theorem~\ref{l03:thm:nthproduct},
$Y(u_{(k)}v) = \Res_{z_1}(z_1-z)^k\cdot0 = 0$, so $u_{(k)}v = (u_{(k)}v)_{(-1)}\lthreevac = 0$.
(3)$\Rightarrow$(1): every term of the commutator formula vanishes. Associativity: by
Theorem~\ref{l03:thm:assocneg} with $k = 0$, $(u_{(-1)}v)_{(-1)}w$ is the $z^0$-coefficient of
${:}Y(u,z)Y(v,z){:}w$, and under (3) only the term $u_{(-1)}v_{(-1)}w$ survives. Commutativity: by
skew-symmetry at $n=-1$, $u_{(-1)}v = \sum_j\frac{(-1)^j}{j!}T^j(v_{(j-1)}u)$, and under (3) only $j = 0$
survives. Leibniz: Theorem~\ref{l03:thm:transcov} at $n = -1$. \qed

\begin{status}
All unconditional, axiom-clean (\texttt{\#print axioms} for \lean{comm_iff_nonneg_modes},
\lean{comm_iff_assoc_alg}, \lean{commRing_of_comm} run for this chapter).
\end{status}

\begin{warning}[title={Theorem 5.1.4 at the VOA level is a statement about the trivial VOA}]
Let $W$ be a commutative \lean{VertexOperatorAlgebra}. By (3), $L_0 = \omega_{(1)} = 0$ on $W$; by
(VOA2), $L_0 = n$ on $W_n$, so $W_n = 0$ for $n\ne0$, and by (VOA4) $W = W_0 = \C\lthreevac$ (or $W = 0$).
Then $\omega\in W_2 = 0$, and if $W\ne0$ the Virasoro relation $[L_2,L_{-2}] = 4L_0+\frac c2$ forces
$c = 0$. So the hypothesis \lean{IsCommutative} is satisfiable only by the one-dimensional VOA
$\C\lthreevac$ (and the zero VOA), and all theorems of this section are, for VOAs, statements about that
algebra. They are true and correctly proved, but they do \emph{not} formalise the classical
content of Gannon's Theorem~5.1.4 (commutative vertex algebras $=$ commutative differential
algebras), which lives at the level of vertex algebras with translation covariance --- a level the
library does not have a correct structure for (Section~\ref{l03:sec:counter}). The source docstring
does not mention this.
\end{warning}

%======================================================================
\section[Reconstruction]{Reconstruction: building a vertex algebra from generating fields}\label{l03:sec:recon}
%======================================================================

The theorems so far start from a vertex algebra. Constructing one --- the Heisenberg, lattice,
affine, Virasoro and Moonshine VOAs of later chapters --- requires the converse direction:
produce a state--field map from a few generating fields. The library has two reconstruction theorems:
\lean{Reconstruction.reconstruct} (Theorem~\ref{l03:thm:reconstruct} below) and its word-field corollary
\lean{NonlinearConformal.reconstruct_of_generators}, which belongs to Chapter~\ref{ch:L11}
(Theorem~\ref{l11:thm:recongen}).
Here we present the ingredients in \leanpath{VOA/\allowbreak Core/\allowbreak ReconstructionExistence.lean}
and say exactly what the theorem delivers.

\begin{definition}[$T$-covariance]\label{l03:def:dcov}
For $T\in\End V$, a distribution $a$ is $T$-covariant if $[T,a(z)] = \partial_za(z)$, i.e.\
$[T,a_m] = (m+1)a_{m+1}$ for the $z^m$-coefficients.
\end{definition}
\begin{leancode}
def DCovariant (T : End ℂ V) (a : FormalDist (End ℂ V)) : Prop :=
  ∀ m : ℤ, T * a m - a m * T = ((m + 1 : ℤ)) • a (m + 1)
\end{leancode}
File \leanpath{VOA/\allowbreak Core/\allowbreak ReconstructionExistence.lean}
(\lean{VOA.Reconstruction.DCovariant}). The same file proves closure under scalars, sums, finite
linear combinations, $\partial$ and ${:}\cdot\,\cdot{:}$ (\lean{DCovariant.smul}, \lean{.add},
\lean{.finsetSum}, \lean{.deriv'}, \lean{.nprod}, the last via the Leibniz rule \lean{deriv_nprod}).

\begin{dictionary}
The $z^m$-coefficient of $\partial_za$ is $(m+1)a_{m+1}$, so the definition is $[T,a(z)] = \partial a(z)$
exactly: Kac's and FBZ's translation axiom for a single field. Here $T$ is an arbitrary operator, not
$u_{(-2)}\lthreevac$.
\end{dictionary}

\begin{physics}
$T$-covariance says that $a$ is a translation-covariant local operator: $e^{\epsilon T}a(z)e^{-\epsilon T}
= a(z+\epsilon)$. This is precisely the property missing from the axioms of
Definition~\ref{l03:def:VA}; in the reconstruction theorem it is \emph{required of the generators}.
\end{physics}

\begin{lemma}[The creation tower is forced]\label{l03:lem:tower}
If $F,G$ are $T$-covariant, $T\lthreevac = 0$, $F(z)\lthreevac,G(z)\lthreevac\in V[[z]]$ and their constant
terms agree, then $F(z)\lthreevac = G(z)\lthreevac$.
\end{lemma}
\begin{leancode}
theorem tower_unique {T : End ℂ V} {F G : FormalDist (End ℂ V)} {vac : V}
    (hF : DCovariant T F) (hG : DCovariant T G) (hTvac : T vac = 0)
    (hbase : F 0 vac = G 0 vac)
    (hFneg : ∀ q : ℤ, q < 0 → F q vac = 0) (hGneg : ∀ q : ℤ, q < 0 → G q vac = 0)
    (p : ℤ) : F p vac = G p vac
\end{leancode}
File \leanpath{VOA/\allowbreak Core/\allowbreak ReconstructionExistence.lean}.

\begin{dictionary}
All hypotheses are as stated; the conclusion is the equality of every $z^p$-coefficient on the vacuum.
\end{dictionary}

\begin{physics}
For a translation-covariant field, $a(z)|0\rangle = e^{zT}a_{(-1)}|0\rangle$: the state at $z$ is the
translate of the state at $0$. This is what fails in the counterexample, whose $Y(x,z)$ is not
$T$-covariant for any $T$.
\end{physics}

\textbf{Proof on paper.} Apply $[T,F_p] = (p+1)F_{p+1}$ to $\lthreevac$ and use $T\lthreevac = 0$:
$(p+1)F_{p+1}\lthreevac = TF_p\lthreevac$. Induction on $p\ge0$. \qed

\begin{theorem}[Operator-first uniqueness]\label{l03:thm:fieldzero}
Let $V$ be a vector space with a vector $\lthreevac$, and let $(s_j,f_j)_{j\in J}$ be states and distributions with
$f_j(z)\lthreevac\in V[[z]]$ and $f_j(z)\lthreevac|_{z=0} = s_j$, such that the $s_j$ span $V$. If a distribution
$A\in\End(V)[[z^{\pm1}]]$ is mutually local with every $f_j$ and $A(z)\lthreevac = 0$, then $A = 0$.
\end{theorem}
\noindent File \leanpath{VOA/\allowbreak Core/\allowbreak ReconstructionExistence.lean}.
\begin{leancode}
theorem field_eq_zero_of_vacuum {ι : Type*} (vac : V) (A : FormalDist (End ℂ V))
    (st : ι → V) (fld : ι → FormalDist (End ℂ V))
    (hloc : ∀ j, Local A (fld j))
    (hcre0 : ∀ j, (fld j) 0 vac = st j)
    (hcreneg : ∀ j, ∀ q : ℤ, q < 0 → (fld j) q vac = 0)
    (hAvac : ∀ p : ℤ, A p vac = 0)
    (hspan : Submodule.span ℂ (Set.range st) = ⊤) :
    A = 0
\end{leancode}

\begin{dictionary}
\lean{ι} is an index type; \lean{st j} is the state $s_j$ and \lean{fld j} the distribution $f_j$. \lean{(fld j) 0 vac = st j}
says the $z^0$-coefficient of $f_j(z)\lthreevac$ is $s_j$, i.e.\ the mode $(f_j)_{(-1)}\lthreevac = s_j$;
\lean{hcreneg} says there are no negative powers of $z$ in $f_j(z)\lthreevac$, so that $f_j(z)\lthreevac\in V[[z]]$;
\lean{hAvac : ∀ p, A p vac = 0} says $A(z)\lthreevac = 0$ coefficientwise; \lean{Local A (fld j)} is
Definition~\ref{l01:def:bracket} (some power of $z_1-z_2$ kills the commutator); \lean{hspan} says the states span $V$
(\lean{Set.range st} is the set of states, \lean{⊤} the whole space). \textbf{Hypotheses:} exactly those of the
mathematical statement. \textbf{Scope.} No vertex algebra is assumed: $V$ is a bare vector space, and neither $A$ nor
the $f_j$ is assumed to be a field. It is the spanning-family strengthening of the abstract Goddard theorem
(Theorem~\ref{l03:thm:goddard}(iii), \lean{goddard_zero}), in which the test fields were required for \emph{every} vector
$b\in V$; here only a spanning family is needed. The conclusion $A=0$ is an equality of distributions. \textbf{Could it be
vacuous?} The hypotheses are satisfiable
non-trivially: in any VOA take the $f_j$ to be the vertex operators $Y(s_j,\cdot)$ of a spanning family $(s_j)$, which
are local with $Y(a,\cdot)$ for every $a$. The theorem has content because locality against a spanning family is a
strong condition; it is used by the reconstruction theorem (Theorem~\ref{l03:thm:reconstruct}) to show that the
extension $Y$ is well defined.
\end{dictionary}

\begin{physics}
An operator that is mutually local with all the fundamental fields and gives zero on the vacuum is the zero
operator. Heuristically this is the chiral cousin of the Reeh--Schlieder property in axiomatic quantum field
theory (the vacuum is cyclic and separating for local operators): a local operator is determined by the state it
creates from the vacuum. Here it is a theorem of formal calculus, and it is the technical heart of the uniqueness part
of the state--field correspondence: the field of a state is determined by the state.
\end{physics}

\textbf{Proof on paper.} Fix $j$ and put $G(z,w) := A(z)f_j(w)\lthreevac\in\End(V)[[z^{\pm1}]]\otimes V[[w]]$. Mutual
locality gives $N$ with $(z-w)^N[A(z),f_j(w)] = 0$; applying this to $\lthreevac$ and using $A(z)\lthreevac = 0$ gives
$(z-w)^NG(z,w) = (z-w)^N\big(A(z)f_j(w)\lthreevac - f_j(w)A(z)\lthreevac\big) = 0$. The series $G$ contains only
nonnegative powers of $w$ (because $f_j(w)\lthreevac\in V[[w]]$), so it is lower truncated in $w$, and by the
cancellation theorem (Theorem~\ref{l01:thm:cancel}: $z-w$ is not a zero divisor on lower-truncated series) $G=0$.
Taking the coefficient of $w^0$, which is $A(z)s_j$, gives $A(z)s_j = 0$ for all $j$. Hence every coefficient
$A_m$ vanishes on the span of the $s_j$, which is $V$, so $A=0$. \qed

\begin{status}
Unconditional; \texttt{\#print axioms} \lean{VOA.Reconstruction.field\_eq\_zero\_of\_vacuum} $=$
\lean{[propext, Classical.}\allowbreak\lean{choice, Quot.}\allowbreak\lean{sound]} (run for this chapter). The PBW theorem \lean{VOA.pbw} is not used.
\end{status}

\begin{theorem}[Operator-first reconstruction]\label{l03:thm:reconstruct}
Let $V$ be a vector space, $\lthreevac\in V$, $T\in\End V$ with $T\lthreevac = 0$, and $(s_i,f_i)_{i\in I}$
states and distributions such that each $f_i$ is a field, the $f_i$ are pairwise local and
$T$-covariant, $f_i(z)\lthreevac\in V[[z]]$ with constant term $s_i$, and the $s_i$ span $V$. Then there is
a linear $Y:V\to\End(V)[[z^{\pm1}]]$ with $Y(s_i,\cdot) = f_i$, all $Y(u,\cdot)$ fields, pairwise
local, $T$-covariant, with $Y(u,z)\lthreevac\in u+zV[[z]]$.
\end{theorem}
\begin{leancode}
theorem reconstruct {ι : Type*} (vac : V) (st : ι → V) (fld : ι → FormalDist (End ℂ V))
    (T : End ℂ V)
    (hspan : Submodule.span ℂ (Set.range st) = ⊤)
    (hfield : ∀ i, IsField (fld i))
    (hlocal : ∀ i j, Local (fld i) (fld j))
    (hcre0 : ∀ i, (fld i) 0 vac = st i)
    (hcreneg : ∀ i, ∀ q : ℤ, q < 0 → (fld i) q vac = 0)
    (hTvac : T vac = 0)
    (hTcov : ∀ i, DCovariant T (fld i)) :
    ∃ YA : V →ₗ[ℂ] FormalDist (End ℂ V),
      (∀ i, YA (st i) = fld i) ∧
      (∀ u : V, IsField (YA u)) ∧
      (∀ u v : V, Local (YA u) (YA v)) ∧
      (∀ u : V, (YA u) 0 vac = u) ∧
      (∀ u : V, ∀ q : ℤ, q < 0 → (YA u) q vac = 0) ∧
      (∀ u : V, DCovariant T (YA u))
\end{leancode}
File \leanpath{VOA/\allowbreak Core/\allowbreak ReconstructionExistence.lean}
(\lean{VOA.Reconstruction.reconstruct}).

\begin{dictionary}
\textbf{What it delivers.} The six conjuncts are: $Y$ extends the given fields; fields (VA5); locality
(VA4); creation (VA3, both halves); $T$-covariance. In particular, through \lean{tower_unique}, the
conclusion includes $Y(u,z)\lthreevac = e^{zT}u$, i.e.\ the full Kac/FBZ axiom system \emph{minus} the
vacuum axiom $Y(\lthreevac) = \mathrm{id}$, which the user must check (it holds as soon as $\lthreevac$ is
among the $s_i$ with $f = \mathrm{id}$).

\textbf{What it does not deliver.} No grading (VA1) and no \lean{PreVertexAlgebra} structure; the user
assembles one. Uniqueness of $Y$ is not part of the statement (the docstring says ``unique''; the
statement is $\exists$), although it follows from \lean{field_eq_zero_of_vacuum}.

\textbf{Comparison with Kac's existence theorem.} Kac (Thm.~4.5) and FBZ (Thm.~4.4.1) start from a
family of generating fields and assume that the \emph{monomials} $a^{i_1}_{(-j_1-1)}\cdots
a^{i_r}_{(-j_r-1)}\lthreevac$ span; the fields of the monomials are then \emph{constructed} as
normal-ordered products. The library's theorem instead takes as input a family whose states already
span, each state with its field. It is therefore a reduction of Kac's theorem to Dong's lemma: one
applies it to the family of all iterated normal-ordered word fields (which are pairwise local by
Dong's lemma and $T$-covariant by \lean{DCovariant.nprod}); this is done in
\lean{reconstruct_of_generators} (Theorem~\ref{l11:thm:recongen} in Chapter~\ref{ch:L11}).

\textbf{Translation covariance is a hypothesis here.} This is the correct design: by
Section~\ref{l03:sec:counter}, VA1--VA5 do not imply it, so a construction must supply it.
\end{dictionary}

\begin{physics}
``A chiral algebra is generated by a few fields and their OPEs.'' One specifies the currents (or the
stress tensor), checks that their OPEs have finite poles (locality), that $\partial$ acts as $[L_{-1},\cdot]$
and that the vacuum is translation invariant; then every composite operator is determined and the
whole operator algebra is consistent.
\end{physics}

\textbf{Proof on paper.} Define $Y$ on formal combinations $\sum c_is_i$ by $\sum c_if_i$. Well-defined:
if $\sum c_is_i = 0$, the field $A = \sum c_if_i$ is local with every $f_j$, annihilates $\lthreevac$ in
degree $<0$, has constant term $\sum c_is_i = 0$, and by $T$-covariance (Lemma~\ref{l03:lem:tower}
applied inductively) annihilates $\lthreevac$ in every degree; so $A = 0$ by Goddard
(Theorem~\ref{l03:thm:goddard}(iii) in its spanning form \lean{field_eq_zero_of_vacuum}). Every listed
property is linear in the field and holds for the $f_i$. \qed

\begin{status}
\lean{reconstruct}, \lean{tower_unique}, \lean{field_eq_zero_of_vacuum}: unconditional, axiom-clean
(\texttt{\#print axioms} \lean{VOA.Reconstruction.reconstruct} run for this chapter).
\end{status}

\subsection{Fields on a tensor product}

To build tensor-product VOAs before either factor is packaged as a vertex algebra, the library lifts
fields to $A\otimes B$.

\begin{definition}[First- and second-factor lifts]\label{l03:def:tensorlift}
For a distribution $f$ on $A$, $f\otimes1$ acts on $A\otimes B$ by $(f\otimes1)_n(a\otimes b) =
f_na\otimes b$; similarly $1\otimes g$. The tensor translation is $T_A\otimes1+1\otimes T_B$.
\end{definition}
\begin{leancode}
noncomputable def firstField (f : FormalDist (End ℂ A)) :
    FormalDist (End ℂ (A ⊗[ℂ] B)) :=
  fun n => LinearMap.rTensor B (f n)

noncomputable def secondField (g : FormalDist (End ℂ B)) :
    FormalDist (End ℂ (A ⊗[ℂ] B)) :=
  fun n => LinearMap.lTensor A (g n)

noncomputable def tensorTranslation (TA : End ℂ A) (TB : End ℂ B) :
    End ℂ (A ⊗[ℂ] B) :=
  LinearMap.rTensor B TA + LinearMap.lTensor A TB
\end{leancode}
File \leanpath{VOA/\allowbreak Core/\allowbreak TensorFieldLift.lean} (namespace
\lean{VOA.TensorFieldLift}).

\begin{dictionary}
\lean{A ⊗[ℂ] B} is Mathlib's tensor product over $\C$; \lean{LinearMap.rTensor B φ} $=\varphi\otimes\mathrm{id}_B$
and \lean{LinearMap.lTensor A ψ} $=\mathrm{id}_A\otimes\psi$.
\end{dictionary}

\begin{physics}
Two decoupled chiral algebras: fields of the first system act on the first factor of the product
Hilbert space, fields of the second on the second, and the total translation generator is the sum.
\end{physics}

\begin{theorem}[Lifts preserve the reconstruction hypotheses]\label{l03:thm:tensorlift}
Lifts of fields are fields; lifts of local pairs are local; a first-factor lift and a second-factor
lift always commute, hence are local; lifts of $T_A$- (resp.\ $T_B$-)covariant fields are covariant for
the tensor translation; and ${:}(f\otimes1)(1\otimes g){:}$ is the ordinary product.
\end{theorem}
\begin{leancode}
theorem isField_firstField {f : FormalDist (End ℂ A)} (hf : IsField f) :
    IsField (firstField (B := B) f)

theorem local_firstField {f g : FormalDist (End ℂ A)} (h : Local f g) :
    Local (firstField (B := B) f) (firstField g)

theorem local_firstField_secondField (f : FormalDist (End ℂ A))
    (g : FormalDist (End ℂ B)) :
    Local (firstField (B := B) f) (secondField (A := A) g) :=
  local_of_bracketDist_eq_zero (bracketDist_firstField_secondField f g)

theorem dcovariant_firstField {TA : End ℂ A} {TB : End ℂ B}
    {f : FormalDist (End ℂ A)} (hf : Reconstruction.DCovariant TA f) :
    Reconstruction.DCovariant (tensorTranslation TA TB) (firstField f)

theorem nprod_firstField_secondField_apply {f : FormalDist (End ℂ A)}
    {g : FormalDist (End ℂ B)} (hf : IsField f) (hg : IsField g) (s : ℤ)
    (x : A ⊗[ℂ] B) :
    nprod (isField_firstField (B := B) hf) (isField_secondField (A := A) hg) s x
      = ∑ᶠ m : ℤ,
          (firstField (B := B) f m * secondField (A := A) g (s - m)) x
\end{leancode}
File \leanpath{VOA/\allowbreak Core/\allowbreak TensorFieldLift.lean}. The second-factor analogues
(\lean{isField_secondField}, \lean{local_secondField}, \lean{dcovariant_secondField},
\lean{nprod_firstField}, \lean{nprod_secondField}) are in the same file.

\begin{dictionary}
\lean{(B := B)} only fixes the implicit second factor. \lean{isField_firstField} needs an argument by
induction on tensors (a general element of $A\otimes B$ is a finite sum of pure tensors). In the last
statement the right side is a \lean{finsum} whose support is finite because $f$, $g$ are fields.
\end{dictionary}

\begin{physics}
Fields of decoupled systems commute (they are mutually local with no singular OPE), and their
normal-ordered product is just the product, e.g.\ the stress tensor of a tensor product is
$T_A\otimes1+1\otimes T_B$ with no cross terms.
\end{physics}

\begin{status}
All unconditional, axiom-clean (\texttt{\#print axioms} \lean{local_firstField_secondField} run for this chapter).
\end{status}

%======================================================================
\section{Local systems of weak vertex operators}\label{l03:sec:local}
%======================================================================

Li (1996) constructs vertex algebras, and their modules, from \emph{local systems}: spaces of
mutually local fields on a fixed vector space $M$, closed under all field products. The space of
such fields is itself a vertex algebra, with $M$ as a module. The library has the first three of
the four steps of this construction. All results here are about an arbitrary complex vector space
$M$.

\subsection{Weak vertex operators and the generated system}

\begin{definition}[Weak vertex operators; the generated system]\label{l03:def:fld}
$\mathrm{Fld}(M)$ is the set of fields on $M$, with the identity field $1$, the normal-ordered product
and $\partial$. For $S\subseteq\mathrm{Fld}(M)$, $\mathrm{Gen}(S,i,a)$ means that $a$ is obtained from
elements of $S$ and $1$ by $i$ applications of ${:}\cdot\,\cdot{:}$ and $\partial$; $\mathrm{Generated}(S)$ is
the set of such $a$ (any $i$).
\end{definition}
\begin{leancode}
def Fld (M : Type*) [AddCommGroup M] [Module ℂ M] : Type _ :=
  {a : FormalDist (End ℂ M) // IsField a}

noncomputable def one : Fld M := ⟨idField, idField_isField⟩

noncomputable def nprod (a b : Fld M) : Fld M :=
  ⟨VOA.nprod a.2 b.2, a.2.nprod b.2⟩

noncomputable def deriv (a : Fld M) : Fld M :=
  ⟨VOA.deriv a.1, a.2.deriv⟩

inductive Gen (S : Set (Fld M)) : ℕ → Fld M → Prop
  | base {a : Fld M} : a ∈ S → Gen S 0 a
  | one : Gen S 0 Fld.one
  | nprod {i j : ℕ} {a b : Fld M} : Gen S i a → Gen S j b → Gen S (i + j + 1) (a.nprod b)
  | deriv {i : ℕ} {a : Fld M} : Gen S i a → Gen S (i + 1) a.deriv

def Generated (S : Set (Fld M)) : Set (Fld M) := {a | ∃ i, Gen S i a}
\end{leancode}
File \leanpath{VOA/\allowbreak Core/\allowbreak LocalSystem.lean} (\lean{Fld.one}, \lean{Fld.nprod},
\lean{Fld.deriv} in namespace \lean{VOA.Fld}).

\begin{dictionary}
\lean{{a // IsField a}} is the subtype: a pair of a distribution \lean{a.1} and a proof \lean{a.2}
that it is a field. \lean{Gen} is an \emph{inductive predicate}: its four constructors are the only
ways to establish \lean{Gen S i a}. The index $i$ counts operations; it is bookkeeping for the
locality induction below. Note that \lean{Generated} is closed under ${:}\cdot\,\cdot{:}$ and $\partial$,
hence under all negative products, but not under linear combinations or the nonnegative products;
the linear, all-product closure is the \lean{LocalFieldSystem} below.
\end{dictionary}

\begin{physics}
The operators one can build from a set of chiral fields by taking derivatives and normal-ordered
products --- the ``composite operators'' of the theory. No vacuum and no grading are involved: this
is the operator algebra acting on a fixed sector $M$.
\end{physics}

\begin{theorem}[The generated system is local]\label{l03:thm:genlocal}
If the elements of $S$ are pairwise local, so are all elements of $\mathrm{Generated}(S)$.
\end{theorem}
\begin{leancode}
theorem gen_local (S : Set (Fld M)) (hS : ∀ a ∈ S, ∀ b ∈ S, Local a.1 b.1) :
    ∀ (N i j : ℕ) (a b : Fld M), i + j ≤ N → Gen S i a → Gen S j b → Local a.1 b.1

theorem generated_local {S : Set (Fld M)} (hS : ∀ a ∈ S, ∀ b ∈ S, Local a.1 b.1)
    {a b : Fld M} (ha : a ∈ Generated S) (hb : b ∈ Generated S) : Local a.1 b.1
\end{leancode}
File \leanpath{VOA/\allowbreak Core/\allowbreak LocalSystem.lean}.

\begin{dictionary}
\lean{hS} includes $a = b$ (self-locality of each generator), which is genuinely required. The bound
$N$ in \lean{gen_local} is the induction parameter; \lean{generated_local} is the clean form.
\end{dictionary}

\begin{physics}
Composite operators of mutually local fields are mutually local: closure of the chiral operator algebra
under the OPE.
\end{physics}

\textbf{Proof on paper.} Strong induction on $i+j$. If $a = {:}a_1a_2{:}$, Dong's lemma
(\lean{dong_neg}) needs $a_1\sim a_2$, $a_1\sim b$, $a_2\sim b$, all of strictly smaller total size;
if $a = \partial a_1$, use \lean{Local.derivL}; if $a$ is a generator or $1$, peel $b$ instead
(using symmetry); two generators are local by \lean{hS}, and $1$ is local with everything. \qed

\begin{status}
Unconditional, axiom-clean (\texttt{\#print axioms} \lean{VOA.gen_local} run for this chapter).
\end{status}

\subsection{All integer products and closed local systems}

\begin{definition}[Integer field products; local field system]\label{l03:def:lfs}
For fields $a,b$ on $M$ and $n\in\Z$: $a_{(n)}b := \Res_{z_1}(z_1-z_2)^n[a(z_1),b(z_2)]$ if $n\ge0$ and
$a_{(n)}b := \frac1{(-n-1)!}{:}\partial^{-n-1}a\,b{:}$ if $n<0$. A \emph{local field system} is a linear
subspace of $\mathrm{Fld}(M)$ containing $1$, pairwise local, and closed under all $a_{(n)}b$.
\end{definition}
\begin{leancode}
noncomputable def prod (n : ℤ) (a b : Fld M) : Fld M :=
  if _hn : 0 ≤ n then
    ⟨cCoeff n.toNat (bracketDist a.1 b.1), isField_cCoeff_bracket a.2 b.2 n.toNat⟩
  else
    ⟨negProd a.2 b.2 (-n - 1).toNat, isField_negProd a.2 b.2 (-n - 1).toNat⟩

structure LocalFieldSystem (M : Type*) [AddCommGroup M] [Module ℂ M] where
  carrier : Submodule ℂ (Fld M)
  one_mem : Fld.one ∈ carrier
  locality : ∀ a ∈ carrier, ∀ b ∈ carrier, Local a.1 b.1
  prod_mem : ∀ (n : ℤ) {a b : Fld M}, a ∈ carrier → b ∈ carrier → Fld.prod n a b ∈ carrier
\end{leancode}
File \leanpath{VOA/\allowbreak Core/\allowbreak LocalSystemBare.lean} (\lean{VOA.Fld.prod},
\lean{VOA.LocalFieldSystem}; $\mathrm{Fld}(M)$ is given its vector-space structure in the same file).

\begin{dictionary}
\lean{n.toNat} is $n$ for $n\ge0$; \lean{(-n - 1).toNat} is $-n-1$ for $n<0$. The two field-property
proofs (\lean{isField_cCoeff_bracket}, \lean{isField_negProd}) show that products of fields are fields.
Same as Li's definition (and Kac's $n$-th products of fields).
\end{dictionary}

\begin{physics}
A maximal consistent set of chiral operators acting on a sector: closed under taking every OPE
coefficient, containing the identity.
\end{physics}

\begin{theorem}[Basic laws of field products]\label{l03:thm:fldlaws}
For fields on $M$: $1_{(n)}a = \delta_{n,-1}a$; $a_{(k)}1 = 0$ ($k\ge0$), $a_{(-1)}1 = a$,
$a_{(-k-1)}1 = \frac1{k!}\partial^ka$; if $(z_1-z_2)^N[a,b] = 0$ then $a_{(n)}b = 0$ for $n\ge N$; and
products of pairwise local fields are local with a third field local with both.
\end{theorem}
\begin{leancode}
theorem prod_one_left_all (n : ℤ) (a : Fld M) :
    prod n one a = if n = -1 then a else 0
\end{leancode}
File \leanpath{VOA/\allowbreak Core/\allowbreak LocalSystemAction.lean} (namespace \lean{VOA.Fld}).
\begin{leancode}
theorem prod_one_right_negsucc (k : ℕ) (a : Fld M) :
    (prod (-(k : ℤ) - 1) a one).1 =
      (k.factorial : ℂ)⁻¹ • (VOA.deriv^[k]) a.1

theorem prod_eventually_zero_of_local {a b : Fld M} (h : Local a.1 b.1) :
    ∃ N : ℤ, ∀ n : ℤ, N ≤ n → prod n a b = 0
\end{leancode}
File \leanpath{VOA/\allowbreak Core/\allowbreak LocalSystemBare.lean} (also \lean{prod_one_right_nat},
\lean{prod_one_right_neg_one}).
\begin{leancode}
theorem prod_local_right_of_pairwise {a b c : Fld M}
    (hab : Local a.1 b.1) (hac : Local a.1 c.1) (hbc : Local b.1 c.1) (n : ℤ) :
    Local (prod n a b).1 c.1
\end{leancode}
File \leanpath{VOA/\allowbreak Core/\allowbreak LocalSystemRegular.lean}.

\begin{dictionary}
These are the vacuum, creation and truncation laws of a vertex algebra, for the system of fields with
$1$ as vacuum; \lean{prod_local_right_of_pairwise} is Dong's lemma for every $n\in\Z$ (negative $n$ via
\lean{dong_neg} and \lean{Local.deriv_iterate_left}). Note the translation built in:
$a_{(-2)}1 = \partial a$.
\end{dictionary}

\begin{physics}
In the space of operators, the identity plays the vacuum and $\partial$ the translation generator:
$a_{(-k-1)}1 = \partial^ka/k!$ is the Taylor expansion of the operator itself.
\end{physics}

\begin{status}
All unconditional, axiom-clean (\texttt{\#print axioms} run for
\lean{Fld.prod_local_right_of_pairwise} and \lean{LocalFieldSystem.creation_deriv}).
\end{status}

\subsection{The Borcherds identity for two local fields}

\begin{theorem}[Borcherds identity for an arbitrary local pair]\label{l03:thm:fieldjacobi}
Let $a,b$ be mutually local fields on an arbitrary vector space $M$. Then for every $x\in M$ and all
$\ell,m,n\in\Z$,
\begin{equation*}
\sum_{i\ge0}(-1)^i\binom\ell i\Big(a_{(\ell+m-i)}b_{(n+i)}x - (-1)^\ell b_{(\ell+n-i)}a_{(m+i)}x\Big)
= \sum_{i\ge0}\binom mi(a_{(\ell+i)}b)_{(m+n-i)}x,
\end{equation*}
where $a_{(\ell+i)}b$ is the field product of Definition~\ref{l03:def:lfs}. Its $m = 0$ case, the iterate
formula, holds for any two fields (locality not needed).
\end{theorem}
\begin{leancode}
theorem mode_prod_iterate (a b : Fld M) (ell N : ℤ) (x : M) :
    mode (prod ell a b).1 N x =
      ∑ᶠ i : ℕ, ((-1 : ℂ) ^ i * gBinom ell i) •
        (mode a.1 (ell - i) (mode b.1 (N + i) x)
          - ((-1 : ℂ) ^ ell) • mode b.1 (ell + N - i) (mode a.1 i x))
\end{leancode}
File \leanpath{VOA/\allowbreak Core/\allowbreak FieldProductIterate.lean}.
\begin{leancode}
theorem field_jacobi_of_local {a b : Fld M} (h : Local a.1 b.1)
    (x : M) (ell m n : ℤ) :
    (∑ᶠ i : ℤ, (negOnePow i * cbz ell i) •
      (mode a.1 (ell + m - i) (mode b.1 (n + i) x) -
        negOnePow ell • mode b.1 (ell + n - i) (mode a.1 (m + i) x))) =
      ∑ᶠ i : ℤ, cbz m i • mode (prod (ell + i) a b).1 (m + n - i) x
\end{leancode}
File \leanpath{VOA/\allowbreak Core/\allowbreak LocalSystemJacobi.lean} (also the packaged form
\lean{field_borcherds_of_local}).

\begin{dictionary}
The left side is literally the left side of \lean{VertexOperatorAlgebra.jacobi} with \lean{W.smul u}
replaced by \lean{mode a.1}; the right side has the field product \lean{prod (ell + i) a b} where
\lean{jacobi} has $Y(u_{(\ell+i)}v)$. The only hypothesis is \lean{Local a.1 b.1}.

\textbf{How this sharpens the counterexample.} Apply the theorem to $a = Y(u)$, $b = Y(v)$ in a
\lean{PreVertexAlgebra}: the Borcherds identity of the algebra holds for all $(\ell,m,n)$ \emph{if and only
if} $Y(u_{(\ell)}v) = Y(u)_{(\ell)}Y(v)$ for all $\ell$ (take $m = 0$ for ``only if''). For $\ell\ge0$ this
needs only Goddard and the creation tower; for $\ell<0$ it needs the creation tower $Y(u,z)\lthreevac =
e^{zT}u$ --- translation covariance. In Section~\ref{l03:sec:counter}, $Y(x)_{(-2)}Y(\lthreevac) = \partial
Y(x) = 2z\,y\cdot{}\ne0 = Y(x_{(-2)}\lthreevac)$.
\end{dictionary}

\begin{physics}
OPE associativity is automatic \emph{for operators}: any two mutually local chiral operators satisfy
the Borcherds identity with their operator OPE coefficients. What a vertex algebra adds is the
state--operator correspondence, and the content of the Borcherds identity of a vertex algebra is that
this correspondence intertwines the OPEs of states with those of operators.
\end{physics}

\textbf{Proof on paper.} Identical in structure to Theorem~\ref{l03:thm:borcherds}: the iterate formula
($m=0$) is a direct coefficient computation from the definitions of the field products (no locality);
the commutator formula ($\ell = 0$) follows from the delta-decomposition of $[a,b]$, whose coefficients
are by construction the nonnegative products (\lean{commutator_prod_bound_of_local}); the Pascal
recursion and the anti-diagonal pin finish. \qed

\begin{status}
Unconditional, axiom-clean (\texttt{\#print axioms} \lean{VOA.Fld.field_jacobi_of_local} and
\lean{VOA.Fld.mode_prod_iterate} run for this chapter).
\end{status}

\subsection{The carrier action, and what is not yet proved}

\begin{definition}[Induced state--field map of a local system]\label{l03:def:statefield}
For a local field system $L$, the vacuum is $1\in L$, the mode action is $a\mapsto(b\mapsto a_{(n)}b)$,
the induced field of $a\in L$ is $\mathcal Y(a,z) = \sum_n(a_{(n)}\cdot)z^{-n-1}$ acting on $L$, and the
translation is $a\mapsto a_{(-2)}1$.
\end{definition}
\begin{leancode}
noncomputable def vacuum : L.carrier := ⟨Fld.one, L.one_mem⟩

noncomputable def stateField : L.carrier →ₗ[ℂ] FormalDist (End ℂ L.carrier) where
  toFun a := fun q => modeAction L (-q - 1) a

theorem stateField_vacuum :
    stateField L (vacuum L) = (idField : FormalDist (End ℂ L.carrier))

theorem creation_deriv (a : L.carrier) (k : ℕ) :
    (modeAction L (-(k : ℤ) - 1) a (vacuum L) : Fld M).1 =
      (k.factorial : ℂ)⁻¹ • (VOA.deriv^[k]) a.1.1

theorem translation_val (a : L.carrier) :
    (translation L a : Fld M).1 = VOA.deriv a.1.1
\end{leancode}
File \leanpath{VOA/\allowbreak Core/\allowbreak LocalSystemAction.lean} (namespace
\lean{VOA.LocalFieldSystem}; \lean{modeAction L n a b} is $a_{(n)}b$).

\begin{dictionary}
\lean{L.carrier} is the subspace, used as a type. \lean{stateField} has $z^q$-coefficient
\lean{modeAction L (-q - 1)}, so its modes are the field products, as in any vertex algebra.
\lean{stateField_vacuum} is VA2, \lean{creation_deriv} the creation tower, \lean{translation_val} says the
induced translation is $\partial$. The file also proves \lean{creation_neg_one}, \lean{creation_nonneg}
(VA3) and \lean{modeAction_eventually_zero} (VA5).
\end{dictionary}

\begin{physics}
The operators of a local system themselves form a chiral algebra: the identity operator is its vacuum,
$\partial$ its translation, and OPE coefficients its products.
\end{physics}

\begin{warning}[title={Li's construction is not complete at this commit}]
The library proves VA2, VA3, VA5 and the creation tower for the induced fields of a local field
system, Dong closure (\lean{LocalFieldSystem.prod_local}), and the Borcherds identity for the action
on $M$. It does \textbf{not} prove that the induced fields \lean{stateFld L a} are mutually local on
$L$ (VA4 for the local system), nor the compatibility $\mathcal Y(a_{(n)}b) = \mathcal Y(a)_{(n)}\mathcal Y(b)$,
and it packages neither a \lean{BarePreVA} nor a \lean{PreVertexAlgebra} on $L$, nor a module structure on
$M$. The file \leanpath{LocalSystemRegular.lean} names these as ``the intended endpoints''. In
particular no result of the library is conditional on them; they are open work. (If the endpoint is
a \lean{BarePreVA}, the warning of Section~\ref{l03:sec:counter} applies: a Borcherds identity over
\lean{BarePreVA} would require adding translation covariance, which a local system does have,
\lean{translation_val}.)
\end{warning}

%======================================================================
\section{Localising the first-order Jacobi identity}\label{l03:sec:jacobilocal}
%======================================================================

The file \leanpath{VOA/\allowbreak Core/\allowbreak JacobiLocal.lean} serves the deformation theory of
Chapter~\ref{ch:L07}, where a first-order deformation $Y+tY_1$ of a VOA is encoded by a
\emph{cochain} $f:V\to\End(V)[[z^{\pm1}]]$ (\lean{VOA.Deform.Cochain}), with deformed products
$u^1_{(n)}v = $ \lean{cmode f u n v}, and the cocycle condition \lean{JacobiAx W f} is the linearisation
of the Borcherds identity in $f$ (four insertions of $f$ on the left, two on the right). Chapter~\ref{ch:L07}
gives the definitions of \lean{Cochain}, \lean{cmode}, \lean{JacobiAx}, \lean{gradedCochains} and
\lean{IsRigid} with their Lean source; we state only the localisation theorem proved here.

\begin{lemma}[Products vanish above the grading bound]\label{l03:lem:gradingbound}
In a VOA, if $u\in V_a$, $v\in V_b$ and $n>a+b-1$, then $u_{(n)}v = 0$.
\end{lemma}
\begin{leancode}
theorem smul_eq_zero_of_grading_lt (W : VertexOperatorAlgebra 𝒱) {a b : ℤ} {u v : 𝒱}
    (hu : u ∈ W.grading a) (hv : v ∈ W.grading b) {n : ℤ} (hn : a + b - 1 < n) :
    W.smul u n v = 0
\end{leancode}
File \leanpath{VOA/\allowbreak Core/\allowbreak JacobiLocal.lean} (\lean{VOA.Deform.smul_eq_zero_of_grading_lt}).

\begin{dictionary}
VA1 puts $u_{(n)}v$ in $V_{a+b-n-1}$, which is $0$ by (VOA4) when $a+b-n-1<0$. This explicit bound is
what makes the Borcherds sums finite with an explicit range once the grades are fixed.
\end{dictionary}

\begin{physics}
In a theory with nonnegative scaling dimensions, the OPE $u(z)v(w)$ has a pole of order at most
$\Delta_u+\Delta_v$, the pole being fixed by the dimension of the lowest state that can appear.
\end{physics}

\begin{theorem}[Localisation of the cocycle condition]\label{l03:thm:jacobilocal}
For a graded cochain $f$ on a VOA, \lean{JacobiAx W f} holds if and only if, for every
$(a,b,d,\ell,m,n)\in\Z^6$, the explicit finite equation \lean{LocalJacobiAt W f a b d ℓ m n} holds:
for all $u\in V_a$, $v\in V_b$, $x\in V_d$,
$\sum_{i=0}^{B_L}\mathrm{jac1L}_i = \sum_{i=0}^{B_R}\mathrm{jac1R}_i$ with
$B_L = \max(b+d-n-1,\,a+d-m-1)$ and $B_R = a+b-\ell-1$.
\end{theorem}
\begin{leancode}
def LocalJacobiAt (W : VertexOperatorAlgebra 𝒱) (f : Cochain 𝒱) (a b d ℓ m n : ℤ) : Prop :=
  ∀ u ∈ W.grading a, ∀ v ∈ W.grading b, ∀ x ∈ W.grading d,
    (∑ i ∈ Finset.Icc (0 : ℤ) (jacBoundL a b d m n), jac1L W f u v x ℓ m n i)
      = ∑ i ∈ Finset.Icc (0 : ℤ) (jacBoundR a b ℓ), jac1R W f u v x ℓ m n i

theorem jacobiAx_iff_localJacobiAt (W : VertexOperatorAlgebra 𝒱) {f : Cochain 𝒱}
    (hf : f ∈ gradedCochains W) :
    JacobiAx W f ↔ ∀ a b d ℓ m n : ℤ, LocalJacobiAt W f a b d ℓ m n
\end{leancode}
File \leanpath{VOA/\allowbreak Core/\allowbreak JacobiLocal.lean} (namespace \lean{VOA.Deform}). The
summands are
\begin{leancode}
noncomputable def jac1L (W : VertexOperatorAlgebra 𝒱) (f : Cochain 𝒱) (u v x : 𝒱)
    (ℓ m n i : ℤ) : 𝒱 :=
  (negOnePow i * cbz ℓ i) •
    (cmode f u (ℓ + m - i) (W.smul v (n + i) x)
      + W.smul u (ℓ + m - i) (cmode f v (n + i) x)
      - negOnePow ℓ •
        (cmode f v (ℓ + n - i) (W.smul u (m + i) x)
          + W.smul v (ℓ + n - i) (cmode f u (m + i) x)))

noncomputable def jac1R (W : VertexOperatorAlgebra 𝒱) (f : Cochain 𝒱) (u v x : 𝒱)
    (ℓ m n i : ℤ) : 𝒱 :=
  cbz m i •
    (W.smul (cmode f u (ℓ + i) v) (m + n - i) x
      + cmode f (W.smul u (ℓ + i) v) (m + n - i) x)

def jacBoundL (a b d m n : ℤ) : ℤ := max (b + d - n - 1) (a + d - m - 1)

def jacBoundR (a b ℓ : ℤ) : ℤ := a + b - ℓ - 1
\end{leancode}
File \leanpath{VOA/\allowbreak Core/\allowbreak JacobiLocal.lean}.

\begin{dictionary}
\lean{Finset.Icc 0 N} is $\{0,\dots,N\}$ (empty if $N<0$). \lean{jac1L} and \lean{jac1R} are the $i$-th
summands of the two sides of the linearised Borcherds identity: the four (resp.\ two) ways of replacing
one product by a deformed product \lean{cmode f}. The theorem replaces each \lean{finsum} by an explicit
finite sum, using Lemma~\ref{l03:lem:gradingbound} and its analogue for graded cochains, and uses the
direct-sum decomposition to reduce from arbitrary to homogeneous $u,v,x$. \textbf{Scope (stated in the
source):} for fixed input grades $(a,b,d)$ there are still infinitely many equations ($\ell\to-\infty$),
so the system is not finite per tridegree; the theorem is a reformulation, not a finiteness theorem.
The same file proves the consumer \lean{not_isRigid_of_local}: a solution of the local system whose
component at some tridegree is not a coboundary component refutes \lean{IsRigid}; by
Chapter~\ref{ch:L07}, \lean{IsRigid} is the ungraded rigidity notion that is false for every Virasoro VOA
for definitional reasons, so this consumer is of limited use.
\end{dictionary}

\begin{physics}
First-order deformations of a chiral algebra are infinitesimal changes of all OPE coefficients
preserving associativity; the theorem says that, sector by sector in scaling dimensions, the
associativity constraints on the deformed OPE coefficients are finitely many linear equations each ---
the form in which conformal-bootstrap-style constraints on marginal deformations are actually
solved.
\end{physics}

\begin{status}
\lean{jacobiAx_iff_localJacobiAt}, \lean{smul_eq_zero_of_grading_lt}, \lean{not_isRigid_of_local}:
unconditional, axiom-clean (\texttt{\#print axioms} run for this chapter).
\end{status}

\section{Combinatorial lemmas}\label{l03:sec:combinatorics}

The file \leanpath{VOA/\allowbreak Core/\allowbreak FieldCombinatorics.lean} (namespace
\lean{VOA.FieldCombinatorics}) collects coefficient identities used by the field-product files; it
is independent of vertex algebras. The one with mathematical content is the negative binomial:

\begin{lemma}\label{l03:lem:negbinom}
For $s,n\in\N$, $\binom{-s-1}{n} = (-1)^n\frac{(s+n)!}{s!\,n!} = (-1)^n\binom{s+n}{n}$.
\end{lemma}
\begin{leancode}
theorem gBinom_negsucc (s n : ℕ) :
    gBinom (-(s : ℤ) - 1) n =
      (-1) ^ n * ((s + n).factorial : ℂ) /
        ((s.factorial : ℂ) * (n.factorial : ℂ))
\end{leancode}
File \leanpath{VOA/\allowbreak Core/\allowbreak FieldCombinatorics.lean}.

\begin{dictionary}
Exactly the displayed identity; it converts the coefficients of normal-ordered products into the
binomials of the iterate formula. Other lemmas there: the expansion of $(z_1-z_2)^N$ coefficientwise
(\lean{mulBySubPow_apply_sum}), $\binom nk$ for natural $n$ (\lean{gBinom_natCast_eq_choose}), binomial
symmetry (\lean{gBinom_symm_nat}) and the coefficients of $\partial^j$ (\lean{deriv_pow_apply}).
\end{dictionary}

\begin{physics}
$\binom{-s-1}{n}$ is the Taylor coefficient of $(z-w)^{-s-1}$ in powers of $w/z$: it governs the
expansion of an OPE pole of order $s+1$ in one radial ordering.
\end{physics}

%======================================================================
\section{Referee's checklist}\label{l03:sec:checklist}
%======================================================================

\begingroup\small
\begin{longtable}{>{\raggedright\arraybackslash}p{0.28\textwidth}>{\raggedright\arraybackslash}p{0.31\textwidth}>{\raggedright\arraybackslash}p{0.2\textwidth}>{\raggedright\arraybackslash}p{0.1\textwidth}}
\toprule
Lean declaration & Intended meaning & Hypotheses / status & Check\\
\midrule\endhead
\lean{PreVertexAlgebra} & Gannon Def.~5.1.3(a): grading, VA1--VA5 & definition; \textbf{no translation axiom}; strictly weaker than Kac/FBZ/LL in that direction & Def.~\ref{l03:def:VA}, §\ref{l03:sec:compare}\\
\lean{BarePreVA} & VA2--VA5 without grading & definition; docstring now warns that it does \textbf{not} suffice for Borcherds (an earlier docstring claimed it did) & Def.~\ref{l03:def:bare}\\
\lean{GradedVertexOperatorAlgebra} & LL/FLM VOA (graded, truncated below), without CFT type & definition; no structure theory proved at this level & Def.~\ref{l03:def:VOA}\\
\lean{VertexOperatorAlgebra} & LL/FLM VOA of CFT type (a restriction of LL/FLM) & definition; (VOA3) supplies translation; LL $\Rightarrow$ library direction on paper only & Prop.~\ref{l03:prop:equiv}\\
\lean{VertexAlgebra}, \lean{toVertexAlgebra} & \lean{PreVertexAlgebra} + translation covariance; every VOA is one & definition; map from VOAs via \lean{isTranslationCovariant} & §\ref{l03:sec:intro}\\
(scratch) \lean{L03Counter.VA}, \lean{jacobi_fails} & a \lean{PreVertexAlgebra} violating Borcherds & machine-checked outside the library (not in this release); library copy \lean{TranslationCovariantCounterexample} & Thm.~\ref{l03:thm:counter}\\
\lean{PreVertexAlgebra.T} & $Tu = u_{(-2)}\lthreevac$ & definition & Def.~\ref{l03:def:smulT}\\
\lean{Yprod_ext}, \lean{eq_}\allowbreak\lean{zero_}\allowbreak\lean{of_}\allowbreak\lean{mulBySubPow_}\allowbreak\lean{of_}\allowbreak\lean{bddBelow} & $z_1-z_2$ non-zero-divisor on $V((z_1))((z_2))$ & none; unconditional & Thm.~\ref{l03:thm:yprodext}\\
\lean{bracket_deltaSum}, \lean{mode_}\allowbreak\lean{commutator_}\allowbreak\lean{structural} & OPE decomposition; structural commutator & VA level; coefficients not identified & Thm.~\ref{l03:thm:ope}\\
\lean{goddard_uniqueness}, \lean{goddard_uniqueness_vac}, \lean{goddard_zero} & state determines local field & VA level / abstract; unconditional & Thm.~\ref{l03:thm:goddard}\\
\lean{dong_nonneg}, \lean{dong_neg} & locality closed under field products & pairwise locality (+ field); unconditional & Thm.~\ref{l03:thm:dong}\\
\lean{creation_pow}, \lean{Y_Tpow}, \lean{T_eq_Lneg1} & $Y(u,z)\lthreevac = e^{zT}u$, $T = L_{-1}$ & VOA; unconditional & Thm.~\ref{l03:thm:creation}\\
\lean{Y_T}, \lean{mode_Y_T} & $Y(Tu,z) = \partial_zY(u,z)$; $(Tu)_{(n)} = -n\,u_{(n-1)}$ & VOA; unconditional & Thm.~\ref{l03:thm:TL}\\
\lean{translation_covariance} & $[L_{-1},Y(u)] = Y(L_{-1}u)$ & VOA; unconditional & Thm.~\ref{l03:thm:transcov}\\
\lean{nth_product}, \lean{assoc_nonneg} & $Y(u_{(k)}v) = \Res(z_1-z)^k[Y(u),Y(v)]$ & VOA; unconditional & Thm.~\ref{l03:thm:nthproduct}\\
\lean{mode_commutator}, \lean{commBound} & commutator formula (complete sum) & VOA; $\exists N$ made exhaustive by \lean{commBound} & Thm.~\ref{l03:thm:comm}\\
\lean{assoc_neg} & $Y(u_{(-k-1)}v) = \frac1{k!}{:}\partial^kY(u)Y(v){:}$ & VOA; unconditional & Thm.~\ref{l03:thm:assocneg}\\
\lean{mode_iterate} & iterate formula, all $\ell\in\Z$ & VOA; vector form; unconditional & Thm.~\ref{l03:thm:iterate}\\
\lean{borcherds}, \lean{jacobi} & Borcherds identity, all $(\ell,m,n)$ & VOA; unconditional, axiom-clean & Thm.~\ref{l03:thm:borcherds}\\
\lean{jacobi_commutator} & commutator formula, $\ell=0$ slice of Borcherds (alias of \lean{mode_commutator}) & VOA; $\exists N$ (see \lean{commBound}); unconditional & Thm.~\ref{l03:thm:jacobicomm}\\
\lean{skewVacLocal}, \lean{weakAssocLoc} & reflected field is local & VOA; unconditional & Thm.~\ref{l03:thm:skewvac}\\
\lean{skew}, \lean{skew_mode} & $Y(u,z)v = e^{zT}Y(v,-z)u$ & VOA; unconditional & Thm.~\ref{l03:thm:skew}\\
\lean{flm} & FLM associativity after $z_1 = z_0+z_2$ & VOA, explicit locality order; unconditional & Thm.~\ref{l03:thm:flm}\\
\lean{comm_iff_nonneg_modes}, \lean{comm_iff_assoc_alg}, \lean{commRing_of_comm} & Gannon Thm.~5.1.4 & VOA; \textbf{hypothesis forces $W = \C\lthreevac$} & Thm.~\ref{l03:thm:commalg}, warning\\
\lean{Reconstruction.}\allowbreak\lean{reconstruct} & existence of $Y$ from spanning local $T$-covariant fields & all hypotheses explicit; no VA2/VA1 in conclusion & Thm.~\ref{l03:thm:reconstruct}\\
\lean{Reconstruction.}\allowbreak\lean{field_eq_zero_of_vacuum} & a field local with a spanning family of creating fields and killing the vacuum is $0$ & bare vector space; all hypotheses explicit; unconditional & Thm.~\ref{l03:thm:fieldzero}\\
\lean{TensorFieldLift.*} & lifts to $A\otimes B$ preserve hypotheses & unconditional & Thm.~\ref{l03:thm:tensorlift}\\
\lean{gen_local}, \lean{generated_local} & generated system is local & pairwise local generators & Thm.~\ref{l03:thm:genlocal}\\
\lean{Fld.prod}, \lean{LocalFieldSystem} & Li's field products and local systems & definitions & Def.~\ref{l03:def:lfs}\\
\lean{Fld.}\allowbreak\lean{field_}\allowbreak\lean{jacobi_}\allowbreak\lean{of_}\allowbreak\lean{local}, \lean{Fld.mode_prod_iterate} & Borcherds for any local pair of fields & \lean{Local a b}; unconditional & Thm.~\ref{l03:thm:fieldjacobi}\\
\lean{LocalFieldSystem.}\allowbreak\lean{stateField} etc. & induced VA on a local system & partial: \textbf{VA4 for induced fields not proved} & Def.~\ref{l03:def:statefield}, warning\\
\lean{Deform.}\allowbreak\lean{jacobiAx_}\allowbreak\lean{iff_}\allowbreak\lean{localJacobiAt} & cocycle condition = local finite equations & graded cochain; unconditional & Thm.~\ref{l03:thm:jacobilocal}\\
\bottomrule
\end{longtable}
\endgroup

\textbf{Machine verification for this chapter.} When the chapter was first written, \texttt{\#print axioms} was run on a scratch file
importing the chapter's modules, with \texttt{lake env lean -DmaxSynthPendingDepth=3} (never the
bare command), against the build products of an earlier commit (not re-run for this edition); a full \texttt{lake build} was
\emph{not} run for this chapter. Every one of the following reported exactly
\lean{[propext, Classical.}\allowbreak\lean{choice, Quot.}\allowbreak\lean{sound]}:
\lean{goddard_uniqueness}, \lean{goddard_zero}, \lean{Yprod_ext}, \lean{bracket_deltaSum},
\lean{dong_nonneg}, \lean{dong_neg}, \lean{T_eq_Lneg1}, \lean{creation_pow}, \lean{Y_Tpow},
\lean{translation_covariance}, \lean{nth_product}, \lean{mode_commutator}, \lean{assoc_neg},
\lean{mode_iterate}, \lean{skew}, \lean{skew_mode}, \lean{weakAssocLoc}, \lean{flm}, \lean{borcherds},
\lean{jacobi}, \lean{jacobi_commutator}, \lean{Y_T}, \lean{field_eq_zero_of_vacuum}, \lean{comm_iff_nonneg_modes}, \lean{comm_iff_assoc_alg}, \lean{commRing_of_comm},
\lean{reconstruct}, \lean{gen_local}, \lean{field_jacobi_of_local}, \lean{mode_prod_iterate},
\lean{prod_local_right_of_pairwise}, \lean{creation_deriv}, \lean{jacobiAx_iff_localJacobiAt},
\lean{smul_eq_zero_of_grading_lt}, \lean{not_isRigid_of_local}, \lean{local_firstField_secondField},
\lean{BarePreVA.isField_Y} (full names as in the sections above).
The PBW theorem \lean{VOA.pbw} occurs nowhere in this chapter.

\textbf{Release audit 2026-10-27:} \texttt{Notes/\allowbreak Textbook/\allowbreak referee/\allowbreak RELEASE\_AUDIT\_FOUNDATIONS.md}
(at an earlier commit; no \texttt{lake build}, \texttt{\#print axioms} on the existing build products). It re-ran
\texttt{\#print axioms} on \lean{VOA.VertexOperatorAlgebra.borcherds} and \lean{VOA.VertexOperatorAlgebra.jacobi} (three
standard axioms each), confirmed that the Borcherds identity is stated for all $(\ell,m,n)\in\Z^3$, that its three
\lean{finsum}s are proved finite (\lean{BR_fin}, \lean{Puv_fin}, \lean{Pvu_fin}), and that it is derived from locality,
not assumed. It recompiled the counterexample \lean{L03Counter.VA} (\lean{translation_fails}, \lean{jacobi_fails}:
three standard axioms) and adopted this chapter's finding: \lean{PreVertexAlgebra}/\lean{BarePreVA} are not vertex
algebras in the standard sense; \lean{VertexOperatorAlgebra} is the LL/FLM VOA of CFT type; the direction
LL $\Rightarrow$ library of Proposition~\ref{l03:prop:equiv} is paper-only; the zero space is a VOA.

%======================================================================
\section{Notes and references}\label{l03:sec:notes}
%======================================================================

\textbf{Original sources.} R.~E.~Borcherds, ``Vertex algebras, Kac--Moody algebras, and the Monster'',
\emph{Proc.\ Nat.\ Acad.\ Sci.\ USA} 83 (1986), 3068--3071 (the definition by the Borcherds identity).
I.~B.~Frenkel, J.~Lepowsky and A.~Meurman, \emph{Vertex Operator Algebras and the Monster}, Pure and
Applied Mathematics 134, Academic Press (1988), Chapter~8 (the Jacobi identity, weak associativity,
the definition of a VOA).

\textbf{Textbooks.} V.~G.~Kac, \emph{Vertex Algebras for Beginners}, 2nd ed., University Lecture
Series 10, AMS (1998): Definition~1.3 (translation operator as data), §§2--3 (fields, locality,
Dong's lemma), §4 (uniqueness, existence, $n$-th product identity, Borcherds identity).
E.~Frenkel and D.~Ben-Zvi, \emph{Vertex Algebras and Algebraic Curves}, 2nd ed., Mathematical
Surveys and Monographs 88, AMS (2004): Definition~1.3.1 and §4.4 (reconstruction).
J.~Lepowsky and H.~Li, \emph{Introduction to Vertex Operator Algebras and Their Representations},
Progress in Mathematics 227, Birkh\"auser (2004): Chapter~3 (Jacobi identity, commutator and iterate
formulas, skew-symmetry, weak commutativity and associativity, Definition~3.1.22 of a VOA).
T.~Gannon, \emph{Moonshine beyond the Monster}, Cambridge University Press (2006), §5.1: Definition
5.1.3 (the library's axioms), Proposition~5.1.2, (5.1.7)--(5.1.9), Theorem~5.1.4.

\textbf{Further.} H.~Li, ``Local systems of vertex operators, vertex superalgebras and modules'',
\emph{J.\ Pure Appl.\ Algebra} 109 (1996), 143--195.
For
Hasse--Schmidt derivations (the structure behind the counterexample of Section~\ref{l03:sec:counter})
see H.~Hasse and F.~K.~Schmidt, \emph{J.\ Reine Angew.\ Math.} 177 (1937), 215--237.
Physics background: A.~A.~Belavin, A.~M.~Polyakov and A.~B.~Zamolodchikov, \emph{Nucl.\ Phys.\ B} 241
(1984), 333--380; P.~Di~Francesco, P.~Mathieu and D.~S\'en\'echal, \emph{Conformal Field Theory},
Springer (1997), Chapter~6 (OPE, radial ordering, mode algebras).

\textbf{Project records.} The history of the skew-symmetry gate (once the sanctioned axiom
\lean{weakAssocLoc_gate}, then a single \lean{sorry} in \lean{Gfield_eq_jumpSum}) is recorded only in
the docstrings of \leanpath{SkewSymmetry.lean} and \leanpath{WeakAssoc.lean}. The translation gap of
\lean{PreVertexAlgebra} and the vacuity of Theorem~5.1.4 for VOAs are recorded for the first time here.

\endgroup
